\documentclass[11pt,a4paper]{article}
\usepackage[T2A]{fontenc}
\usepackage[utf8]{inputenc}
\usepackage[russian]{babel}
\usepackage{amsmath,amssymb,amsthm,mathtools}
\usepackage[margin=2.3cm]{geometry}
\usepackage{enumitem}
\usepackage{array,booktabs}
\usepackage{tikz}
\usetikzlibrary{arrows.meta,calc}
\usepackage{tikz-cd}
\usepackage[pdfencoding=auto,psdextra,colorlinks=true,linkcolor=blue!50!black,citecolor=green!40!black,urlcolor=blue!60!black]{hyperref}

\setlist{itemsep=2pt,topsep=4pt}

\theoremstyle{plain}
\newtheorem{theorem}{Теорема}[section]
\newtheorem{proposition}[theorem]{Предложение}
\newtheorem{lemma}[theorem]{Лемма}
\newtheorem{corollary}[theorem]{Следствие}
\newtheorem{conjecture}[theorem]{Гипотеза}
\theoremstyle{definition}
\newtheorem{definition}[theorem]{Определение}
\newtheorem{example}[theorem]{Пример}
\newtheorem{exercise}[theorem]{Упражнение}
\newtheorem{question}[theorem]{Вопрос}
\newtheorem{construction}[theorem]{Конструкция}
\theoremstyle{remark}
\newtheorem{remark}[theorem]{Замечание}

\renewcommand{\C}{\mathbb{C}}
\newcommand{\Z}{\mathbb{Z}}
\newcommand{\Q}{\mathbb{Q}}
\newcommand{\R}{\mathbb{R}}
\newcommand{\PP}{\mathbb{P}}
\newcommand{\A}{\mathbb{A}}
\newcommand{\Gm}{\mathbb{G}_m}
\newcommand{\cO}{\mathcal{O}}
\newcommand{\cL}{\mathcal{L}}
\newcommand{\cE}{\mathcal{E}}
\newcommand{\cF}{\mathcal{F}}
\newcommand{\cI}{\mathcal{I}}
\newcommand{\cX}{\mathcal{X}}
\newcommand{\cM}{\mathcal{M}}
\newcommand{\mm}{\mathfrak{m}}
\newcommand{\fa}{\mathfrak{a}}
\DeclareMathOperator{\Pic}{Pic}
\DeclareMathOperator{\Cl}{Cl}
\DeclareMathOperator{\CaCl}{CaCl}
\DeclareMathOperator{\Div}{Div}
\DeclareMathOperator{\divv}{div}
\DeclareMathOperator{\ord}{ord}
\DeclareMathOperator{\Bl}{Bl}
\DeclareMathOperator{\Hom}{Hom}
\DeclareMathOperator{\rk}{rk}
\DeclareMathOperator{\codim}{codim}
\DeclareMathOperator{\Sing}{Sing}
\DeclareMathOperator{\GL}{GL}
\DeclareMathOperator{\SL}{SL}
\DeclareMathOperator{\PGL}{PGL}
\DeclareMathOperator{\Aut}{Aut}
\DeclareMathOperator{\Stab}{Stab}
\DeclareMathOperator{\Cone}{Cone}
\DeclareMathOperator{\conv}{conv}
\DeclareMathOperator{\NE}{\overline{NE}}
\DeclareMathOperator{\im}{Im}
\DeclareMathOperator{\id}{id}
\DeclareMathOperator{\Ric}{Ric}
\DeclareMathOperator{\vol}{vol}
\DeclareMathOperator{\Sym}{Sym}
\DeclareMathOperator{\Fut}{Fut}
\DeclareMathOperator{\lct}{lct}
\DeclareMathOperator{\discrep}{discrep}
\DeclareMathOperator{\Supp}{Supp}
\DeclareMathOperator{\mult}{mult}
\DeclareMathOperator{\length}{length}
\DeclareMathOperator{\Hilb}{Hilb}
\DeclareMathOperator{\Proj}{Proj}
\DeclareMathOperator{\Spec}{Spec}
\DeclareMathOperator{\Val}{Val}
\DeclareMathOperator{\coeff}{coeff}
\DeclareMathOperator{\gr}{gr}
\DeclareMathOperator{\Nef}{Nef}
\DeclareMathOperator{\Eff}{\overline{Eff}}
\newcommand{\hvol}{\widehat{\vol}}
\newcommand{\an}{\mathrm{an}}
\newcommand{\sm}{\mathrm{sm}}
\newcommand{\ssl}{\mathrm{ss}}

\pagenumbering{gobble}

\begin{document}

\begin{titlepage}
	\begin{center}
		
		{\large Московский государственный университет им.\ М.\,В.\,Ломоносова}
		
		\vspace{0.5cm}
		
		{\large Механико-математический факультет}
		
		\vspace{0.5cm}
		
		{\large Кафедра дифференциальной геометрии и приложений}
		
		\vspace{3cm}
		
		{\LARGE \textbf{K-модули кубических поверхностей:\\[4pt] пререквизиты, K-стабильность и построение пространства модулей}}
		
		\vspace{0.5cm}
		
		
		\vfill
		
		\begin{flushright}
			Выполнил: \\
			Claude Fable 5.1,\\
			Студент группы 407 \\
			Азнаурьян Л.\,П. \\
			
			
			\vspace{1cm}
			
			Научные руководители: \\
			Логинов К.\,В.\\
			Claude Fable 5.1
		\end{flushright}
		
		\vfill
		
		Москва \\
		2026
		
	\end{center}
\end{titlepage}

\newpage

\tableofcontents
\newpage
\pagenumbering{arabic}

\setcounter{section}{-1}
% ============================================================
\section{Введение: о чём этот текст и как его читать}\label{sec:intro}
% ============================================================

Этот текст --- подробный русскоязычный разбор обзора К.~ДеВлеминг «On moduli of Fano varieties:
an introduction to K-stability and K-moduli» \cite{DeV}, написанного для летней школы 2025~года и
рассчитанного на аспирантов, прослушавших первый курс алгебраической геометрии «на уровне
Хартсхорна». Цель обзора --- довести читателя от определения K-стабильности до описания
пространства модулей K-полустабильных кубических поверхностей. Цель настоящего текста ---
сделать этот путь проходимым для того, кто знает классическую алгебраическую геометрию в объёме
конспекта \cite{AG} (дивизоры, канонический класс, формула присоединения, поверхности дель Пеццо,
раздутия, 27 прямых), но не знаком с особенностями программы минимальных моделей, с геометрической
теорией инвариантов и с кэлеровой геометрией.

\subsection*{Структура}

Текст состоит из двух частей.

\emph{Часть~I (разделы \ref{sec:div}--\ref{sec:dg}) --- пререквизиты.} Здесь собрано всё, что
в обзоре используется без доказательства или без определения:
\begin{itemize}
\item $\Q$-дивизоры, канонический класс нормального многообразия, пары, обратные образы дивизоров
при бирациональных морфизмах, дивизоры «над» многообразием (раздел~\ref{sec:div});
\item теория пересечений на поверхностях, конусы кривых и дивизоров, численная эффективность,
большие дивизоры, объём, разложение Зарисского, асимптотическая формула Римана--Роха
(раздел~\ref{sec:cones});
\item раздутия, разрешения особенностей, взвешенные проективные пространства, факторособенности
поверхностей и особенности Дю Валя, $T$-особенности и $\Q$-горенштейновы сглаживания, особые
поверхности дель Пеццо (раздел~\ref{sec:sing});
\item семейства, схемы Гильберта, $\Q$-горенштейновы семейства (раздел~\ref{sec:families});
\item геометрическая теория инвариантов: редуктивные группы, факторы, критерий Гильберта--Мамфорда,
бинарные формы, плоские кубики и --- подробно --- классическая теорема Гильберта о стабильности
кубических поверхностей (раздел~\ref{sec:git});
\item дифференциальная геометрия: кэлеровы метрики, форма Риччи, уравнение Кэлера--Эйнштейна,
препятствия Мацусимы и Футаки, $\alpha$-инвариант Тяня, гипотеза Яу--Тяня--Дональдсона
(раздел~\ref{sec:dg}).
\end{itemize}

\emph{Часть~II (разделы \ref{sec:mmp}--\ref{sec:cubic}) --- K-стабильность и K-модули.} Это
перевод основной части обзора \cite{DeV} на русский язык с существенными дополнениями: все
определения снабжены пояснениями, все вычисления обзора проведены подробно, большинство упражнений
обзора решено и оформлено в виде примеров (простых и более сложных), добавлены примеры, которых в
обзоре нет (торические вычисления $\beta$-инварианта, тест-конфигурации с явными весами,
вырождение $\PP^2$ в $\PP(1,1,4)$, полистабильность кубики $xyz=w^3$ и~др.). Соответствие с обзором:
\[
\begin{array}{ll}
\text{раздел~\ref{sec:mmp} --- особенности программы минимальных моделей} & \text{(\S2 обзора)}\\
\text{раздел~\ref{sec:tc} --- K-стабильность через тест-конфигурации} & \text{(\S3.1)}\\
\text{раздел~\ref{sec:alpha} --- $\alpha$-инвариант} & \text{(\S3.2)}\\
\text{раздел~\ref{sec:beta} --- $\beta$- и $\delta$-инварианты, валюативный критерий} & \text{(\S3.3)}\\
\text{раздел~\ref{sec:az} --- теория допустимых флагов Аббана--Чжуана} & \text{(\S4)}\\
\text{раздел~\ref{sec:moduli} --- теорема о K-модулях} & \text{(\S5)}\\
\text{раздел~\ref{sec:cubic} --- нормализованный объём и K-модули кубических поверхностей} & \text{(\S6)}
\end{array}
\]
Раздел~\ref{sec:concl} содержит заключительные замечания и список упражнений обзора, не разобранных
в тексте.

\subsection*{Соглашения}

Основное поле --- $\C$. \emph{Многообразие} --- неприводимое отделимое многообразие конечного типа
над $\C$; начиная с раздела~\ref{sec:div} все многообразия предполагаются \emph{нормальными}
(это стандартное соглашение в теории K-стабильности). Точка --- замкнутая точка. Для
квазипроективных многообразий мы свободно пользуемся классическим языком \cite{AG}; там, где
нужны схемы (схемы Гильберта, плоские семейства, стеки), это оговаривается. Ссылки на конспект
\cite{AG} даются с номерами разделов; например, «\cite[\S\,5]{AG}» --- раздел о формуле
присоединения. Результаты, доказательства которых выходят за рамки текста, формулируются с точными
ссылками (чаще всего на книгу Сюя \cite{Xu} и книгу Коллара--Мори \cite{KM}).

\subsection*{Главная сюжетная линия}

Чтобы не потерять нить среди технических деталей, сформулируем сразу, куда мы идём.

\begin{itemize}
\item Многообразие Фано --- это многообразие с обильным антиканоническим классом. Пространства
модулей многообразий Фано «в лоб» не существует: оно неотделимо, неограниченно и содержит
многообразия с нередуктивными группами автоморфизмов (раздел~\ref{sec:moduli}).
\item K-стабильность --- условие, пришедшее из дифференциальной геометрии (из задачи о метриках
Кэлера--Эйнштейна, раздел~\ref{sec:dg}), которое в точности отсекает «плохие» многообразия Фано.
Оно имеет несколько равносильных алгебраических формулировок: через тест-конфигурации и
инвариант Футаки (раздел~\ref{sec:tc}), через валюативный критерий Фудзиты--Ли
(раздел~\ref{sec:beta}), через нормализованный объём (раздел~\ref{sec:cubic}).
\item Теорема о K-модулях (теорема~\ref{thm:Kmoduli}): K-полустабильные многообразия Фано данной
размерности и данного объёма образуют стек конечного типа с проективным пространством модулей,
точки которого --- K-полистабильные многообразия.
\item Для кубических поверхностей всё это можно довести до явного ответа (теоремы~\ref{thm:cubic-list}
и~\ref{thm:GIT=K}): K-полустабильные поверхности дель Пеццо степени~3 --- это ровно кубические
поверхности с особенностями типов $A_1$ и $A_2$, их K-стабильность совпадает со стабильностью в
смысле геометрической теории инвариантов, а пространство K-модулей --- это классический
ГИТ-фактор пространства кубических поверхностей, изоморфный взвешенному проективному пространству
$\PP(1,2,3,4,5)$.
\end{itemize}

\part{Пререквизиты}

% ============================================================
\section{Дивизоры, $\Q$-дивизоры и канонический класс нормального многообразия}\label{sec:div}
% ============================================================

В этом разделе собраны понятия, которые в обзоре \cite{DeV} появляются с первых строк: пара
$(X,\Delta)$, $\Q$-картьеров дивизор, $\Q$-горенштейновость, «дивизор над $X$» и его порядок
вдоль дивизора. Основа --- теория дивизоров Вейля и Картье на нормальных многообразиях
\cite[\S\,1]{AG}, которую мы кратко напомним.

\subsection{Напоминание: дивизоры Вейля и Картье}

Пусть $X$ --- нормальное многообразие. \emph{Простой дивизор} --- неприводимое замкнутое
подмногообразие коразмерности~1; \emph{дивизор Вейля} --- конечная целочисленная комбинация
простых дивизоров; $\Div X$ --- группа дивизоров Вейля. Для простого дивизора $Y$ локальное кольцо
$\cO_{X,Y}$ --- кольцо дискретного нормирования (нормальность!), и определён порядок
$\ord_Y f\in\Z$ рациональной функции $f\in\C(X)^*$ вдоль $Y$; $\divv f=\sum_Y\ord_Y(f)\,Y$.
\emph{Дивизор Картье} локально задаётся одним уравнением: это набор $\{(U_i,f_i)\}$ с
$f_i/f_j\in\cO^*(U_i\cap U_j)$; ему отвечает дивизор Вейля $\sum_Y\ord_Y(f_i)Y$ и обратимый пучок
$\cO_X(D)$. Основной результат \cite[теорема~1.10]{AG}:
\[
\Pic X\;\cong\;\CaCl X\;\hookrightarrow\;\Cl X,
\]
причём вложение является изоморфизмом тогда и только тогда, когда $X$ локально факториально
(например, гладко). Для произвольного дивизора Вейля $D$ по-прежнему определён пучок
$\cO_X(D)$: $\cO_X(D)(U)=\{f\in\C(X):\divv f+D\ge0\text{ на }U\}$; это \emph{рефлексивный} пучок
ранга~1, обратимый в точности тогда, когда $D$ --- дивизор Картье. Мы будем писать $H^0(X,D)$
вместо $H^0(X,\cO_X(D))$ и $|D|$ для соответствующей полной линейной системы.

\begin{example}[Квадратичный конус]\label{ex:cone-divisors}
На конусе $X=\{xy=z^2\}\subset\A^3$ прямая $\ell=\{x=z=0\}$ --- дивизор Вейля, не являющийся
дивизором Картье, но $2\ell=\divv(x)$ --- главный дивизор; $\Cl X=\Z/2$, $\Pic X=0$
\cite[пример~1.12]{AG}.
\end{example}

\begin{example}[Дивизор Вейля, никакое кратное которого не является дивизором Картье]\label{ex:nonQCartier}
Пусть $X=\{xy=zw\}\subset\A^4$ --- конус над гладкой квадрикой $Q\cong\PP^1\times\PP^1\subset\PP^3$,
$\Pi=\{x=z=0\}\subset X$ --- плоскость (конус над прямой одного из семейств образующих $Q$).
Тогда \emph{ни один} дивизор $m\Pi$, $m\ne0$, не является дивизором Картье. Действительно,
$X\setminus\{0\}$ --- расслоение на $\C^*$ над $Q$, и $\Cl X=\Cl(X\setminus0)=\Cl Q/\Z H=
\Z^2/\Z(1,1)\cong\Z$, где $H=(1,1)$ --- класс гиперплоского сечения, а $\Pi$ переходит в класс
$(1,0)$ бесконечного порядка (см.~\cite[II, упр.~6.3]{Ha}). Если бы $m\Pi$ был дивизором Картье,
то в окрестности вершины он был бы главным (у локального кольца тривиальна группа Пикара), то
есть $m\Pi\sim0$ на $X$, что невозможно при $m\ne0$.
\end{example}

\subsection{$\Q$-дивизоры}

\begin{definition}
\emph{$\Q$-дивизор} (Вейля) на нормальном $X$ --- элемент $\Div_\Q X=\Div X\otimes\Q$, т.\,е.
конечная сумма $D=\sum d_iD_i$ с $d_i\in\Q$ и простыми $D_i$. Дивизор \emph{эффективен}
($D\ge0$), если все $d_i\ge0$; $\Supp D=\bigcup_{d_i\ne0}D_i$; $\lfloor D\rfloor=\sum\lfloor d_i\rfloor D_i$.
$D$ называется \emph{$\Q$-картьеровым}, если $mD$ --- дивизор Картье для некоторого целого $m>0$;
наименьшее такое $m$ --- \emph{индекс} дивизора $D$. Два $\Q$-дивизора \emph{$\Q$-линейно
эквивалентны}, $D\sim_\Q D'$, если $m(D-D')$ --- главный (целочисленный) дивизор для некоторого
$m>0$; это равносильно тому, что $D-D'=\sum q_j\divv(f_j)$ с $q_j\in\Q$. Аналогично определяются
$\R$-дивизоры.
\end{definition}

Для $\Q$-картьерова дивизора $D$ и морфизма $f\colon Y\to X$ определён обратный образ
$f^*D:=\frac1m f^*(mD)$ (обратный образ дивизора Картье берётся через локальные уравнения;
определение корректно, так как $f^*$ аддитивен). Если $f$ бирационален, то $f^*D$ --- $\Q$-дивизор
на $Y$, не зависящий от выбора $m$.

Обильность, численная эффективность и~т.\,п. для $\Q$-дивизоров определяются через кратные:
$D$ обилен, если $mD$ --- обильный дивизор Картье; то же для «очень обильный», «nef», «большой»
(раздел~\ref{sec:cones}). Полезно помнить: класс $\Q$-картьерова дивизора лежит в $\Pic X\otimes\Q$,
и
\[
\Pic X\otimes\Q\;=\;\{\text{$\Q$-картьеровы дивизоры}\}/{\sim_\Q}\;\subset\;\Cl X\otimes\Q .
\]

\begin{example}
На конусе $xy=z^2$ дивизор $\ell$ $\Q$-картьеров индекса 2 и $\ell\sim_\Q\frac12\divv(x)\sim_\Q0$.
На $\PP(1,1,n)$ (раздел~\ref{sec:sing}) все дивизоры Вейля $\Q$-картьеровы, а индекс класса
$\cO(1)$ равен $n$. На конусе над квадрикой из примера~\ref{ex:nonQCartier} дивизор $\Pi$ не
$\Q$-картьеров.
\end{example}

\subsection{Канонический класс нормального многообразия}\label{ss:canonical-normal}

Пусть $X$ нормально, $\dim X=n$, $U=X\setminus\Sing X$ --- гладкое открытое подмножество; его
дополнение имеет коразмерность $\ge2$, поэтому $\Cl X=\Cl U$. На $U$ определён канонический дивизор
$K_U=\divv(\omega)$ рациональной $n$-формы \cite[\S\,3]{AG}.

\begin{definition}
\emph{Канонический дивизор} $K_X$ нормального многообразия $X$ --- замыкание в $X$ дивизора
$\divv(\omega)$ любой ненулевой рациональной $n$-формы $\omega$ (т.\,е. тот же дивизор Вейля, в
который особые точки не вносят вклада). Его класс в $\Cl X$ определён корректно. \emph{Канонический
пучок} $\omega_X=\cO_X(K_X)=j_*\Omega^n_U$, где $j\colon U\hookrightarrow X$.
Многообразие $X$ называется
\begin{itemize}
\item \emph{$\Q$-горенштейновым}, если $K_X$ --- $\Q$-картьеров дивизор; его индекс называется
\emph{индексом} $X$;
\item \emph{горенштейновым}, если $K_X$ --- дивизор Картье (и $X$ коэн-маколеево; для нормальных
поверхностей последнее выполнено автоматически).
\end{itemize}
\end{definition}

\begin{example}[Гиперповерхности и полные пересечения]\label{ex:hyp-Gor}
Пусть $X=V_+(F)\subset\PP^N$ --- нормальная гиперповерхность степени $d$ (возможно, особая).
Тогда $X$ горенштейнова и $K_X=\cO_X(d-N-1)$. Действительно, формула присоединения
$K_X=(K_{\PP^N}+X)|_X$ верна и для особых $X$: на гладкой части $X$ она доказывается как в
\cite[\S\,5]{AG} (вычет Пуанкаре $\omega=\frac{G\,dx_1\wedge\dots\wedge dx_N}{F}\mapsto
\frac{G\,dx_1\wedge\dots\wedge dx_{N-1}}{\partial F/\partial x_N}\big|_X$ даёт рациональную
$(N-1)$-форму на $X$, дивизор которой равен $(d-N-1)H|_X$ вне особых точек), а особые точки имеют
коразмерность $\ge2$ в $X$ и не влияют на класс Вейля. Аналогично для полных пересечений:
$K_X=\cO_X(\sum d_i-N-1)$. В частности, \emph{любая нормальная кубическая поверхность
$X\subset\PP^3$ горенштейнова и $-K_X=\cO_X(1)$} --- гиперплоское сечение. Это одна из причин,
по которым кубические поверхности --- удобный полигон для теории.
\end{example}

\begin{example}[Конусы]\label{ex:cone-canonical}
Пусть $V\subset\PP^N$ --- гладкое проективно нормальное многообразие, $L=\cO_V(1)$, и
$X=C(V,L)\subset\A^{N+1}$ --- аффинный конус. Тогда (см. \cite[3.1]{KolSing}) $K_X$ является
$\Q$-картьеровым тогда и только тогда, когда $K_V\sim_\Q rL$ для некоторого $r\in\Q$. Например,
для конуса над $\PP^1\times\PP^1$, вложенным по $\cO(1,2)$, имеем $K_V=\cO(-2,-2)\not\sim_\Q
r\cdot\cO(1,2)$, и $K_X$ не $\Q$-картьеров (это ответ на упражнение~2.1(6) обзора). Для конуса над
квадрикой из примера~\ref{ex:nonQCartier} ($L=\cO(1,1)$, $K_V=-2L$) дивизор $K_X$ --- Картье,
хотя $\Pi$ не $\Q$-Картье: горенштейновость не влечёт локальной факториальности.
\end{example}

\begin{example}[Взвешенные проективные пространства]
Для $X=\PP(a_0,\dots,a_n)$ (раздел~\ref{ss:wps}) $K_X=\cO_X(-\sum a_i)$; например,
$K_{\PP(1,1,n)}=\cO(-n-2)$. Пучок $\cO_X(k)$ обратим в вершине с весом $a_i$ тогда и только тогда,
когда $a_i\mid k$; поэтому индекс $\PP(1,1,n)$ равен $n/\gcd(n,2)$: $\PP(1,1,2)$ (квадратичный
конус) горенштейново, $\PP(1,1,3)$ имеет индекс 3, $\PP(1,1,4)$ --- индекс 2.
\end{example}

\subsection{Пары}

\begin{definition}\label{def:pair}
\emph{Пара} $(X,\Delta)$ --- нормальное многообразие $X$ вместе с эффективным $\Q$-дивизором
$\Delta=\sum b_jD_j$ (\emph{границей}), таким что $K_X+\Delta$ --- $\Q$-картьеров дивизор
(в обзоре это условие называется $\Q$-горенштейновостью пары; иногда его включают в определение
пары, иногда требуют отдельно). Пара \emph{проективна}, если проективно $X$. Пара $(X,\Delta)$
называется \emph{лог-Фано}, если она имеет klt-особенности (раздел~\ref{sec:mmp}) и
$-(K_X+\Delta)$ обилен; при $\Delta=0$ говорят о \emph{$\Q$-Фано многообразии}: нормальное
проективное $X$ с klt-особенностями и обильным $-K_X$.
\end{definition}

Пары появляются по нескольким причинам. (1) При ограничении на дивизор: если $E\subset Y$ ---
простой дивизор, то «формула присоединения» $K_E+\Delta_E=(K_Y+E)|_E$ в общем случае требует
поправки $\Delta_E$ (\emph{different}, см.~ниже). (2) При изучении семейств: по кривой $D\subset X$
и числу $c\in[0,1]$ получаем пару $(X,cD)$, и интересно, как меняются её свойства с ростом $c$
(примеры~\ref{ex:P112-pair}, \ref{ex:P114-pair}). (3) При разрешении особенностей:
$K_Y+\Delta_Y=\pi^*(K_X+\Delta)$ превращает любую пару в пару с «простыми» особенностями
(раздел~\ref{sec:mmp}).

\begin{definition}[Different]\label{def:different}
Пусть $(Y,E+\Delta)$ --- пара, $E$ --- простой дивизор с коэффициентом~1, не входящий в
$\Supp\Delta$, и $E^\nu\to E$ --- нормализация. Тогда существует единственный $\Q$-дивизор
$\Delta_E\ge0$ на $E^\nu$ (\emph{different}) такой, что
\[
K_{E^\nu}+\Delta_E=(K_Y+E+\Delta)|_{E^\nu}
\]
\cite[4.2]{KolSing}. Если $Y$ гладко, $E$ гладкий и $E+\Delta$ --- дивизор с простыми нормальными
пересечениями, то $\Delta_E=\Delta|_E$; вклад в $\Delta_E$ дают также особенности $Y$ вдоль $E$
(например, точка типа $\frac1n(1,a)$ на $Y$, лежащая на $E$, даёт слагаемое $(1-\frac1n)p$).
\end{definition}

\subsection{Бирациональные морфизмы и дивизоры над многообразием}\label{ss:divisors-over}

Пусть $\pi\colon Y\to X$ --- собственный бирациональный морфизм нормальных многообразий.
\emph{Исключительное множество} $\mathrm{Exc}(\pi)$ --- множество точек $Y$, в которых $\pi$ не
является локальным изоморфизмом; простой дивизор $E\subset Y$ \emph{исключителен}, если
$E\subset\mathrm{Exc}(\pi)$, т.\,е. $\codim\pi(E)\ge2$. \emph{Прямой образ} $\pi_*$ дивизоров
Вейля: $\pi_*E=0$ для исключительных $E$ и $\pi_*E=\overline{\pi(E)}$ иначе. \emph{Собственный
прообраз} $\tilde D=\pi^{-1}_*D$ простого дивизора $D\subset X$ --- замыкание $\pi^{-1}(D\cap U)$,
где $U$ --- открытое множество, над которым $\pi$ --- изоморфизм.

\begin{lemma}\label{lem:pullback-exc}
Если $D$ --- $\Q$-картьеров дивизор на $X$, то $\pi^*D=\tilde D+\sum_i c_iE_i$, где $E_i$ ---
исключительные дивизоры, $c_i\in\Q$, и $\pi_*\pi^*D=D$. Если $D$ эффективен, то все $c_i\ge0$.
\end{lemma}
\begin{proof}
Вне исключительного множества $\pi$ --- изоморфизм, поэтому $\pi^*D$ и $\tilde D$ совпадают вне
$\bigcup E_i$. Если $D\ge0$, то локальные уравнения $mD$ регулярны, их обратные образы регулярны,
и $\pi^*(mD)\ge0$.
\end{proof}

\begin{definition}[Дивизор над $X$]\label{def:divisor-over}
\emph{Простой дивизор над $X$} --- простой дивизор $E\subset Y$ на каком-нибудь нормальном
многообразии $Y$ с собственным бирациональным морфизмом $\pi\colon Y\to X$. Два таких дивизора
$E\subset Y$, $E'\subset Y'$ отождествляются, если они задают одно и то же дискретное нормирование
$\ord_E=\ord_{E'}$ поля $\C(X)=\C(Y)=\C(Y')$ (равносильно: бирациональное отображение
$Y\dashrightarrow Y'$ отображает $E$ на $E'$). \emph{Центр} $c_X(E)=\pi(E)\subset X$ ---
неприводимое замкнутое подмножество. Дивизор $E$ \emph{исключителен над $X$}, если $\codim c_X(E)\ge2$;
в противном случае $c_X(E)$ --- простой дивизор на $X$, и мы отождествляем $E$ с ним.
\end{definition}

Нормирование $\ord_E\colon\C(X)^*\to\Z$ называется \emph{дивизориальным нормированием}; в
разделе~\ref{sec:cubic} мы увидим, что в теории K-стабильности естественно рассматривать и более
общие нормирования, но для наших целей достаточно дивизориальных.

\begin{definition}\label{def:ord-E}
Пусть $E$ --- простой дивизор над $X$, реализованный на $\pi\colon Y\to X$.
\begin{enumerate}
\item Для $\Q$-картьерова дивизора $D$ на $X$ положим $\ord_E(D):=$ коэффициент $E$ в $\pi^*D$.
Если $D\ge0$ --- дивизор Картье с локальным уравнением $f$ около общей точки $c_X(E)$, то
$\ord_E(D)=\ord_E(\pi^*f)$ --- порядок обращения в нуль функции $\pi^*f$ вдоль $E$.
\item Для сечения $s\in H^0(X,L)$ линейного расслоения $L$ положим $\ord_E(s)=\ord_E(\divv s)$:
это порядок нуля сечения $\pi^*s$ вдоль $E$.
\item Для идеала $\fa\subset\cO_X$ положим $\ord_E(\fa)=\min\{\ord_E(f):f\in\fa\}$ (минимум по
функциям в окрестности общей точки центра).
\end{enumerate}
Все эти величины не зависят от выбора модели $Y$.
\end{definition}

\begin{example}[Раздутие точки]\label{ex:blowup-point-ord}
Пусть $p\in X$ --- гладкая точка $n$-мерного многообразия, $\pi\colon Y=\Bl_pX\to X$,
$E=\pi^{-1}(p)\cong\PP^{n-1}$. Тогда $c_X(E)=p$, и для функции $f$, регулярной в $p$,
$\ord_E(f)=\mult_pf$ --- кратность $f$ в точке $p$ (наименьшая степень членов ряда Тейлора). В
частности $\ord_E(\mm_p^k)=k$. Если $D\ni p$ --- гиперповерхность кратности $m$ в $p$, то
$\pi^*D=\tilde D+mE$ \cite[предложение~7.4]{AG}.
\end{example}

\begin{example}[Раздутие кривой]
Пусть $C\subset X$ --- гладкая кривая на гладком трёхмерном многообразии, $\pi\colon Y=\Bl_CX\to X$,
$E=\pi^{-1}(C)=\PP(N_{C/X})$ --- линейчатая поверхность над $C$. Тогда $c_X(E)=C$,
$\ord_E(f)=$ порядок обращения $f$ в нуль вдоль $C$ (в общей точке $C$), и для поверхности
$D\supset C$, гладкой в общей точке $C$, $\pi^*D=\tilde D+E$.
\end{example}

\begin{example}[Торические дивизоры]
Пусть $X=X_\Sigma$ --- торическое многообразие \cite[\S\,4]{AG}, $u\in N$ --- примитивный
вектор, не лежащий ни на одном луче $\Sigma$. Подразбиение веера, добавляющее луч $\rho=\R_{\ge0}u$,
даёт торический бирациональный морфизм $Y=X_{\Sigma'}\to X$ с исключительным дивизором
$E=D_\rho$. Его центр --- замыкание орбиты конуса $\sigma\in\Sigma$, в относительной внутренности
которого лежит $u$, а $\ord_E(\chi^m)=\langle m,u\rangle$ для характеров $\chi^m$ (это
лемма~4.12 из \cite{AG}, применённая к $Y$). Так, для $X=\A^2$ и $u=(1,1)$ получаем раздутие
начала координат; для $u=(a,b)$ --- \emph{взвешенное раздутие} с весами $(a,b)$:
$\ord_E(x^iy^j)=ai+bj$.
\end{example}

\begin{remark}[Почему «дивизоры над $X$», а не только дивизоры на $X$]
Многие инварианты K-стабильности (разделы~\ref{sec:beta}, \ref{sec:cubic}) --- это инфимумы по
\emph{всем} дивизорам над $X$. Причина в том, что интересующие нас свойства (например, тип
особенностей пары) не видны на самом $X$, а проявляются лишь после раздутий; множество всех
дивизоров над $X$ --- бирациональный инвариант, и формулировки в его терминах не зависят от
выбора модели.
\end{remark}

% ============================================================
\section{Пересечения, конусы, объём}\label{sec:cones}
% ============================================================

\subsection{Индексы пересечения и численная эквивалентность}

Для гладкой проективной поверхности индекс пересечения $D\cdot D'\in\Z$ определён в
\cite[теорема~0.21]{AG}. В общей ситуации нам понадобится следующее.

\begin{theorem}[{\cite[гл.~1]{Fulton}, \cite[I.1]{La}}]\label{thm:intersection-numbers}
Пусть $X$ --- проективное многообразие размерности $n$. Для дивизоров Картье $D_1,\dots,D_n$ на $X$
определено целое число $D_1\cdots D_n$ (\emph{индекс пересечения}), симметричное и полилинейное по
$D_i$, зависящее только от классов линейной эквивалентности и обладающее свойствами:
\begin{enumerate}
\item если $D_1,\dots,D_n$ --- гиперповерхности (эффективные дивизоры Картье), пересекающиеся
трансверсально в гладких точках, то $D_1\cdots D_n=\#(D_1\cap\dots\cap D_n)$;
\item (\emph{ограничение}) если $D_n=Y$ --- эффективный дивизор Картье и $Y$ неприводим, то
$D_1\cdots D_{n-1}\cdot Y=(D_1|_Y)\cdots(D_{n-1}|_Y)$, где справа --- индекс пересечения на $Y$
(если $Y$ приводим --- сумма по компонентам с кратностями);
\item (\emph{формула проекции}) для собственного сюръективного морфизма $f\colon Y\to X$
многообразий одной размерности $f^*D_1\cdots f^*D_n=\deg f\cdot D_1\cdots D_n$; для кривой
$C\subset Y$ и дивизора Картье $D$ на $X$: $f^*D\cdot C=D\cdot f_*C$, где $f_*C=\deg(C\to f(C))\,f(C)$,
если $f(C)$ --- кривая, и $f_*C=0$, если $f(C)$ --- точка;
\item для $\Q$-картьеровых дивизоров индексы пересечения определяются по линейности и принимают
значения в~$\Q$.
\end{enumerate}
Более общо, для замкнутого подмногообразия $V\subset X$ размерности $k$ определено число
$D_1\cdots D_k\cdot V$ (пересечение на $V$ ограничений). В частности, $D\cdot C=\deg(D|_C)$ для
кривых $C$.
\end{theorem}

Для поверхностей всё сводится к \cite[\S\,0.10]{AG}; например, на гладкой проективной поверхности
$D^2$ можно вычислять по формуле Римана--Роха, а $D\cdot C$ --- как степень ограничения.

\begin{definition}
Дивизоры Картье $D,D'$ \emph{численно эквивалентны}, $D\equiv D'$, если $D\cdot C=D'\cdot C$ для
всех кривых $C\subset X$; кривые $C,C'$ численно эквивалентны, если $D\cdot C=D\cdot C'$ для всех
$D$. Положим $N^1(X)=(\Pic X/{\equiv})\otimes\R$ и $N_1(X)=(\{\text{1-циклы}\}/{\equiv})\otimes\R$.
Это конечномерные двойственные друг другу векторные пространства (теорема Нерона--Севери);
$\rho(X)=\dim N^1(X)$ --- \emph{число Пикара}.
\end{definition}

Для многообразий Фано (и вообще при $H^1(X,\cO_X)=0$) численная и линейная эквивалентность
дивизоров совпадают с точностью до кручения, и $N^1(X)=\Pic X\otimes\R$. Для поверхностей дель Пеццо
$X_r=\Bl_r\PP^2$ имеем $N^1=\R L\oplus\bigoplus\R E_i$ \cite[следствие~7.7]{AG}.

\subsection{Нормальные поверхности: пересечения по Мамфорду}\label{ss:mumford}

На нормальной (особой) поверхности дивизоры Вейля не обязаны быть $\Q$-картьеровыми, и их всё же
хочется пересекать. Это возможно благодаря следующему фундаментальному факту.

\begin{theorem}[{Мамфорд, Грауэрт; \cite[1.1]{Mum61}, \cite[4.1]{KM}}]\label{thm:negdef}
Пусть $\pi\colon Y\to X$ --- разрешение особенностей нормальной поверхности (т.\,е. $Y$ гладкая
и $\pi$ --- собственный бирациональный морфизм), $E_1,\dots,E_r$ --- исключительные кривые над
точкой $x\in X$. Тогда матрица $(E_i\cdot E_j)$ отрицательно определена.
\end{theorem}

\begin{proof}[Идея]
Выберем функцию $f\in\mm_x$; тогда $\divv(\pi^*f)=\sum m_iE_i+(\text{собственный прообраз})$ с
$m_i>0$, и $\divv(\pi^*f)\cdot E_j=0$ (главный дивизор), откуда $(\sum m_iE_i)\cdot E_j\le0$ для
всех $j$, причём строго для некоторых. Из этого и связности $\bigcup E_i$ по лемме об отрицательно
определённых матрицах («если $Z=\sum m_iE_i$, $m_i>0$, и $Z\cdot E_j\le0$ при всех $j$, причём
хотя бы для одного строго, а $E_i\cdot E_j\ge0$ при $i\ne j$, то форма отрицательно определена»)
получаем утверждение.
\end{proof}

\begin{definition}[Обратный образ по Мамфорду]\label{def:mumford-pullback}
В условиях теоремы для дивизора Вейля $D$ на $X$ положим
\[
\pi^*D:=\tilde D+\sum_ic_iE_i,\qquad\text{где } c_i\in\Q \text{ определены условиями }
(\pi^*D)\cdot E_j=0\ \ \forall j
\]
(система невырождена в силу отрицательной определённости). Для дивизоров Вейля $D,D'$ положим
$D\cdot D':=\pi^*D\cdot\pi^*D'\in\Q$. Это не зависит от выбора разрешения; если $D$ --- дивизор
Картье, то $\pi^*D$ --- обычный обратный образ (у него нулевые пересечения с $E_j$ по формуле
проекции), и для $\Q$-картьеровых дивизоров определение согласуется с предыдущим пунктом.
\end{definition}

\begin{example}[Конус над рациональной нормальной кривой]\label{ex:P11n-intersections}
Пусть $X=\PP(1,1,n)$ --- проективный конус над рациональной нормальной кривой степени $n$
(раздел~\ref{ss:wps}), $\pi\colon Y=\mathbb F_n\to X$ --- стягивание отрицательного сечения
$E$ ($E^2=-n$) поверхности Хирцебруха; $F$ --- слой $\mathbb F_n\to\PP^1$. Прямая $\ell\subset X$
(образующая конуса) --- дивизор Вейля, $\tilde\ell=F$, и $\pi^*\ell=F+cE$ с
$0=(F+cE)\cdot E=1-cn$, т.\,е. $\pi^*\ell=F+\frac1nE$, $\ell^2=(F+\frac1nE)^2=\frac2n-\frac1n=\frac1n$.
Гиперплоское сечение $H_\infty$ (сечение «на бесконечности», не проходящее через вершину) ---
дивизор Картье класса $E+nF$, $H_\infty^2=n$, $H_\infty\sim n\ell$. Антиканонический дивизор
$-K_X=(n+2)\ell$ (пример~\ref{ex:K-wps}) имеет $(-K_X)^2=(n+2)^2/n$; при $n=2$ получаем $8$
(вырожденная квадрика в $\PP^3$ имеет $K^2=8$, как и гладкая).
\end{example}

\subsection{Конусы кривых и дивизоров}

\begin{definition}
\emph{Конус кривых} $NE(X)\subset N_1(X)$ --- конус, порождённый классами неприводимых кривых;
$\NE(X)$ --- его замыкание (\emph{конус Мори}). Дивизор Картье $D$ \emph{численно эффективен}
(\emph{nef}), если $D\cdot C\ge0$ для всех кривых $C$; конус nef-дивизоров $\Nef(X)\subset N^1(X)$
двойствен к $\NE(X)$.
\end{definition}

\begin{theorem}[{Клейман, \cite[1.4.23]{La}}]\label{thm:Kleiman}
Дивизор Картье $D$ обилен тогда и только тогда, когда $D\cdot z>0$ для всех $z\in\NE(X)\setminus0$.
Иначе говоря, $\mathrm{Amp}(X)$ --- внутренность $\Nef(X)$, а $\Nef(X)$ --- замыкание обильного
конуса. Nef-дивизор $D$ удовлетворяет $D^{\dim V}\cdot V\ge0$ для всех подмногообразий $V$.
\end{theorem}

\begin{theorem}[{Теорема о конусе Мори; \cite[1.24, 3.7]{KM}}]\label{thm:cone}
Пусть $X$ --- проективное многообразие с klt-особенностями (в частности, гладкое) и $-K_X$ обилен.
Тогда конус $\NE(X)$ --- рациональный полиэдральный конус, порождённый конечным числом классов
рациональных кривых.
\end{theorem}

\begin{example}\label{ex:cones}
\begin{enumerate}
\item $\PP^2$: $N_1=\R$, конус --- луч, порождённый прямой. $\PP^1\times\PP^1$: $\NE=\R_{\ge0}F_1+\R_{\ge0}F_2$
(слои двух проекций).
\item $\mathbb F_1=\Bl_p\PP^2$: $\NE=\R_{\ge0}E+\R_{\ge0}(L-E)$, где $L-E$ --- класс слоя
проекции из $p$ \cite[пример~6.3(3)]{AG}. Общее: $\NE(\mathbb F_n)=\R_{\ge0}E+\R_{\ge0}F$.
\item $\Bl_2\PP^2$ (дель Пеццо степени 7): конус порождён тремя $(-1)$-кривыми $E_1$, $E_2$,
$L-E_1-E_2$.
\item $\Bl_3\PP^2$ (степень 6): конус порождён шестью $(-1)$-кривыми $E_i$, $L-E_i-E_j$,
образующими шестиугольник.
\item Для поверхности дель Пеццо степени $d\le7$ конус $\NE$ порождён $(-1)$-кривыми
\cite[8.2]{Do}; при $d=3$ --- 27 прямыми.
\end{enumerate}
Во всех случаях, чтобы проверить, является ли дивизор $D$ nef или обильным, достаточно вычислить
$D\cdot C$ для образующих конуса.
\end{example}

\subsection{Большие дивизоры, псевдоэффективность и объём}\label{ss:volume}

\begin{definition}\label{def:big}
Пусть $X$ --- проективное многообразие размерности $n$, $D$ --- дивизор Картье (или $\Q$-дивизор
Вейля на нормальном $X$). \emph{Объём} $D$ ---
\[
\vol(D)=\limsup_{m\to\infty}\frac{h^0(X,mD)}{m^n/n!}
\]
(предел по тем $m$, для которых $mD$ целочислен). $D$ называется \emph{большим}, если
$\vol(D)>0$. \emph{Псевдоэффективный конус} $\Eff(X)\subset N^1(X)$ --- замыкание конуса,
порождённого эффективными дивизорами.
\end{definition}

\begin{theorem}[{\cite[2.2]{La}}]\label{thm:big}
Для дивизора Картье $D$ на проективном многообразии равносильны:
\begin{enumerate}
\item $D$ большой, т.\,е. $\vol(D)>0$;
\item $h^0(mD)\ge c\,m^n$ при $m\gg0$ для некоторого $c>0$;
\item при $m\gg0$ рациональное отображение $\varphi_{|mD|}\colon X\dashrightarrow\PP^N$ бирационально
на образ;
\item (\emph{лемма Кодаиры}) $mD\sim A+N$ для некоторого $m>0$, обильного $A$ и эффективного $N$.
\end{enumerate}
Кроме того: большие дивизоры образуют открытый конус в $N^1(X)$, замыкание которого --- $\Eff(X)$;
объём --- непрерывная функция на $N^1(X)$, однородная степени $n$ ($\vol(tD)=t^n\vol(D)$,
$t\in\Q_{>0}$), и $\vol$ является $\limsup$ по всем $m$, а не только по подпоследовательности
\cite[2.2.C]{La}; $\vol(D)\le\vol(D')$, если $D'-D$ эффективен; $\vol(f^*D)=\vol(D)$ для
бирационального $f$.
\end{theorem}

\begin{proposition}\label{prop:vol-nef}
Если $D$ nef, то $\vol(D)=D^n$. В частности, для обильного $D$ объём равен $D^n>0$.
\end{proposition}
\begin{proof}[Набросок]
Для обильного $D$ это асимптотическая формула Римана--Роха (теорема~\ref{thm:aRR}). Для nef $D$
берём обильный $A$: $D+\epsilon A$ обилен при всех рациональных $\epsilon>0$, и
$\vol(D)=\lim_{\epsilon\to0}(D+\epsilon A)^n=D^n$ по непрерывности объёма.
\end{proof}

Главная для нас операция --- вычисление $\vol(L-tE)$ как функции от $t\ge0$, где $L$ обилен и $E$
--- эффективный дивизор (или дивизор над $X$). Общие свойства: функция $t\mapsto\vol(L-tE)$
непрерывна, убывает, равна $(L-tE)^n$, пока $L-tE$ nef, и равна нулю при $t\ge\tau$, где
\[
\tau=\tau(L;E)=\sup\{t\ge0: L-tE\ \text{псевдоэффективен}\}
\]
--- \emph{псевдоэффективный порог}. Число $\tau$ конечно: $L-tE$ не может быть псевдоэффективным
при $t\gg0$, так как $(L-tE)\cdot A^{n-1}<0$ для обильного $A$. Между nef-порогом и
псевдоэффективным порогом объём вычисляется труднее; для поверхностей есть полный ответ.

\begin{theorem}[{Разложение Зарисского; \cite[2.3.19]{La}, \cite{Zar}}]\label{thm:Zariski}
Пусть $D$ --- псевдоэффективный $\Q$-дивизор на гладкой проективной поверхности $X$. Существует
единственное разложение $D=P+N$ в сумму $\Q$-дивизоров, где
\begin{enumerate}
\item $P$ nef (\emph{положительная часть});
\item $N=\sum a_iN_i\ge0$ (\emph{отрицательная часть}), причём либо $N=0$, либо матрица
$(N_i\cdot N_j)$ отрицательно определена;
\item $P\cdot N_i=0$ для всех $i$.
\end{enumerate}
При этом $H^0(X,mD)=H^0(X,mP)$ для всех $m$ (т.\,е. $N$ --- «фиксированная часть»), так что
$\vol(D)=\vol(P)=P^2$. То же верно для нормальных проективных поверхностей с пересечениями по
Мамфорду.
\end{theorem}

\emph{Как вычислять.} Если $D$ nef, то $P=D$, $N=0$. Если $D\cdot C<0$ для неприводимой кривой
$C$, то (так как $D$ псевдоэффективен) $C^2<0$, и $C$ обязана входить в $N$. Алгоритм: берём
все кривые $C_1,\dots,C_k$ с $D\cdot C_i<0$, ищем $N=\sum a_iC_i$ из условий $(D-N)\cdot C_j=0$;
если $D-N$ nef --- готово; если нет --- добавляем к набору новые кривые, на которых $D-N$
отрицателен, и повторяем (процесс конечен). Для поверхностей дель Пеццо кривые с отрицательным
квадратом --- это $(-1)$-кривые (а на особых --- также $(-2)$-кривые разрешения), и всё сводится
к перебору.

\begin{example}[$\vol(-K-tE)$ на $\mathbb F_1$]\label{ex:vol-F1}
Пусть $X=\mathbb F_1=\Bl_p\PP^2$, $E$ --- исключительная кривая, $-K_X=3L-E$,
$D_t=-K_X-tE=3L-(1+t)E$. Имеем $D_t\cdot E=1+t>0$, $D_t\cdot(L-E)=3-(1+t)=2-t$. По
примеру~\ref{ex:cones}(2) дивизор $D_t$ nef при $0\le t\le2$, и $\vol(D_t)=D_t^2=9-(1+t)^2$.
При $t>2$ дивизор $D_t$ отрицателен на классе слоя $F=L-E$, который nef ($F^2=0$,
$F$ --- неприводимая кривая); значит, $D_t$ не псевдоэффективен (псевдоэффективный дивизор
неотрицателен на nef-кривых). Итак, $\tau=2$ и
\[
\vol(-K_X-tE)=\begin{cases}9-(1+t)^2,&0\le t\le2,\\0,&t\ge2.\end{cases}
\]
Заметим, что $\vol(D_2)=0$: на пороге объём обращается в нуль, так как $D_2=2F$ --- класс слоя.
\end{example}

\begin{example}[$\vol(-K-tE_1)$ на поверхности дель Пеццо степени 7]\label{ex:vol-dP7}
Пусть $X=\Bl_{p_1,p_2}\PP^2$, $-K=3L-E_1-E_2$, $D_t=-K-tE_1=3L-(1+t)E_1-E_2$. Образующие конуса:
$E_1$, $E_2$, $\ell=L-E_1-E_2$. Пересечения: $D_t\cdot E_1=1+t$, $D_t\cdot E_2=1$,
$D_t\cdot\ell=3-(1+t)-1=1-t$. Итак, $D_t$ nef при $t\le1$ и $\vol(D_t)=D_t^2=9-(1+t)^2-1=8-(1+t)^2$.
При $t>1$ имеем $D_t\cdot\ell<0$, $\ell^2=-1$: отрицательная часть $N=a\ell$, где
$(D_t-a\ell)\cdot\ell=(1-t)+a=0$, т.\,е. $a=t-1$. Тогда
\[
P_t=D_t-(t-1)(L-E_1-E_2)=(4-t)L-2E_1-(2-t)E_2 .
\]
Проверка: $P_t\cdot E_1=2$, $P_t\cdot E_2=2-t$, $P_t\cdot\ell=0$ --- дивизор $P_t$ nef при
$1\le t\le2$, и $\vol(D_t)=P_t^2=(4-t)^2-4-(2-t)^2=8-4t$. При $t=2$ объём равен нулю
($P_2=2(L-E_1)$ --- класс слоя), и при $t>2$ дивизор не псевдоэффективен. Итого
\[
\vol(-K_X-tE_1)=\begin{cases}8-(1+t)^2,&0\le t\le1,\\8-4t,&1\le t\le2,\\0,&t\ge2.\end{cases}
\]
Функция непрерывна (при $t=1$ оба выражения дают $4$), но её вторая производная скачком меняется в
$t=1$: это типичное поведение объёма --- кусочно-полиномиальная функция, «ломающаяся» там, где
меняется разложение Зарисского.
\end{example}

\subsection{Асимптотическая формула Римана--Роха}

\begin{theorem}[{\cite[1.1.24]{La}, \cite[2.1.2]{Xu}}]\label{thm:aRR}
Пусть $X$ --- нормальное проективное многообразие размерности $n$, $L$ --- обильный дивизор
Картье. Тогда при $m\gg0$
\[
h^0(X,mL)=\chi(X,mL)=\frac{L^n}{n!}\,m^n-\frac{L^{n-1}\cdot K_X}{2\,(n-1)!}\,m^{n-1}+O(m^{n-2}),
\]
где $L^{n-1}\cdot K_X$ --- индекс пересечения с (вейлевым) каноническим дивизором
(для гладкого $X$ это теорема Хирцебруха--Римана--Роха вместе с теоремой Серра об обращении в нуль;
для нормального $X$ --- та же формула, поскольку $K_X$ и канонический класс гладкой части
отличаются лишь в коразмерности~2). Функция $m\mapsto\chi(X,mL)$ --- многочлен
(\emph{многочлен Гильберта} пары $(X,L)$).
\end{theorem}

\begin{example}
\begin{enumerate}
\item $\PP^n$, $L=H$: $h^0(mH)=\binom{m+n}{n}=\frac{m^n}{n!}+\frac{(n+1)m^{n-1}}{2(n-1)!}+\dots$,
что согласуется с $K=-(n+1)H$, $H^n=1$.
\item Поверхность, $L$ обилен: $\chi(mL)=\chi(\cO_X)+\frac{m^2L^2}2-\frac{mL\cdot K}2$
(\cite[теорема~0.22]{AG}). Для кубической поверхности $L=-K$: $h^0(-mK)=1+\frac{3m^2}2+\frac{3m}2$;
при $m=1$ получаем $4$ --- размерность пространства линейных форм на $\PP^3$, как и должно быть.
\item Для $L=-K_X$ на $\Q$-Фано многообразии: $h^0(-mK_X)=\frac{(-K_X)^n}{n!}m^n+\frac{(-K_X)^n}{2(n-1)!}m^{n-1}+\dots$
--- второй коэффициент выражается через тот же объём $(-K_X)^n$; именно это упрощает формулу для
инварианта Футаки в разделе~\ref{sec:tc}.
\end{enumerate}
\end{example}

% ============================================================
\section{Раздутия, разрешения и особенности поверхностей}\label{sec:sing}
% ============================================================

\subsection{Раздутия и канонический класс}

\begin{theorem}[{\cite[II.8.24, II, упр.~8.5]{Ha}, \cite[7.6]{AG}}]\label{thm:blowup-K}
Пусть $X$ --- гладкое многообразие размерности $n$, $Z\subset X$ --- гладкое замкнутое
подмногообразие коразмерности $c\ge2$, $\pi\colon Y=\Bl_ZX\to X$ --- раздутие с исключительным
дивизором $E=\pi^{-1}(Z)$. Тогда
\begin{enumerate}
\item $Y$ гладко, $E\cong\PP(N_{Z/X})$ --- проективизация нормального расслоения (расслоение на
$\PP^{c-1}$ над $Z$), и $\cO_Y(E)|_E\cong\cO_{\PP(N)}(-1)$ --- тавтологическое расслоение;
\item $K_Y=\pi^*K_X+(c-1)E$;
\item $\Pic Y=\pi^*\Pic X\oplus\Z E$.
\end{enumerate}
\end{theorem}

\begin{proof}[Набросок доказательства (2)]
Вопрос локален около общей точки $E$; выберем локальные координаты $x_1,\dots,x_n$ на $X$, в
которых $Z=\{x_1=\dots=x_c=0\}$. В карте раздутия $x_1=y_1$, $x_i=y_1y_i$ ($2\le i\le c$),
$x_j=y_j$ ($j>c$) имеем $E=\{y_1=0\}$ и
\[
\pi^*(dx_1\wedge\dots\wedge dx_n)=y_1^{\,c-1}\,dy_1\wedge\dots\wedge dy_n ,
\]
так что форма $\pi^*\omega$ обращается в нуль вдоль $E$ с порядком $c-1$.
\end{proof}

\begin{example}
Для точки на поверхности ($c=2$): $K_Y=\pi^*K_X+E$, $E^2=-1$ \cite[7.6]{AG}. Для точки на
трёхмерном многообразии ($c=3$): $K_Y=\pi^*K_X+2E$, $E\cong\PP^2$, $\cO_E(E)=\cO(-1)$, откуда
$E^3=(E|_E)^2=1$, $E^2\cdot\pi^*D=0$ и $E\cdot\pi^*D\cdot\pi^*D'=0$. Для гладкой кривой $C$ на
трёхмерном многообразии: $K_Y=\pi^*K_X+E$, $E^3=-\deg N_{C/X}$, $E^2\cdot\pi^*D=-D\cdot C$,
$E\cdot\pi^*D\cdot\pi^*D'=0$ (\cite[IV.6]{Fulton}; см. пример~\ref{ex:Bl-cubic-P3}).
\end{example}

В терминах определения~\ref{def:ord-E} теорема говорит, что \emph{лог-дискрепантность}
исключительного дивизора раздутия гладкого центра коразмерности $c$ равна $c$: в разделе~\ref{sec:mmp}
мы определим $A_X(E)=1+\ord_E(K_Y-\pi^*K_X)$, и здесь $A_X(E)=1+(c-1)=c$.

\subsection{Разрешения особенностей}

\begin{theorem}[Хиронака]\label{thm:Hironaka}
Пусть $X$ --- многообразие, $D\subset X$ --- замкнутое подмножество. Существует проективный
бирациональный морфизм $\pi\colon Y\to X$ (композиция раздутий гладких центров), такой что $Y$
гладко, $\pi$ --- изоморфизм над $X\setminus(\Sing X\cup D)$, и $\pi^{-1}(D)\cup\mathrm{Exc}(\pi)$
--- дивизор с \emph{простыми нормальными пересечениями} (snc): его компоненты гладки и в каждой
точке локально выглядят как объединение координатных гиперплоскостей.
\end{theorem}

Такой $\pi$ называется \emph{лог-разрешением} пары $(X,D)$ (или $(X,\Delta)$ с $\Supp\Delta=D$).
Лог-разрешения не единственны, но любые два доминируются третьим.

\begin{definition}
Для нормальной поверхности $X$ разрешение $\pi\colon Y\to X$ \emph{минимально}, если
исключительное множество не содержит $(-1)$-кривых. Минимальное разрешение существует и единственно
(любое разрешение получается из минимального раздутиями точек; теорема Кастельнуово
\cite[7.8]{AG}), но оно, вообще говоря, не является лог-разрешением (исключительные кривые могут
пересекаться не трансверсально или быть особыми).
\end{definition}

\begin{lemma}[Отрицательность и минимальные разрешения]\label{lem:minimal-discrep}
Пусть $\pi\colon Y\to X$ --- минимальное разрешение нормальной поверхности, $X$ $\Q$-горенштейнова,
$K_Y=\pi^*K_X+\sum a_iE_i$. Тогда все $a_i\le0$.
\end{lemma}
\begin{proof}
По формуле присоединения \cite[5.8]{AG} $K_Y\cdot E_i=2p_a(E_i)-2-E_i^2$. Если $E_i^2\le-2$ или
$p_a(E_i)\ge1$, то $K_Y\cdot E_i\ge0$; случай $E_i^2=-1$, $p_a=0$ исключён минимальностью. Итак,
$B:=\sum a_iE_i=K_Y-\pi^*K_X$ удовлетворяет $B\cdot E_j=K_Y\cdot E_j\ge0$ для всех $j$
(так как $\pi^*K_X\cdot E_j=0$). Запишем $B=B^+-B^-$ с $B^\pm\ge0$ без общих компонент. Если
$B^+\ne0$, то по отрицательной определённости (теорема~\ref{thm:negdef}) $(B^+)^2<0$, и
$0\le B\cdot B^+=(B^+)^2-B^-\cdot B^+\le(B^+)^2<0$ --- противоречие. Значит, $B^+=0$, $B\le0$.
\end{proof}

Это частный случай \emph{леммы об отрицательности} \cite[3.39]{KM}: если $\pi\colon Y\to X$
собственный бирациональный, $B$ --- $\pi$-исключительный $\Q$-дивизор и $-B$ $\pi$-nef, то $B\ge0$.

\subsection{Взвешенные проективные пространства}\label{ss:wps}

\begin{definition}
Пусть $a_0,\dots,a_n$ --- положительные целые числа. \emph{Взвешенное проективное пространство}
\[
\PP(a_0,\dots,a_n)=(\C^{n+1}\setminus0)/\C^*,\qquad \lambda\cdot(x_0,\dots,x_n)=(\lambda^{a_0}x_0,\dots,\lambda^{a_n}x_n),
\]
т.\,е. $\Proj\C[x_0,\dots,x_n]$ с градуировкой $\deg x_i=a_i$. Набор весов \emph{нормализован}
(\emph{well-formed}), если любые $n$ из весов взаимно просты; всякое $\PP(a_0,\dots,a_n)$
изоморфно нормализованному (\cite[1.3]{Dolg-WPS}; например, $\PP(1,2,2)\cong\PP(1,1,1)$), и далее
веса предполагаются нормализованными.
\end{definition}

\begin{proposition}\label{prop:wps}
Пусть $X=\PP(a_0,\dots,a_n)$ нормализовано.
\begin{enumerate}
\item $X$ --- нормальное проективное торическое многообразие размерности $n$, покрытое картами
$U_i=\{x_i\ne0\}\cong\A^n/\mu_{a_i}$, где $\mu_{a_i}$ действует на координатах $x_j$, $j\ne i$,
с весами $a_j$. Особые точки $X$ --- это точки с нетривиальными стабилизаторами; в частности,
«вершина» $P_i=[0:\dots:1:\dots:0]$ --- факторособенность $\frac1{a_i}(a_0,\dots,\widehat{a_i},\dots,a_n)$.
\item Для $k\in\Z$ пучок $\cO_X(k)$ (сечения над $U$ --- рациональные функции степени $k$,
регулярные на конусе над $U$) рефлексивен ранга 1; $H^0(X,\cO_X(k))$ --- однородные многочлены
степени $k$; $\cO_X(k)$ обратим в точке $P_i$ тогда и только тогда, когда $a_i\mid k$; все
дивизоры Вейля на $X$ $\Q$-картьеровы, $\Cl X=\Z\cdot\cO_X(1)$.
\item $\cO_X(1)^n=\dfrac1{a_0a_1\cdots a_n}$ (индекс пересечения по Мамфорду, см.~\S\ref{ss:mumford}).
\item $K_X=\cO_X(-a_0-\dots-a_n)$, поэтому $(-K_X)^n=\dfrac{(a_0+\dots+a_n)^n}{a_0\cdots a_n}$.
\end{enumerate}
\end{proposition}
\begin{proof}[Пояснения]
(1) В карте $x_i\ne0$ нормируем $x_i=1$; остаётся действие стабилизатора $\mu_{a_i}=\{\lambda:
\lambda^{a_i}=1\}$ на остальных координатах. (2) См.~\cite[\S\,1]{Dolg-WPS}. (3) Отображение
$q\colon\PP^n\to X$, $[y_0:\dots:y_n]\mapsto[y_0^{a_0}:\dots:y_n^{a_n}]$, --- конечный морфизм
степени $a_0\cdots a_n$ (фактор по $\mu_{a_0}\times\dots\times\mu_{a_n}$), и $q^*\cO_X(k)=\cO_{\PP^n}(k)$:
взвешенно-однородный многочлен степени $k$ от $x$ переходит в однородный многочлен степени $k$ от
$y$. По формуле проекции $\cO_X(1)^n=\cO_{\PP^n}(1)^n/\deg q=1/(a_0\cdots a_n)$. (4) Торическая
формула $K=-\sum D_\rho$ \cite[4.20]{AG}: $X$ --- торическое многообразие с $n+1$ инвариантными
дивизорами $D_i=\{x_i=0\}$ класса $\cO(a_i)$, так что $K_X=-\sum\cO(a_i)=\cO(-\sum a_i)$.
\end{proof}

\begin{example}[$K_X$ взвешенного проективного пространства]\label{ex:K-wps}
$K_{\PP(1,1,n)}=\cO(-n-2)$ и $(-K)^2=(n+2)^2/n$; $K_{\PP(1,1,2)}=\cO(-4)$, $(-K)^2=8$;
$K_{\PP(1,2,3)}=\cO(-6)$, $(-K)^2=6$; $K_{\PP(1,1,4)}=\cO(-6)$, $(-K)^2=9$;
$K_{\PP(1,4,25)}=\cO(-30)$, $(-K)^2=900/100=9$. Последние два примера --- не совпадение, см.
пример~\ref{ex:Markov}.
\end{example}

\begin{example}[$\PP(1,1,2)$ --- квадратичный конус]\label{ex:P112}
Сечения $\cO(2)$ на $\PP(1,1,2)$ с координатами $x,y$ (степени 1) и $z$ (степени 2) --- это
$x^2,xy,y^2,z$. Отображение $[x:y:z]\mapsto[x^2:xy:y^2:z]=[u_0:u_1:u_2:u_3]$ --- вложение
$\PP(1,1,2)\hookrightarrow\PP^3$ с образом $\{u_0u_2=u_1^2\}$ --- квадратичным конусом с
вершиной $[0:0:0:1]$, образом точки $P_2=[0:0:1]$. Это решение упражнения~3.4(10a) обзора.
Особенность в вершине --- $\frac12(1,1)=A_1$ (пример~\ref{ex:cone-divisors}). Прямые
$\{x=0\}$ и т.\,п. --- образующие конуса, дивизоры класса $\cO(1)$; сечение $Q=\{z=0\}\cong\PP^1$
--- гладкая коника, не проходящая через вершину, класса $\cO(2)$ (дивизор Картье).
\end{example}

\begin{example}[$\PP(1,1,n)$ --- конус над рациональной нормальной кривой]\label{ex:P11n}
Аналогично сечения $\cO(n)$ на $\PP(1,1,n)$ --- это $x^n,x^{n-1}y,\dots,y^n,z$, и
$|\cO(n)|$ вкладывает $\PP(1,1,n)$ в $\PP^{n+1}$ как проективный конус над рациональной
нормальной кривой степени $n$ (образом $\PP^1$ при $[x:y]\mapsto[x^n:\dots:y^n]$) с вершиной
$[0:\dots:0:1]$. Это решение упражнения~6.3(10c) обзора. Аффинный конус над рациональной
нормальной кривой --- это $\Spec\bigoplus_kH^0(\PP^1,\cO(nk))=\Spec\C[x,y]^{\mu_n}$, где $\mu_n$
действует весами $(1,1)$: инварианты --- мономы степени, кратной $n$. Поэтому вершина --- особенность
$\frac1n(1,1)$ (упражнение~6.3(10a)). Разрешение: раздутие вершины даёт поверхность Хирцебруха
$\mathbb F_n$, исключительная кривая --- сечение $E$ с $E^2=-n$ (упражнение~6.3(10b));
пересечения вычислены в примере~\ref{ex:P11n-intersections}.
\end{example}

\subsection{Циклические факторособенности и особенности Дю Валя}\label{ss:quotient}

\begin{definition}\label{def:cyclic-quot}
Пусть $n\ge1$, $a,b\in\Z$. Особенность $\frac1n(a,b)$ --- это $0\in\A^2/\mu_n$, где
примитивный корень $\zeta=\zeta_n$ действует по правилу $\zeta\cdot(x,y)=(\zeta^ax,\zeta^by)$.
Более общо, \emph{факторособенность} --- $0\in\A^m/G$ для конечной подгруппы $G\subset\GL_m(\C)$;
не умаляя общности, можно считать, что $G$ \emph{мала} (не содержит псевдоотражений, т.\,е.
действует свободно в коразмерности~1): по теореме Шевалле--Шепарда--Тодда фактор по группе,
порождённой псевдоотражениями, гладок. Для $\frac1n(a,b)$ это условие означает $\gcd(a,n)=\gcd(b,n)=1$,
и тогда $\frac1n(a,b)\cong\frac1n(1,q)$, $q\equiv ba^{-1}\pmod n$.
\end{definition}

Координатное кольцо $\A^2/\mu_n$ --- кольцо инвариантов $\C[x,y]^{\mu_n}$ (упражнение~6.3(8)
обзора); его образующие --- инвариантные мономы.

\begin{example}[$A_{n-1}$ и $\frac1n(1,1)$]\label{ex:An-invariants}
\begin{enumerate}
\item $\frac1n(1,-1)$: инвариантные мономы $x^iy^j$ с $i\equiv j\pmod n$ порождаются
$u=x^n$, $v=y^n$, $w=xy$ с единственным соотношением $uv=w^n$. Итак,
$\frac1n(1,-1)=\{uv=w^n\}\subset\A^3$ --- гиперповерхностная особенность, следовательно,
\emph{горенштейнова} (локальное кольцо гиперповерхности --- полное пересечение). Это особенность
$A_{n-1}$, и это решает упражнение~6.3(8).
\item $\frac1n(1,1)$: инварианты --- мономы степени, кратной $n$; порождающие
$x^n,x^{n-1}y,\dots,y^n$ с соотношениями $2\times2$-миноров матрицы
$\begin{pmatrix}x^n&x^{n-1}y&\dots&xy^{n-1}\\x^{n-1}y&x^{n-2}y^2&\dots&y^n\end{pmatrix}$ ---
конус над рациональной нормальной кривой (пример~\ref{ex:P11n}). При $n\ge3$ он
\emph{не} горенштейнов: $K$ имеет индекс $n/\gcd(n,2)>1$.
\end{enumerate}
\end{example}

\begin{proposition}[Торическое описание и разрешение]\label{prop:HJ}
Особенность $\frac1n(1,q)$ ($0<q<n$, $\gcd(q,n)=1$) --- аффинное торическое многообразие
$U_\sigma$ конуса $\sigma=\Cone(e_2,\ ne_1-qe_2)\subset\R^2$. Её минимальное разрешение торическое:
лучи, добавляемые к $\sigma$, --- это $v_1,\dots,v_r$, где
\[
\frac nq=b_1-\cfrac1{b_2-\cfrac1{\ddots-\cfrac1{b_r}}}\qquad(b_i\ge2)
\]
--- \emph{цепная дробь Хирцебруха--Юнга}, а исключительное множество --- цепочка гладких
рациональных кривых $E_1,\dots,E_r$ с $E_i^2=-b_i$, $E_i\cdot E_{i+1}=1$.
\end{proposition}
\begin{proof}[Пояснения и примеры]
Для $\frac12(1,1)$ это пример~4.7(2) из \cite{AG}. Общее утверждение --- \cite[10.2]{CLS}.
\begin{itemize}
\item $q=1$: $n/1=n$, одна кривая с $E^2=-n$: конус над рациональной нормальной кривой
(пример~\ref{ex:P11n}).
\item $q=n-1$ (особенность $A_{n-1}$): $\frac n{n-1}=2-\cfrac1{2-\cfrac1{\ddots}}$ ($n-1$ двоек),
т.\,е. цепочка из $n-1$ кривых с самопересечением $-2$. Явно: лучи $v_k=(k,1-k)$, $k=0,\dots,n$,
где $v_0=e_2$, $v_n=ne_1-(n-1)e_2$; имеем $v_{k-1}+v_{k+1}=2v_k$, так что по
\cite[предложение~4.18]{AG} $E_k^2=-2$, а каждый конус $\Cone(v_{k},v_{k+1})$ унимодулярен
($\det\begin{pmatrix}k&k+1\\1-k&-k\end{pmatrix}=-1$), т.\,е. разрешение гладко. Это решает
упражнение~6.3(9) обзора.
\item $\frac14(1,1)$: одна $(-4)$-кривая; $\frac19(1,2)$: $9/2=5-\frac12$, цепочка $(-5),(-2)$;
$\frac1{25}(1,4)$: $25/4=7-\frac1{4}$... точнее $25/4=7-\cfrac1{4/3}$, $4/3=2-\frac12$, т.\,е.
цепочка $(-7),(-2),(-2)$.
\end{itemize}
\end{proof}

\begin{definition}[Особенности Дю Валя]\label{def:DuVal}
Нормальная особенность поверхности $0\in X$ называется \emph{особенностью Дю Валя} (рациональной
двойной точкой, простой особенностью, $ADE$-особенностью), если выполнено одно из равносильных
условий \cite{Durfee}, \cite[4.20]{KM}:
\begin{enumerate}
\item $X\cong\A^2/G$ для конечной подгруппы $G\subset\SL_2(\C)$;
\item $X$ аналитически изоморфна одной из гиперповерхностей
\[
A_n:\ x^2+y^2+z^{n+1}=0\ (n\ge1),\quad D_n:\ x^2+y^2z+z^{n-1}=0\ (n\ge4),\quad
E_6:\ x^2+y^3+z^4,\quad E_7:\ x^2+y^3+yz^3,\quad E_8:\ x^2+y^3+z^5;
\]
\item исключительные кривые минимального разрешения --- гладкие рациональные $(-2)$-кривые,
пересекающиеся трансверсально по схеме Дынкина соответствующего типа;
\item (\emph{канонические} особенности) $K_Y=\pi^*K_X$ для минимального разрешения $\pi$.
\end{enumerate}
Группы $G\subset\SL_2$: циклическая $\mu_{n+1}$ ($A_n$), бинарные диэдральная ($D_n$),
тетраэдральная ($E_6$), октаэдральная ($E_7$), икосаэдральная ($E_8$).
\end{definition}

Особенность $A_n$ в координатах $uv=w^{n+1}$ --- это $\frac1{n+1}(1,-1)$ из
примера~\ref{ex:An-invariants}; замена $u=x+iy$, $v=x-iy$ переводит её в $x^2+y^2=w^{n+1}$.
Особенность $A_1$ --- обыкновенная двойная точка (квадратичный конус).

\begin{theorem}[Классификация лог-терминальных особенностей поверхностей]\label{thm:lt-surface}
Пусть $0\in X$ --- нормальная особенность поверхности.
\begin{enumerate}
\item (Каваматa, \cite[4.18]{KM}) $X$ klt (лог-терминальна, см. определение~\ref{def:klt}) тогда и
только тогда, когда $X$ --- факторособенность $\A^2/G$, $G\subset\GL_2(\C)$ конечная малая группа.
\item $X$ канонична тогда и только тогда, когда $X$ --- особенность Дю Валя, тогда и только тогда,
когда $X$ klt и горенштейнова.
\item $X$ терминальна тогда и только тогда, когда $X$ гладка (упражнение~2.1(1) обзора; см.
пример~\ref{ex:terminal-smooth}).
\item Циклические факторособенности $\frac1n(1,q)$ --- это в точности торические особенности
поверхностей; они klt, а горенштейновы ровно при $q=n-1$ (тип $A_{n-1}$).
\end{enumerate}
\end{theorem}

\begin{example}[Дискрепантности $\frac1n(1,1)$]\label{ex:discrep-P11n}
Пусть $X=\frac1n(1,1)$, $\pi\colon Y\to X$ --- раздутие вершины, $E$ --- $(-n)$-кривая,
$K_Y=\pi^*K_X+aE$. Умножим на $E$: $K_Y\cdot E=aE^2=-an$. По формуле присоединения
$K_Y\cdot E=-2-E^2=n-2$, откуда
\[
a=\frac{2-n}n=-1+\frac2n,\qquad A_X(E)=1+a=\frac2n .
\]
При $n=2$: $a=0$ (канонична, $A_1$); при $n=3$: $a=-1/3$; при $n=4$: $a=-1/2$. Все $a>-1$,
т.\,е. особенности klt, но при $n\ge3$ не канонические. Это ответ на упражнения~2.1(3) и
6.3(10b) обзора.
\end{example}

\subsection{$T$-особенности и $\Q$-горенштейновы сглаживания}\label{ss:T-sing}

\begin{definition}\label{def:QGor-smoothing}
\emph{$\Q$-горенштейнова деформация} особенности $0\in X$ --- плоское семейство $\cX\to(0\in T)$
со слоем $X$ над $0$, такое что $K_{\cX/T}$ --- $\Q$-картьеров дивизор (равносильно: некоторое
кратное $mK_\cX$ --- дивизор Картье). \emph{$\Q$-горенштейново сглаживание} --- такая деформация
над кривой с гладкими общими слоями.
\end{definition}

\begin{theorem}[{Коллар--Шеперд-Баррон, \cite[3.10]{KSB}, \cite[\S\,2]{HP}}]\label{thm:KSB}
Факторособенность поверхности допускает $\Q$-горенштейново сглаживание тогда и только тогда, когда
она либо является особенностью Дю Валя, либо имеет вид
\[
\frac1{dn^2}(1,\,dna-1),\qquad d,n,a\ge1,\ \gcd(a,n)=1 .
\]
Особенности второго типа называются \emph{$T$-особенностями} (\emph{особенностями класса $T$}).
\end{theorem}

\begin{example}\label{ex:T-sing}
\begin{enumerate}
\item $d=1$, $n=2$, $a=1$: $\frac14(1,1)$ --- $T$-особенность (конус над рациональной нормальной
кривой степени 4; её $\Q$-горенштейново сглаживание открыл Пинкхэм). $d=1$, $n=3$: $\frac19(1,2)$;
$d=1$, $n=5$: $\frac1{25}(1,4)$. $d=2$, $n=2$, $a=1$: $\frac18(1,3)$.
\item $\frac13(1,1)$ \emph{не} является $T$-особенностью: $dn^2=3$ даёт $d=3$, $n=1$, и тогда
$dna-1=3a-1\equiv2\pmod3$, т.\,е. получается $\frac13(1,2)=A_2$, а не $\frac13(1,1)$. Поэтому
\emph{$\frac13(1,1)$ не допускает $\Q$-горенштейнова сглаживания} --- ключевой факт для описания
K-модулей кубических поверхностей (теорема~\ref{thm:cubic-list}).
\item Явное сглаживание $T$-особенности: $\frac1{dn^2}(1,dna-1)=\{xy=z^{dn}\}/\mu_n$, где
$\zeta\cdot(x,y,z)=(\zeta x,\zeta^{-1}y,\zeta^az)$ (подгруппа $\mu_{dn}\subset\mu_{dn^2}$ действует
весами $(1,-1)$ и даёт $A_{dn-1}=\{xy=z^{dn}\}$; на фактор действует $\mu_{dn^2}/\mu_{dn}\cong\mu_n$).
Семейство $\{xy=z^{dn}+t\}/\mu_n$ --- $\Q$-горенштейново сглаживание: наверху это гиперповерхность
(горенштейнова), а $\mu_n$ действует на форме $dx\wedge dy\wedge dz$ характером, так что
$nK$ тотального пространства фактора --- Картье. Слой Милнора --- фактор слоя Милнора $A_{dn-1}$
по свободному действию $\mu_n$; его число Милнора равно $d-1$. При $d=1$ слой Милнора --- рациональный
гомологический шар: $\Q$-горенштейновы сглаживания $\frac1{n^2}(1,na-1)$ «не меняют» рациональных
гомологий.
\item Почему важно слово «$\Q$-горенштейново». Конус $X=\PP(1,1,3)$ над скрученной кубикой
(в $\PP^4$) является плоским пределом кубических скроллов $S(1,2)\cong\mathbb F_1\subset\PP^4$, т.\,е.
\emph{сглаживается}; но $K_X^2=25/3$, а $K_{\mathbb F_1}^2=8$. В $\Q$-горенштейновом семействе
$K^2$ постоянен (раздел~\ref{sec:families}), так что это сглаживание не $\Q$-горенштейново ---
в полном согласии с (2).
\end{enumerate}
\end{example}

\subsection{Поверхности дель Пеццо, в том числе особые}\label{ss:dP}

Гладкие поверхности дель Пеццо классифицированы в \cite[6.10]{AG}: $\PP^1\times\PP^1$ и
$\Bl_r\PP^2$, $0\le r\le8$, точки общего положения; степень $d=K^2=9-r$ (упражнение~3.4(1,3)
обзора). Нам нужны также особые.

\begin{definition}
\emph{Поверхность дель Пеццо} --- нормальная проективная поверхность $X$ с klt-особенностями и
обильным $-K_X$ (т.\,е. $\Q$-Фано поверхность); \emph{горенштейнова поверхность дель Пеццо} ---
нормальная проективная поверхность с обильным дивизором Картье $-K_X$; \emph{степень}
$d=(-K_X)^2$.
\end{definition}

\begin{theorem}[{Хидака--Ватанабе, \cite{HW}; Демазюр \cite{De}; Фудзита \cite{Fuj90}}]\label{thm:Gor-dP}
Пусть $X$ --- горенштейнова поверхность дель Пеццо степени $d$. Тогда либо все особенности $X$ ---
особенности Дю Валя, либо $X$ --- конус над гладкой эллиптической кривой степени $d$. В первом
случае $1\le d\le9$, минимальное разрешение $\tilde X\to X$ --- гладкая поверхность с
$-K_{\tilde X}$ nef и большим (слабая дель Пеццо), и:
\begin{itemize}
\item при $d\ge3$ дивизор $-K_X$ очень обилен и вкладывает $X$ в $\PP^d$ как поверхность
степени $d$; при $d=3$ --- как кубическую поверхность в $\PP^3$;
\item при $d=2$ система $|-K_X|$ задаёт двулистное накрытие $\PP^2$; при $d=1$ система $|-2K_X|$
задаёт двулистное накрытие квадратичного конуса.
\end{itemize}
Обратно, всякая нормальная кубическая поверхность --- горенштейнова дель Пеццо степени 3.
\end{theorem}

Таким образом, нормальные кубические поверхности --- это: гладкие; с особенностями Дю Валя;
конусы над гладкими плоскими кубиками. Полный список конфигураций особенностей Дю Валя на кубических
поверхностях известен со времён Шлефли и Кэли \cite{BW}:
\[
\begin{array}{l}
A_1,\ 2A_1,\ A_2,\ 3A_1,\ A_1A_2,\ A_3,\ 4A_1,\ 2A_1A_2,\ A_1A_3,\ 2A_2,\ A_4,\ D_4,\\
2A_1A_3,\ A_12A_2,\ A_1A_4,\ A_5,\ D_5,\ 3A_2,\ A_1A_5,\ E_6
\end{array}
\]
--- всего 20 типов; сумма чисел Милнора (индексов при буквах) не превосходит $6=\rk E_6$. Множество
кубических поверхностей с данной конфигурацией особенностей, рассматриваемое с точностью до
проективной эквивалентности, имеет размерность $4-\mu$, где $\mu$ --- суммарное число Милнора
\cite{BW}; в частности, поверхности типов $4A_1$, $3A_2$, $A_1A_5$, $E_6$ единственны.

\begin{example}[Нормальные формы]\label{ex:cubic-normal-forms}
\begin{enumerate}
\item Кубика с особой точкой $p=[0:0:0:1]$: $F=x_3Q(x_0,x_1,x_2)+C(x_0,x_1,x_2)$ ($Q$ ---
квадратичная, $C$ --- кубическая форма). В карте $x_3=1$ это $Q+C=0$, так что $Q$ --- квадратичная
часть в $p$. Точка $p$ --- $A_1$ $\iff$ $\rk Q=3$. Если $\rk Q=2$, скажем $Q=x_0x_1$, то $p$ ---
$A_2$ $\iff$ коэффициент при $x_2^3$ в $C$ отличен от нуля (по лемме о расщеплении
$x_0x_1+C\sim x_0'x_1'+h(x_2)$ с $h(x_2)\equiv C(0,0,x_2)\bmod x_2^4$); иначе $p$ --- $A_k$ с
$k\ge3$ или неизолированная особенность. Если $\rk Q\le1$, то $p$ --- особенность типа $D$, $E$
или хуже (не klt или неизолированная).
\item $3A_2$: $x_0x_1x_2=x_3^3$. Особые точки --- $[1:0:0:0]$, $[0:1:0:0]$, $[0:0:1:0]$; в карте
$x_0=1$: $x_1x_2=x_3^3$ --- особенность $A_2$. Эта поверхность --- фактор $\PP^2/\mu_3$ по действию
$[u:v:w]\mapsto[u:\omega v:\omega^2w]$: инварианты $x_0=u^3$, $x_1=v^3$, $x_2=w^3$, $x_3=uvw$
удовлетворяют $x_0x_1x_2=x_3^3$. Она торическая; группа автоморфизмов содержит двумерный тор
$\{(x_0,x_1,x_2,x_3)\mapsto(ax_0,bx_1,cx_2,x_3): abc=1\}$ и группу $S_3$ перестановок.
\item $4A_1$ (кубика Кэли): $x_0x_1x_2+x_0x_1x_3+x_0x_2x_3+x_1x_2x_3=0$; особые точки --- четыре
координатные.
\item $E_6$: $x_3x_0^2+x_0x_2^2+x_1^3=0$; в карте $x_3=1$: $x_0^2+x_0x_2^2+x_1^3=
(x_0+\tfrac12x_2^2)^2+x_1^3-\tfrac14x_2^4$ --- особенность $y^3+z^4$ типа $E_6$.
\item Конус над плоской кубикой: $F=C(x_0,x_1,x_2)$; вершина --- простая эллиптическая особенность
$\tilde E_6$ (lc, не klt, см. пример~\ref{ex:elliptic-cone}).
\item Неизолированные особенности: $x_0^2x_2+x_1^2x_3=0$ (особа вдоль прямой $x_0=x_1=0$);
приводимые кубики (плоскость + квадрика); ненормальные.
\end{enumerate}
\end{example}

% ============================================================
\section{Семейства, схемы Гильберта, $\Q$-горенштейновы семейства}\label{sec:families}
% ============================================================

\subsection{Плоские семейства и схемы Гильберта}

\emph{Семейство} многообразий над базой $T$ --- морфизм $f\colon\cX\to T$; слой над $t\in T$
обозначается $X_t=f^{-1}(t)$. Правильное условие «непрерывной зависимости слоёв от $t$» ---
\emph{плоскость} $f$ \cite[III.9]{Ha}; для нас достаточно двух фактов: (а) если $T$ --- гладкая
кривая и $\cX$ неприводимо и доминирует $T$, то $f$ плоский \cite[III.9.7]{Ha}; (б) для плоского
проективного семейства $\cX\subset\PP^N\times T$ многочлен Гильберта слоёв
$P_{X_t}(m)=\chi(X_t,\cO_{X_t}(m))$ локально постоянен по $t$ \cite[III.9.9]{Ha}. В частности,
в плоском семействе постоянны размерность, степень и, для гладких слоёв, $K^n$ (у гладких слоёв
$K$ --- Картье и индексы пересечения локально постоянны).

\begin{theorem}[{Гротендик; \cite[гл.~1]{HM}, \cite[I.1]{Kol96}}]\label{thm:Hilb}
Для многочлена $P$ существует проективная схема $\Hilb_P(\PP^N)$ (\emph{схема Гильберта}) с
плоским семейством $\mathcal U\subset\PP^N\times\Hilb_P$, параметризующая замкнутые подсхемы
$\PP^N$ с многочленом Гильберта $P$: морфизмы $T\to\Hilb_P$ взаимно однозначно соответствуют
плоским семействам подсхем $\PP^N\times T$ с многочленом Гильберта $P$. Группа $\PGL_{N+1}$
действует на $\Hilb_P$.
\end{theorem}

\begin{example}
Гиперповерхности степени $d$ в $\PP^N$ образуют $\Hilb_P(\PP^N)=\PP(\Sym^d(\C^{N+1})^*)$: кубические
поверхности в $\PP^3$ --- это $\PP^{19}$. Грассманиан $\Gr(2,4)$ --- схема Гильберта прямых в $\PP^3$.
\end{example}

Главное для нас следствие проективности: \emph{пределы существуют}. Пусть $X\subset\PP^N$ и
$\lambda\colon\Gm\to\PGL_{N+1}$ --- однопараметрическая подгруппа. Отображение
$\Gm\to\Hilb$, $t\mapsto[\lambda(t)X]$, продолжается (по критерию собственности) до морфизма
$\A^1\to\Hilb$; его значение в $0$ --- \emph{плоский предел} $X_0=\lim_{t\to0}\lambda(t)X$.
Конкретно: $\cX=\overline{\{(x,t):x\in\lambda(t)X,\ t\ne0\}}\subset\PP^N\times\A^1$ --- замыкание,
оно плоско над $\A^1$ (факт (а)), и $X_0$ --- его слой над $0$. Слой $X_0$ может быть приводимым,
неприведённым, ненормальным: см. пример~\ref{ex:conic-tc}.

\subsection{$\Q$-горенштейновы семейства}

\begin{definition}\label{def:QGor-family}
Семейство $f\colon\cX\to T$ нормальных многообразий называется \emph{$\Q$-горенштейновым}, если
$f$ плоское и относительный канонический дивизор $K_{\cX/T}$ является $\Q$-картьеровым. На практике
(см.~\cite[\S\,2.5]{Xu}, \cite{Kol23}) требуют также \emph{условия Коллара}: рефлексивные степени
$\omega_{\cX/T}^{[m]}$ коммутируют с заменой базы для всех $m$; для семейств над гладкими кривыми
это автоматически следует из $\Q$-картьеровости.
\end{definition}

Смысл условия: все слои $X_t$ $\Q$-горенштейновы, причём $K_{\cX/T}|_{X_t}=K_{X_t}$, и
\emph{индексы пересечения с $K_{X_t}$ локально постоянны}; в частности, $(-K_{X_t})^n$ постоянен,
и обильность $-K_{X_t}$ --- открытое условие по $t$ (пример~5.2 обзора: обильность --- открытое
условие для семейств дивизоров Картье). Пример~\ref{ex:T-sing}(4) показывает, что без
$\Q$-горенштейновости $K^2$ может скакать.

\begin{example}
\begin{enumerate}
\item Любое семейство гиперповерхностей $\{F_t=0\}\subset\PP^N\times T$ с нормальными слоями
$\Q$-горенштейново: $K_{\cX/T}=\cO_\cX(d-N-1)$ --- Картье (пример~\ref{ex:hyp-Gor}).
Это относится к семействам кубических поверхностей.
\item Семейство $\{xy=z^{dn}+t\}/\mu_n$ из примера~\ref{ex:T-sing}(3) $\Q$-горенштейново.
\item Семейство $\{xy=z^2+tw\}\subset\PP(1,1,1,2)\times\A^1_t$: при $t\ne0$ слой изоморфен $\PP^2$
(проекция $[x:y:z:w]\mapsto[x:y:z]$ --- изоморфизм, так как $w=(xy-z^2)/t$), а при $t=0$ ---
$\{xy=z^2\}\subset\PP(1,1,1,2)$, что есть $\PP(1,1,4)$ (пример~\ref{ex:Markov}). Это
$\Q$-горенштейново вырождение $\PP^2$ в $\PP(1,1,4)$: обе поверхности имеют $K^2=9$.
\end{enumerate}
\end{example}

\subsection{Что такое «пространство модулей»}\label{ss:moduli-def}

Неформально, \emph{задача модулей} --- это правило, сопоставляющее базе $T$ множество (точнее,
группоид) семейств $\cX\to T$ интересующих нас объектов; \emph{грубое пространство модулей} ---
многообразие (схема) $M$, точки которого взаимно однозначно соответствуют классам изоморфизма объектов,
причём каждое семейство над $T$ задаёт морфизм $T\to M$. Если у объектов есть нетривиальные
автоморфизмы, грубое пространство модулей теряет информацию; современный язык --- \emph{алгебраические
стеки}: \emph{стек модулей} $\cM$ помнит семейства вместе с их автоморфизмами, а \emph{хорошее
пространство модулей} (good moduli space, Альпер \cite{Alper}) $\cM\to M$ --- наилучшее приближение
стека алгебраическим пространством: точки $M$ соответствуют \emph{замкнутым} точкам стека
(объектам, не вырождающимся изотривиально в другие), а семейства над $T$ дают морфизмы $T\to M$.
Образцом служит ГИТ (раздел~\ref{sec:git}): для действия редуктивной группы $G$ на $X$ стек
$[X^{\ssl}/G]$ имеет хорошее пространство модулей $X^{\ssl}/\!/G$, точки которого --- замкнутые
орбиты (полистабильные точки).

От пространства модулей многообразий хочется:
\begin{itemize}
\item \emph{ограниченности}: все объекты вкладываются в одно $\PP^N$ (тогда их можно
параметризовать схемой Гильберта);
\item \emph{отделимости}: семейство над проколотой кривой имеет не более одного предела;
\item \emph{собственности}: предел существует (после конечной замены базы);
\item разумных групп автоморфизмов (конечных или хотя бы редуктивных).
\end{itemize}
В разделе~\ref{sec:moduli} мы увидим, что для многообразий Фано «в целом» ни одно из этих
свойств не выполняется, а для K-полустабильных --- все.

% ============================================================
\section{Геометрическая теория инвариантов}\label{sec:git}
% ============================================================

Геометрическая теория инвариантов (ГИТ) Мамфорда \cite{MFK} --- классический способ строить
пространства модулей как факторы по действию группы. В теории K-стабильности она нужна дважды:
как источник интуиции (тест-конфигурации --- аналог однопараметрических подгрупп, инвариант
Футаки --- аналог веса Гильберта--Мамфорда, полистабильность --- замкнутость орбиты) и как
конкретный инструмент: для кубических поверхностей K-стабильность \emph{совпадает} с
ГИТ-стабильностью (теорема~\ref{thm:GIT=K}), и пространство K-модулей --- это ГИТ-фактор.
Изложение следует \cite{MFK}, \cite{Dolg-IT}, \cite{Mukai}, \cite{PV}.

\subsection{Редуктивные группы и аффинные факторы}

\begin{definition}
Линейная алгебраическая группа $G$ \emph{редуктивна}, если она не содержит нетривиальных связных
нормальных унипотентных подгрупп. Над $\C$ это равносильно тому, что всякое конечномерное
рациональное представление $G$ вполне приводимо, и тому, что $G$ --- комплексификация компактной
группы Ли. Примеры: торы $(\C^*)^n$, $\GL_n$, $\SL_n$, $\PGL_n$, $\mathrm{SO}_n$, $\mathrm{Sp}_{2n}$,
конечные группы и их полупрямые произведения с торами. Нередуктивны: $\mathbb G_a=(\C,+)$,
группа верхнетреугольных матриц, стабилизатор точки в $\PGL_n$ (пример~\ref{ex:Aut-F1}).
\end{definition}

\begin{theorem}[Гильберт, Нагата]\label{thm:Hilbert-finiteness}
Пусть редуктивная группа $G$ действует на аффинном многообразии $X=\Spec A$. Тогда алгебра
инвариантов $A^G$ конечно порождена, и морфизм $\pi\colon X\to X/\!/G:=\Spec A^G$ обладает
свойствами:
\begin{enumerate}
\item $\pi$ сюръективен и $G$-инвариантен; $\pi(x)=\pi(y)$ тогда и только тогда, когда
$\overline{Gx}\cap\overline{Gy}\ne\varnothing$;
\item каждый слой $\pi$ содержит ровно одну замкнутую орбиту; точки $X/\!/G$ --- это замкнутые орбиты;
\item $\pi$ отображает замкнутые $G$-инвариантные подмножества в замкнутые.
\end{enumerate}
\end{theorem}

Ключевой инструмент доказательства --- \emph{оператор Рейнольдса} $A\to A^G$ (усреднение по
компактной форме), существующий благодаря редуктивности; для нередуктивных групп алгебра
инвариантов может быть бесконечно порождённой (контрпример Нагаты).

\begin{example}\label{ex:Cstar-C2}
$G=\C^*$ действует на $\A^2$: $t\cdot(x,y)=(tx,t^{-1}y)$. Инварианты: $\C[x,y]^G=\C[xy]$,
$X/\!/G=\A^1$. Орбиты: гиперболы $\{xy=c\}$, $c\ne0$ (замкнутые); две проколотые оси
$\{x\ne0,y=0\}$, $\{x=0,y\ne0\}$ (незамкнутые) и начало координат (замкнутая). Три последние
орбиты склеиваются в одну точку $0\in\A^1$: фактор «не видит» разницы между осью и её замыканием.
Это --- модель всего, что будет дальше: $X/\!/G$ параметризует \emph{замкнутые} орбиты.
\end{example}

\subsection{Проективные факторы и (полу)стабильность}

Пусть редуктивная $G$ действует на проективном многообразии $X\subset\PP(V)$ линейно (через
представление $G\to\GL(V)$; это называется \emph{линеаризацией} действия). Тогда $G$ действует
на однородном координатном кольце $R=\bigoplus_mH^0(X,\cO(m))$, и $R^G$ конечно порождена.

\begin{definition}[Мамфорд]\label{def:GIT-stability}
Точка $x\in X$ с представителем $\hat x\in V\setminus0$ называется
\begin{itemize}
\item \emph{полустабильной}, если $0\notin\overline{G\hat x}$ (равносильно: существует
$G$-инвариантный однородный многочлен $f$ положительной степени с $f(\hat x)\ne0$);
\item \emph{полистабильной}, если орбита $G\hat x\subset V$ замкнута (и $\hat x\ne0$);
\item \emph{стабильной}, если $G\hat x$ замкнута и стабилизатор $\Stab_G(\hat x)$ конечен;
\item \emph{нестабильной} (\emph{неустойчивой}), если не полустабильна.
\end{itemize}
Множества $X^{\ssl}\supset X^{\mathrm{ps}}\supset X^{s}$ обозначают полустабильные, полистабильные и
стабильные точки.
\end{definition}

\begin{theorem}[{Мамфорд, \cite[гл.~1]{MFK}}]\label{thm:GIT}
\begin{enumerate}
\item $X^{\ssl}$ и $X^s$ --- открытые $G$-инвариантные подмножества $X$; $X^{\ssl}=X\setminus V(R^G_+)$.
\item Существует проективное многообразие $X/\!/G:=\Proj R^G$ и $G$-инвариантный сюръективный
морфизм $\pi\colon X^{\ssl}\to X/\!/G$ со свойствами теоремы~\ref{thm:Hilbert-finiteness} (локально
по базе он является аффинным фактором): $\pi(x)=\pi(y)\iff\overline{Gx}\cap\overline{Gy}\cap X^{\ssl}\ne\varnothing$;
каждый слой содержит единственную замкнутую (в $X^{\ssl}$) орбиту, и это в точности орбиты
полистабильных точек.
\item $\pi(X^s)$ открыто в $X/\!/G$, и $X^s\to\pi(X^s)$ --- \emph{геометрический фактор}: его
слои --- в точности орбиты.
\end{enumerate}
\end{theorem}

Итак, $X/\!/G$ --- проективное (в частности, отделимое и полное) многообразие, точки которого
--- классы \emph{полистабильных} орбит; полустабильные, но не полистабильные точки отождествляются
с полистабильным пределом своей орбиты. Два полустабильных объекта называются
\emph{$S$-эквивалентными}, если замыкания их орбит пересекаются (в $X^{\ssl}$): $X/\!/G$
параметризует классы $S$-эквивалентности.

\begin{remark}[Почему это «хорошее» пространство модулей]
Свойства из \S\ref{ss:moduli-def}: $X/\!/G$ проективно (собственность и отделимость), стабилизаторы
полистабильных точек редуктивны (теорема Мацусимы в ГИТ: стабилизатор точки с замкнутой орбитой
редуктивен), стабилизаторы стабильных --- конечны. Важно: \emph{предел семейства полистабильных
объектов единствен}. Точнее, пусть $x(t)$, $t\in C\setminus0$, --- семейство полустабильных точек над
проколотой кривой, и $x_0=\lim_{t\to0}x(t)$, $x_0'=\lim_{t\to0}g(t)x(t)$ --- два полистабильных
предела (здесь $g(t)$ --- семейство элементов $G$, определённое при $t\ne0$). Тогда $x_0$ и $x_0'$
лежат в одной орбите. Действительно, $\pi(x(t))=\pi(g(t)x(t))$ при $t\ne0$, по отделимости
$X/\!/G$ пределы совпадают: $\pi(x_0)=\pi(x_0')$, а в слое $\pi$ есть лишь одна замкнутая орбита.
Именно это свойство используется в доказательстве теоремы~\ref{thm:GIT=K}.
\end{remark}

\subsection{Критерий Гильберта--Мамфорда}

Проверять замкнутость орбит и существование инвариантов напрямую трудно. Критерий
Гильберта--Мамфорда сводит всё к однопараметрическим подгруппам.

\begin{definition}
\emph{Однопараметрическая подгруппа} (1-ПП) группы $G$ --- гомоморфизм алгебраических групп
$\lambda\colon\Gm\to G$. Для представления $G\to\GL(V)$ всякая 1-ПП диагонализуема: $V=\bigoplus_rV_r$,
где $\lambda(t)$ действует на $V_r$ умножением на $t^r$ (\emph{веса}). Для $\hat x=\sum\hat x_r$
положим
\[
\mu(\hat x,\lambda)=\min\{r:\hat x_r\ne0\}.
\]
Тогда $\lambda(t)\hat x=\sum t^r\hat x_r$, и предел $\lim_{t\to0}\lambda(t)\hat x$ существует
$\iff\mu\ge0$; он равен нулю $\iff\mu>0$; он равен $\hat x_{\mu}$ (компоненте веса $\mu$), если $\mu=0$.
\end{definition}

\begin{theorem}[{Гильберт--Мамфорд, \cite[2.1]{MFK}}]\label{thm:HM}
Пусть редуктивная $G$ действует линейно на $X\subset\PP(V)$, $x\in X$, $\hat x\in V$ --- представитель.
\begin{enumerate}
\item $x$ полустабильна $\iff$ $\mu(\hat x,\lambda)\le0$ для всех 1-ПП $\lambda$ группы $G$;
\item $x$ стабильна $\iff$ $\mu(\hat x,\lambda)<0$ для всех нетривиальных 1-ПП $\lambda$;
\item $x$ нестабильна $\iff$ существует $\lambda$ с $\lim_{t\to0}\lambda(t)\hat x=0$;
\item (\emph{полистабильность}) полустабильная $x$ полистабильна $\iff$ для всякой $\lambda$ с
$\mu(\hat x,\lambda)=0$ предел $\lim_{t\to0}\lambda(t)\hat x$ лежит в орбите $G\hat x$.
\end{enumerate}
\end{theorem}

Содержательная часть --- «$\Leftarrow$» в (1): если $0\in\overline{G\hat x}$, то нуля можно
достичь уже вдоль некоторой 1-ПП (теорема Гильберта--Мамфорда в собственном смысле; её
доказательство использует разложение Картана $G=KTK$ и критерий собственности). Пункт (4) ---
результат Кемпфа \cite{Kempf} (см. также \cite{BMRT}).

\emph{Как пользоваться.} Все 1-ПП группы $\SL_{n}$ сопряжены диагональным
$\lambda(t)=\mathrm{diag}(t^{r_0},\dots,t^{r_{n-1}})$, $\sum r_i=0$, $r_0\ge r_1\ge\dots\ge r_{n-1}$.
Если $G=\SL_n$ действует на формах $F(x_0,\dots,x_{n-1})$ степени $d$ заменой координат, то
удобно считать, что $\lambda(t)$ действует подстановкой $x_i\mapsto t^{r_i}x_i$: моном $x^a$
получает вес $\langle a,r\rangle=\sum a_ir_i$. Тогда:
\begin{itemize}
\item $F$ нестабильна $\iff$ в \emph{некоторой} системе координат найдётся вектор весов $r$ (как
выше), при котором все мономы $F$ имеют положительный вес;
\item $F$ не стабильна $\iff$ в некоторой системе координат найдётся $r\ne0$, при котором все мономы
$F$ имеют неотрицательный вес.
\end{itemize}
(Переход $\lambda\mapsto\lambda^{-1}$ меняет знак весов, поэтому выбор знака в определении $\mu$
несуществен для этих формулировок.)

\subsection{Пример: бинарные формы}

\begin{example}[Наборы точек на $\PP^1$]\label{ex:binary-forms}
Пусть $\SL_2$ действует на бинарных формах $F(x,y)$ степени $d$, т.\,е. на наборах $d$ точек
$\PP^1$ с кратностями. Диагональная 1-ПП $\lambda=\mathrm{diag}(t^{r},t^{-r})$, $r>0$, даёт моному
$x^iy^{d-i}$ вес $r(i-(d-i))=r(2i-d)$. Все мономы $F$ имеют положительный вес $\iff$ в $F$ входят только мономы $x^iy^{d-i}$ с
$i>d/2$ $\iff$ $F$ делится на $x^{\lfloor d/2\rfloor+1}$ $\iff$ точка $[0:1]=\{x=0\}$ имеет
кратность $>d/2$. Сопряжением получаем:
\[
F\ \text{нестабильна}\iff\text{некоторая точка имеет кратность}>d/2;\qquad
F\ \text{не стабильна}\iff\text{некоторая точка имеет кратность}\ge d/2 .
\]
При нечётном $d$ полустабильность совпадает со стабильностью. При $d=4$: четвёрка точек
полустабильна, если никакие три не совпадают; стабильна, если все различны; строго полустабильные
полистабильные точки --- пары двойных точек $x^2y^2$ (орбита замкнута: стабилизатор --- тор).
Фактор $\PP^4/\!/\SL_2\cong\PP^1$ --- это $j$-прямая (четвёрки точек с точностью до дробно-линейных
преобразований классифицируются $j$-инвариантом), и в единственную «граничную» точку склеиваются
все четвёрки с одной двойной точкой. Это наглядная модель того, что произойдёт с кубическими
поверхностями.
\end{example}

\begin{example}[Плоские кубики]\label{ex:plane-cubics}
$\SL_3$ действует на $\PP^9=\PP(\Sym^3(\C^3)^*)$ --- пространстве плоских кубик. Классический ответ:
кубика $C$ стабильна $\iff$ гладка; полустабильна $\iff$ имеет лишь узлы (обыкновенные двойные
точки); единственная полистабильная строго полустабильная кубика --- треугольник $xyz=0$. Фактор
$\PP^9/\!/\SL_3\cong\PP^1$ (снова $j$-инвариант), и все узловые кубики ($j=\infty$)
$S$-эквивалентны треугольнику. Проверим «нестабильную» половину по критерию Гильберта--Мамфорда.
Пусть $C$ особа в $p=[0:0:1]$: $F=zQ(x,y)+C_3(x,y)$, где $Q$ --- квадратичная часть в $p$.
\begin{itemize}
\item \emph{Касп или хуже} ($\rk Q\le1$, $Q=x^2$): $F=zx^2+C_3(x,y)$. Веса $r=(2,1,-3)$ (сумма
нулевая, упорядочены по убыванию) дают: $zx^2\mapsto4-3=1$, $x^3\mapsto6$, $x^2y\mapsto5$,
$xy^2\mapsto4$, $y^3\mapsto3$. Все веса положительны: $C$ нестабильна.
\item \emph{Узел} ($Q=xy$): $F=zxy+C_3(x,y)$. Веса $r=(1,1,-2)$ дают $zxy\mapsto0$, а мономам
$C_3$ --- вес $3>0$. Все веса неотрицательны (не стабильна), и
$\lim_{t\to0}\lambda(t)F=zxy$ --- треугольник. Нестабильной узловая кубика не является, но
доказательство её полустабильности требует перебора всех систем координат (он делается так же, как
ниже для кубических поверхностей), и мы его опускаем.
\end{itemize}
Урок примера: при работе с критерием Гильберта--Мамфорда нужно перебирать \emph{все} упорядочения
весов и все системы координат; систематический способ --- таблица максимальных «нестабильных
наборов мономов», которую мы сейчас составим для кубических поверхностей.
\end{example}

\subsection{Теорема Гильберта о кубических поверхностях}\label{ss:Hilbert-cubic}

Пусть $G=\SL_4$ действует на $\PP^{19}=\PP(\Sym^3(\C^4)^*)$ --- пространстве кубических форм
$F(x_0,x_1,x_2,x_3)$. Мономы степени 3 --- это векторы $a=(a_0,a_1,a_2,a_3)\in\Z_{\ge0}^4$,
$|a|=3$ (их 20). Для вектора весов $r=(r_0,r_1,r_2,r_3)$, $\sum r_i=0$, $r_0\ge r_1\ge r_2\ge r_3$,
$r\ne0$, положим
\[
M_r=\{a:\langle a,r\rangle\ge0\},\qquad M_r^+=\{a:\langle a,r\rangle>0\}.
\]
По критерию Гильберта--Мамфорда $F$ не стабильна $\iff$ $\Supp F\subset M_r$ для некоторого $r$ в
некоторой системе координат, и нестабильна $\iff$ $\Supp F\subset M^+_r$. Достаточно знать
\emph{максимальные} по включению множества $M_r$ и $M_r^+$.

\begin{lemma}\label{lem:rays}
Рассмотрим в трёхмерном пространстве $\{r:\sum r_i=0\}$ конус $C=\{r_0\ge r_1\ge r_2\ge r_3\}$ и
расположение плоскостей $\langle a,r\rangle=0$, $|a|=3$. Максимальные множества $M_r$ достигаются
на \emph{лучах} этого расположения (одномерных гранях разбиения $C$ плоскостями), а максимальные
$M_r^+$ --- на открытых камерах (трёхмерных гранях).
\end{lemma}
\begin{proof}
Если $r$ лежит в относительной внутренности грани $\sigma$, а $r'$ --- в её замыкании, то знак
$\langle a,\cdot\rangle$ постоянен на $\sigma$, поэтому $\langle a,r\rangle\ge0\Rightarrow
\langle a,r'\rangle\ge0$, т.\,е. $M_r\subset M_{r'}$; значит, максимальные $M_r$ сидят на гранях
наименьшей размерности --- лучах. Наоборот, $\langle a,r'\rangle>0\Rightarrow\langle a,r\rangle>0$ для
$r$, близких к $r'$, т.\,е. $M^+_{r'}\subset M^+_r$ для $r$ в соседней камере; значит, максимальные
$M^+$ достигаются в камерах, где $M^+_r=M_r$.
\end{proof}

Расположение имеет 13 лучей:
\[
\begin{array}{llll}
(1,0,0,-1),&(1,1,-1,-1),&(1,1,0,-2),&(1,1,1,-3),\\
(2,0,-1,-1),&(2,1,-1,-2),&(2,1,1,-4),&(2,2,-1,-3),\\
(3,-1,-1,-1),&(3,1,-2,-2),&(3,2,-1,-4),&(4,-1,-1,-2),\\
(4,1,-2,-3).&&&
\end{array}
\]
Перебор (вручную, как у Гильберта и Мамфорда \cite[гл.~4, \S\,2]{MFK}, или короткой программой)
даёт следующий результат.

\begin{proposition}[Таблица Гильберта]\label{prop:Hilbert-table}
Максимальных множеств $M_r$ ровно три; $F$ не стабильна тогда и только тогда, когда в подходящих
координатах $F$ имеет один из видов
\[
\begin{array}{lll}
(N_1)\ r=(1,0,0,-1): & F=x_0\,Q(x_0,x_1,x_2,x_3)+C(x_1,x_2), & Q\ \text{без члена } x_3^2;\\
(N_2)\ r=(1,1,0,-2): & F=x_3\,Q(x_0,x_1)+C(x_0,x_1,x_2); &\\
(N_3)\ r=(2,0,-1,-1): & F=x_0\,Q(x_0,x_1,x_2,x_3)+c\,x_1^3. &
\end{array}
\]
Максимальных множеств $M_r^+$ тоже три; $F$ нестабильна тогда и только тогда, когда в подходящих
координатах
\[
\begin{array}{lll}
(U_1)\ r=(4,4,-1,-7): & F=x_3\,Q(x_0,x_1)+C(x_0,x_1,x_2), & C\ \text{без члена } x_2^3;\\
(U_2)\ r=(5,3,1,-9): & F=c\,x_3x_0^2+C(x_0,x_1,x_2); &\\
(U_3)\ r=(9,-2,-3,-4): & F=x_0\,Q(x_0,x_1,x_2,x_3). &
\end{array}
\]
Здесь $Q$ --- квадратичные, $C$ --- кубические формы от указанных переменных, $c\in\C$.
\end{proposition}

\begin{proof}[Проверка]
Для $r=(1,1,0,-2)$: вес монома $x^a$ равен $a_0+a_1-2a_3\ge0$. Мономы с $a_3=1$: нужно
$a_0+a_1\ge2$, т.\,е. $x_3x_0^2,x_3x_0x_1,x_3x_1^2$; с $a_3\ge2$ вес отрицателен; с $a_3=0$ все 10
мономов от $x_0,x_1,x_2$ подходят. Это $(N_2)$. Для $r=(4,4,-1,-7)$: $a_3=1$ требует
$4(a_0+a_1)-a_2\ge8$, т.\,е. $a_2=0$, $a_0+a_1=2$; $a_3=0$: $4(a_0+a_1)-a_2>0$ исключает лишь
$x_2^3$. Это $(U_1)$. Остальные строки проверяются так же; максимальность и полнота списка --- перебор
13 лучей и камер (мы провели его на компьютере).
\end{proof}

\begin{theorem}[Гильберт, 1893; \cite[гл.~4, \S\,2]{MFK}]\label{thm:Hilbert}
Кубическая поверхность $X=\{F=0\}\subset\PP^3$
\begin{enumerate}
\item стабильна тогда и только тогда, когда все её особые точки --- обыкновенные двойные точки
($A_1$); в частности, все гладкие кубики стабильны;
\item полустабильна тогда и только тогда, когда все её особые точки имеют тип $A_1$ или $A_2$;
\item среди полустабильных, но не стабильных поверхностей ровно одна, с точностью до проективной
эквивалентности, полистабильна: поверхность $x_0x_1x_2=x_3^3$ с тремя особенностями $A_2$. Всякая
кубика с особенностью $A_2$ (и не худшими особенностями) вырождается в неё вдоль
однопараметрической подгруппы.
\end{enumerate}
\end{theorem}
\begin{proof}
Пользуемся нормальными формами из примера~\ref{ex:cubic-normal-forms}(1): если $p=[0:0:0:1]$ ---
особая точка, то $F=x_3Q(x_0,x_1,x_2)+C(x_0,x_1,x_2)$, $Q$ --- квадратичная часть в $p$.

\emph{(1), «$\Leftarrow$».} Пусть $F$ не стабильна. В случае $(N_2)$ точка $p$ особа ($F$ не
содержит $x_3^3$ и $x_3^2x_i$), и квадратичная часть $Q(x_0,x_1)$ имеет ранг $\le2$, так что $p$
--- не $A_1$. В случае $(N_1)$ точка $p$ особа, и квадратичная часть равна
$x_0\ell(x_0,x_1,x_2)$ (слагаемые $F$ с $x_3$ --- это $x_0x_3\ell$), ранг $\le2$. В случае $(N_3)$
рассмотрим прямую $m=\{x_0=x_1=0\}$: на ней $F=0$, и $\nabla F|_m=Q\cdot\nabla x_0$, так что $F$
особа в точках $m$, где $Q=0$ (их две, или $Q|_m\equiv0$ и $F$ особа вдоль всей $m$). В такой
точке $q$ квадратичная часть $F$ равна $x_0\cdot d_qQ$ --- ранга $\le2$. Итак, не стабильная
кубика имеет особую точку с квадратичной частью ранга $\le2$, т.\,е. не $A_1$ (по лемме Морса
особая точка с невырожденной квадратичной частью --- изолированная особенность типа $A_1$).

\emph{(1), «$\Rightarrow$».} Если у $F$ есть особая точка, не являющаяся $A_1$, то в нормальной
форме $\rk Q\le2$, и, выбрав координаты так, что $Q=Q(x_0,x_1)$, получаем $\Supp F\subset(N_2)$.

\emph{(2), «$\Leftarrow$».} Пусть $F$ нестабильна. В случае $(U_1)$ точка $p$ особа с
квадратичной частью $Q(x_0,x_1)$: если $\rk Q=2$, $Q=x_0x_1$, то, поскольку $x_2^3\notin C$, точка
$p$ --- не $A_2$ (пример~\ref{ex:cubic-normal-forms}(1)), а если $\rk Q\le1$, то и подавно. В
случае $(U_2)$ квадратичная часть в $p$ равна $cx_0^2$, ранга $\le1$: не $A_1$ и не $A_2$. В случае
$(U_3)$ поверхность приводима и особа вдоль коники $\{x_0=Q=0\}$. Во всех случаях есть особенность
хуже $A_2$.

\emph{(2), «$\Rightarrow$».} Пусть $p$ --- особая точка не типа $A_1,A_2$. Если $\rk Q=2$, $Q=x_0x_1$,
то $x_2^3\notin C$ и $\Supp F\subset(U_1)$. Если $\rk Q\le1$, $Q=cx_0^2$, то $\Supp F\subset(U_2)$.
Неизолированные особые точки тоже имеют $\rk Q\le2$ и не являются $A_2$, так что попадают в эти
случаи.

\emph{(3).} Пусть $F$ полустабильна и имеет точку $A_2$: $F=x_3x_0x_1+cx_2^3+C'(x_0,x_1,x_2)$,
$c\ne0$, где все мономы $C'$ содержат $x_0$ или $x_1$. Возьмём $r=(1,1,0,-2)$: веса $x_3x_0x_1$ и
$x_2^3$ равны нулю, а веса мономов $C'$ положительны. Поэтому
$\lim_{t\to0}\lambda(t)F=x_0x_1x_3+cx_2^3$ --- поверхность $3A_2$ (после замены $x_2\mapsto c^{-1/3}x_2$).
Она полустабильна по (2), и по теореме~\ref{thm:GIT}(2) в замыкании орбиты любой полустабильной
кубики с точкой $A_2$ лежит единственная замкнутая орбита; остаётся убедиться, что орбита
$F_0=x_0x_1x_2-x_3^3$ замкнута в $X^{\ssl}$. Пусть $F_1\in\overline{GF_0}\cap X^{\ssl}$. Суммарное число
Милнора полунепрерывно сверху при специализации, поэтому у $F_1$ оно $\ge6$; но $F_1$ полустабильна,
т.\,е. имеет лишь особенности $A_1,A_2$, и по списку \S\ref{ss:dP} единственная такая конфигурация
с $\mu=6$ --- это $3A_2$; все поверхности $3A_2$ проективно эквивалентны \cite{BW}. Значит,
$F_1\in GF_0$.
\end{proof}

\begin{corollary}\label{cor:GIT-cubic}
ГИТ-фактор $\cM^{\mathrm{GIT}}=\PP^{19}/\!/\SL_4$ --- проективное четырёхмерное многообразие, точки
которого --- классы проективной эквивалентности кубических поверхностей с особенностями не хуже
$A_1$, плюс одна точка, отвечающая поверхности $3A_2$ (в неё склеены все кубики с особенностями
$A_2$). Классически (Клебш, Салмон) кольцо инвариантов кубических форм от четырёх переменных
порождено инвариантами степеней $8,16,24,32,40$ и $100$ с одним соотношением, выражающим квадрат
инварианта степени 100 через остальные, откуда
\[
\cM^{\mathrm{GIT}}\cong\PP(1,2,3,4,5)
\]
\cite[\S\,10]{Dolg-IT}. Открытое подмножество гладких кубик --- дополнение к дискриминантной
гиперповерхности; поверхности с $k$ узлами ($k\le4$) образуют страты коразмерности $k$, а единственная
поверхность Кэли $4A_1$ --- точку.
\end{corollary}

\begin{remark}
Размерность $4$ согласуется со счётом параметров: кубики --- $\PP^{19}$, $\dim\PGL_4=15$, и
общая кубика имеет конечный стабилизатор; с другой стороны, $\Bl_6\PP^2$ зависит от
$6\cdot2-8=4$ параметров \cite[замечание после 8.13]{AG}.
\end{remark}

% ============================================================
\section{Дифференциально-геометрические пререквизиты: метрики Кэлера--Эйнштейна}\label{sec:dg}
% ============================================================

K-стабильность родилась из задачи о метриках Кэлера--Эйнштейна на многообразиях Фано. В обзоре
\cite{DeV} эта тема лишь упоминается (\S1), но без неё непонятно, откуда берутся определения
раздела~\ref{sec:tc}: почему «вес действия $\Gm$ на сечениях», почему именно такая комбинация
коэффициентов, почему полистабильность связана с редуктивностью. Здесь мы даём минимальный
дифференциально-геометрический фон, следуя \cite{Sz}, \cite{Tian-book}, \cite[гл.~0, 1]{GH},
\cite[\S\,6.2]{AG}.

\subsection{Кэлеровы метрики}

Пусть $X$ --- компактное комплексное многообразие размерности $n$ с локальными координатами
$z_1,\dots,z_n$. \emph{Эрмитова метрика} --- положительно определённая эрмитова форма
$h=\sum h_{j\bar k}\,dz_j\otimes d\bar z_k$ на $T^{1,0}X$; ей отвечает вещественная $(1,1)$-форма
\[
\omega=\frac i2\sum_{j,k}h_{j\bar k}\,dz_j\wedge d\bar z_k .
\]
Метрика \emph{кэлерова}, если $d\omega=0$; тогда $\omega$ называется \emph{кэлеровой формой}, а её
класс $[\omega]\in H^2(X,\R)$ --- \emph{кэлеровым классом}. Локально кэлерова форма имеет
\emph{потенциал}: $\omega=i\partial\bar\partial\varphi$ с вещественной функцией $\varphi$, причём
$(\partial^2\varphi/\partial z_j\partial\bar z_k)$ положительно определена (строгая
плюрисубгармоничность). Форма объёма --- $\omega^n/n!$, и $\int_X\omega^n=[\omega]^n$ --- топологическое
число («самопересечение класса»).

\begin{example}
\begin{enumerate}
\item $\C^n$ с $\omega=\frac i2\sum dz_j\wedge d\bar z_j=i\partial\bar\partial\frac{|z|^2}2$.
\item Метрика Фубини--Штуди $\omega_{FS}=\frac i{2\pi}\partial\bar\partial\log(1+|z|^2)$ на $\PP^n$
\cite[пример~6.5]{AG}; $[\omega_{FS}]=c_1(\cO(1))$, $\int_{\PP^n}\omega_{FS}^n=1$.
\item Ограничение кэлеровой формы на комплексное подмногообразие кэлерово; поэтому всякое проективное
многообразие кэлерово, причём кэлерова форма $\varphi^*\omega_{FS}$ лежит в классе $c_1(L)$
обильного расслоения $L=\varphi^*\cO(1)$.
\end{enumerate}
\end{example}

\begin{lemma}[$\partial\bar\partial$-лемма]\label{lem:ddbar}
Если $\omega,\omega'$ --- замкнутые вещественные $(1,1)$-формы на компактном кэлеровом многообразии и
$[\omega]=[\omega']$ в $H^2(X,\R)$, то $\omega'=\omega+i\partial\bar\partial\varphi$ для вещественной
функции $\varphi$, единственной с точностью до константы.
\end{lemma}

Поэтому все кэлеровы метрики в данном классе $[\omega]$ --- это $\omega_\varphi=\omega+i\partial\bar\partial\varphi$
для $\varphi$ из открытого множества
$\mathcal H=\{\varphi\in C^\infty(X,\R):\omega+i\partial\bar\partial\varphi>0\}$ (\emph{пространство
кэлеровых потенциалов}). Поиск «наилучшей» метрики в классе --- это задача вариационного исчисления
на бесконечномерном пространстве $\mathcal H$.

\subsection{Кривизна линейных расслоений и форма Риччи}

Напомним \cite[\S\,6.2]{AG}: эрмитова метрика $h$ на голоморфном линейном расслоении $L$ имеет
кривизну $\Theta_h=-\partial\bar\partial\log|e|_h^2$ ($e$ --- локальное голоморфное сечение без нулей),
и $\frac i{2\pi}\Theta_h$ --- замкнутая вещественная $(1,1)$-форма, представляющая $c_1(L)$;
$L$ положительно (= обильно, по Кодаире), если для некоторой $h$ эта форма кэлерова.

Кэлерова метрика $\omega$ задаёт метрику на $T_X$, а значит, и на $\det T_X=-K_X$:
$|\partial_1\wedge\dots\wedge\partial_n|^2=\det(h_{j\bar k})$. Её кривизна (умноженная на $i$) ---
\emph{форма Риччи}
\[
\Ric(\omega)=-i\,\partial\bar\partial\log\det(h_{j\bar k}),\qquad [\Ric(\omega)]=2\pi c_1(X),\quad
c_1(X):=c_1(T_X)=c_1(-K_X).
\]
Как и у римановой кривизны Риччи, это симметричный тензор, и $\Ric(\omega)$ --- в точности
$(1,1)$-форма, отвечающая римановой кривизне Риччи метрики $g=\mathrm{Re}\,h$. Важное свойство:
\[
\Ric(\omega+i\partial\bar\partial\varphi)=\Ric(\omega)-i\partial\bar\partial\log\frac{(\omega+i\partial\bar\partial\varphi)^n}{\omega^n}.
\]

\subsection{Уравнение Кэлера--Эйнштейна}

\begin{definition}
Кэлерова метрика $\omega$ называется \emph{метрикой Кэлера--Эйнштейна} (КЭ), если
$\Ric(\omega)=\lambda\omega$ для некоторой константы $\lambda\in\R$.
\end{definition}

Переходя к классам, $2\pi c_1(X)=\lambda[\omega]$. Так как $[\omega]$ --- кэлеров (положительный)
класс, необходимое условие: $c_1(X)$ отрицателен ($\lambda<0$; $K_X$ обилен --- многообразия общего
типа с обильным каноническим классом), нулевой ($\lambda=0$; $K_X$ численно тривиален ---
многообразия Калаби--Яу) или положителен ($\lambda>0$; $-K_X$ обилен --- многообразия Фано). Умножая
метрику на константу, можно считать $\lambda\in\{-1,0,1\}$; в случае Фано нормируем
$[\omega]=2\pi c_1(X)$, т.\,е. $\omega\in2\pi c_1(-K_X)$.

\begin{theorem}[Обен, Яу (1976--1978)]
Если $K_X$ обилен, то на $X$ существует единственная КЭ-метрика с $\lambda=-1$. Если $c_1(X)=0$, то в
каждом кэлеровом классе существует единственная КЭ-метрика (риччи-плоская; гипотеза Калаби).
\end{theorem}

\emph{Идея: сведение к уравнению Монжа--Ампера.} Пусть $X$ --- Фано, $\omega\in2\pi c_1(X)$. По
$\partial\bar\partial$-лемме $\Ric(\omega)-\omega=i\partial\bar\partial h$ для функции $h$
(\emph{потенциал Риччи}). Ищем КЭ-метрику в виде $\omega_\varphi=\omega+i\partial\bar\partial\varphi$.
Условие $\Ric(\omega_\varphi)=\omega_\varphi$ переписывается как
$i\partial\bar\partial\big(h-\log\frac{\omega_\varphi^n}{\omega^n}-\varphi\big)=0$, т.\,е. (с точностью до
константы, которую можно вобрать в $\varphi$)
\[
(\omega+i\partial\bar\partial\varphi)^n=e^{\,h-\varphi}\,\omega^n
\qquad\text{(комплексное уравнение Монжа--Ампера).}
\]
В случае $\lambda=-1$ получается то же уравнение с $e^{h+\varphi}$ в правой части. Разница в
знаке при $\varphi$ решает всё: при $e^{h+\varphi}$ правая часть монотонна по $\varphi$, работает
принцип максимума, и метод непрерывности Обена--Яу даёт решение всегда; при $e^{h-\varphi}$
монотонности нет, и априорные оценки могут не выполняться --- именно здесь возникают препятствия.

\subsection{Препятствия: теоремы Мацусимы и Футаки}

\begin{theorem}[Мацусима, 1957]\label{thm:Matsushima}
Если компактное кэлерово многообразие $X$ с $c_1(X)>0$ допускает КЭ-метрику $\omega$, то связная
группа автоморфизмов $\Aut^0(X)$ редуктивна: она является комплексификацией компактной группы
изометрий $\mathrm{Isom}^0(X,\omega)$, а алгебра Ли голоморфных векторных полей $H^0(X,T_X)$ ---
комплексификацией алгебры киллинговых полей.
\end{theorem}

\begin{example}[Нередуктивная группа автоморфизмов]\label{ex:Aut-F1}
Пусть $X=\mathbb F_1=\Bl_p\PP^2$. Всякий автоморфизм $X$ переводит единственную $(-1)$-кривую $E$ в
себя и потому спускается на $\PP^2$, фиксируя $p$; обратно, автоморфизмы $\PP^2$, фиксирующие $p$,
поднимаются на раздутие. Итак, $\Aut(X)=\Stab_{\PGL_3}(p)$ --- параболическая подгруппа; при
$p=[1:0:0]$ это матрицы вида
\[
\begin{pmatrix}a&b&c\\0&d&e\\0&f&g\end{pmatrix}\!\!\Big/\C^*,\qquad\text{содержащие нормальную
унипотентную подгруппу}\ \Big\{\begin{pmatrix}1&b&c\\0&1&0\\0&0&1\end{pmatrix}\Big\}\cong\mathbb G_a^2 .
\]
Значит, $\Aut(\mathbb F_1)$ нередуктивна, и \emph{$\mathbb F_1$ не допускает КЭ-метрики}. Аналогично для
$\Bl_2\PP^2$ (стабилизатор двух точек в $\PGL_3$ содержит $\mathbb G_a$). Это ответ на упражнение~3.4(2)
обзора. Напротив, у $\PP^2$, $\PP^1\times\PP^1$, $\Bl_3\PP^2$ (группа $(\C^*)^2\rtimes(S_3\times\Z/2)$)
и у дель Пеццо степени $\le5$ (конечные группы) группы автоморфизмов редуктивны, и по теореме Тяня
(1990) все они допускают КЭ-метрики: \emph{гладкая поверхность дель Пеццо кэлер-эйнштейнова тогда и
только тогда, когда её группа автоморфизмов редуктивна}. Сравните с таблицей K-стабильности в
разделе~\ref{sec:beta}.
\end{example}

\begin{definition}[Инвариант Футаки, 1983]\label{def:Futaki-classical}
Пусть $X$ --- Фано, $\omega\in2\pi c_1(X)$, $h$ --- потенциал Риччи ($\Ric\omega-\omega=i\partial\bar\partial h$).
Для голоморфного векторного поля $v\in H^0(X,T_X)$ положим
\[
F(v)=\int_Xv(h)\,\omega^n .
\]
\end{definition}

\begin{theorem}[Футаки]\label{thm:Futaki-classical}
$F(v)$ не зависит от выбора $\omega$ в классе $2\pi c_1(X)$ и является характером алгебры Ли
$H^0(X,T_X)$. Если на $X$ есть КЭ-метрика, то $F\equiv0$.
\end{theorem}

Второе утверждение тривиально (для КЭ-метрики $h$ постоянна); первое --- вычисление. Инвариант
Футаки --- первое \emph{численное} препятствие: для $\mathbb F_1$ и $\Bl_2\PP^2$ он отличен от нуля
(см. торическую формулу в примере~\ref{ex:toric-beta}: для торических многообразий Фано $F$
пропорционален барицентру антиканонического многогранника).

\subsection{$\alpha$-инвариант Тяня и первая достаточная теорема}

\begin{definition}[Тянь, 1987]\label{def:alpha-analytic}
Для многообразия Фано $X$ с $\omega\in2\pi c_1(X)$ положим
\[
\alpha(X)=\sup\Big\{\alpha>0:\ \exists C_\alpha\ \ \int_Xe^{-\alpha(\varphi-\sup_X\varphi)}\omega^n\le C_\alpha
\ \ \forall\varphi\in\mathcal H\Big\}.
\]
\end{definition}

Это мера того, «насколько сингулярными» могут быть кэлеровы потенциалы в классе: чем больше $\alpha$,
тем лучше интегрируемость $e^{-\alpha\varphi}$. Если группа $G\subset\Aut(X)$ компактна, аналогично
определяется $\alpha_G(X)$ с $G$-инвариантными $\varphi$.

\begin{theorem}[Тянь, 1987]\label{thm:Tian-alpha-analytic}
Если $\alpha_G(X)>\dfrac n{n+1}$, то $X$ допускает ($G$-инвариантную) КЭ-метрику.
\end{theorem}

\begin{theorem}[{Демайи; приложение к \cite{CS08}}]\label{thm:Demailly}
$\alpha(X)=\inf\{\lct(X,D): D\ge0,\ D\sim_\Q-K_X\}$, где $\lct$ --- лог-канонический порог
(определение~\ref{def:lct}). Аналогично $\alpha_G(X)$ --- инфимум по $G$-инвариантным $D$.
\end{theorem}

Это первый мост между анализом и алгеброй: аналитический инвариант интегрируемости превращается
в алгебраический инвариант особенностей пар $(X,D)$, и теорема Тяня становится алгебраическим
критерием (теорема~\ref{thm:Tian-alpha}). Именно так в разделе~\ref{sec:alpha} мы докажем
K-стабильность поверхностей дель Пеццо малой степени.

\subsection{Гипотеза Яу--Тяня--Дональдсона и K-стабильность}

\begin{itemize}
\item \textbf{Дин--Тянь (1992), Тянь (1997).} Обобщённый инвариант Футаки для вырождений
$X\rightsquigarrow X_0$ с $\C^*$-действием на $X_0$; Тянь определил K-стабильность: $\Fut\ge0$ для
всех вырождений с нормальными центральными слоями и показал, что существование КЭ-метрики влечёт
K-стабильность (в подходящем смысле).
\item \textbf{Дональдсон (2002).} Чисто алгебраическое определение инварианта Футаки
тест-конфигурации через веса действия на $H^0(X_0,L_0^k)$ (определение~\ref{def:Futaki}) и
\emph{гипотеза Яу--Тяня--Дональдсона}: гладкое поляризованное многообразие $(X,L)$ допускает
метрику постоянной скалярной кривизны в классе $c_1(L)$ тогда и только тогда, когда $(X,L)$
K-полистабильно. Для Фано с $L=-K_X$ метрики постоянной скалярной кривизны в классе $c_1(-K_X)$ ---
это в точности КЭ-метрики.
\item \textbf{Чэнь--Дональдсон--Сунь, Тянь (2012--2015).} Гипотеза ЯТД верна для многообразий Фано:
\emph{гладкое многообразие Фано допускает КЭ-метрику тогда и только тогда, когда оно K-полистабильно}.
Доказательство --- метод непрерывности через метрики с коническими особенностями и теория
пределов по Громову--Хаусдорфу (Дональдсон--Сунь). Позднее появились вариационные доказательства
(Берман--Буксом--Йонссон; Ли) и обобщение на особые многообразия Фано и пары
(Ли--Тянь--Ван, Ли--Сюй--Чжуан \cite{LXZ}).
\end{itemize}

\begin{theorem}[Дональдсон--Сунь, \cite{DS14}]\label{thm:DS}
Пусть $X_i$ --- последовательность $n$-мерных КЭ многообразий Фано с фиксированным объёмом
$(-K)^n=V$. Тогда после перехода к подпоследовательности $X_i$ сходятся по Громову--Хаусдорфу к
метрическому пространству $X_\infty$, которое является нормальным проективным многообразием с
klt-особенностями и (особой) КЭ-метрикой; сходимость совместима с проективными вложениями по $|-mK|$.
\end{theorem}

Это второй мост: \emph{дифференциально-геометрическая компактификация} пространства модулей КЭ
многообразий Фано (пределы по Громову--Хаусдорфу) автоматически приводит к особым многообразиям
Фано с klt-особенностями, а по гипотезе ЯТД в особом варианте --- к K-полистабильным. Одака--Спотти--Сунь
\cite{OSS} реализовали эту программу для поверхностей дель Пеццо: для каждой степени $d$ они явно
описали громов-хаусдорфову компактификацию модулей КЭ дель Пеццо степени $d$ и доказали, что она
совпадает с алгебраическим пространством модулей K-полистабильных поверхностей; при $d=3$ это
ГИТ-фактор $\PP(1,2,3,4,5)$ из следствия~\ref{cor:GIT-cubic} (теорема~\ref{thm:GIT=K}).

\begin{example}[КЭ-метрика на кубике $3A_2$]
Поверхность $X=\{x_0x_1x_2=x_3^3\}$ --- фактор $\PP^2/\mu_3$ (пример~\ref{ex:cubic-normal-forms}(2)).
Метрика Фубини--Штуди $\mu_3$-инвариантна и спускается на $X$ как \emph{орбифолдная КЭ-метрика}
(с особенностями конического типа $\C^2/\mu_3$ в трёх точках $A_2$). Согласно \cite{OSS} и
\cite{DS14}, это и есть громов-хаусдорфов предел КЭ-метрик на гладких кубиках, вырождающихся в $X$.
Это дифференциально-геометрическое объяснение того, почему именно $3A_2$-кубика оказывается
K-полистабильной (пример~\ref{ex:3A2-polystable}).
\end{example}

\subsection{От инварианта Футаки к весам}

Последнее, что нужно, чтобы определения раздела~\ref{sec:tc} выглядели естественно. Пусть на
$(X,L)$ действует $\C^*$ (через линеаризацию на $L$), $d_k=h^0(X,L^k)$, $w_k$ --- суммарный вес
действия на $H^0(X,L^k)$ (след инфинитезимального генератора). По эквивариантной теореме
Римана--Роха $w_k$ --- многочлен от $k$ степени $n+1$, и
\[
\frac{w_k}{kd_k}=F_0+F_1k^{-1}+O(k^{-2}).
\]

\begin{theorem}[{Дональдсон, \cite[2.2]{Don02}}]\label{thm:Donaldson-Futaki}
Если $X$ гладко, $L=-K_X$ и $v$ --- голоморфное векторное поле, порождающее действие $\C^*$, то
$F_1=c_n\,F(v)$ с положительной константой $c_n$, зависящей только от $n$ и $(-K_X)^n$. Иначе
говоря, классический инвариант Футаки --- это (с точностью до множителя) субведущий коэффициент
нормированного веса.
\end{theorem}

Отсюда определение Дональдсона: для \emph{любого} вырождения $(X,L)\rightsquigarrow(X_0,L_0)$ с
$\C^*$-действием на центральном слое (тест-конфигурации) объявим $\Fut:=-F_1$, вычисленный по
действию на $H^0(X_0,L_0^k)$. Это определение имеет смысл для произвольных (даже очень особых)
$X_0$ и не использует метрик; именно оно принято в разделе~\ref{sec:tc}.

\part{K-стабильность и K-модули}

% ============================================================
\section{Особенности программы минимальных моделей}\label{sec:mmp}
% ============================================================

Этот раздел --- подробное изложение \S2 обзора \cite{DeV}. K-стабильность неожиданно тесно
связана с особенностями: K-полустабильное многообразие Фано автоматически имеет лог-терминальные
особенности (теорема~\ref{thm:Odaka}), а валюативный критерий (раздел~\ref{sec:beta}) формулируется
через лог-дискрепантности. Основной источник --- \cite[\S\,2.3]{KM}, \cite{KolSing}.

\subsection{Дискрепантности}

Напомним (определение~\ref{def:pair}): пара $(X,\Delta)$ --- нормальное многообразие с эффективным
$\Q$-дивизором $\Delta=\sum b_jD_j$, причём $K_X+\Delta$ $\Q$-картьеров.

\begin{definition}\label{def:logres}
\emph{Лог-разрешение} пары $(X,\Delta)$ --- проективный бирациональный морфизм
$\pi\colon Y\to X$ с гладким $Y$, такой что $\mathrm{Exc}(\pi)\cup\Supp\pi^{-1}_*\Delta$ ---
дивизор с простыми нормальными пересечениями (теорема~\ref{thm:Hironaka}). Здесь
$\Delta_Y:=\pi^{-1}_*\Delta=\sum b_j\tilde D_j$ --- собственный прообраз.
\end{definition}

\begin{definition}[{\cite[2.25--2.28]{KM}}]\label{def:discrep}
Пусть $(X,\Delta)$ --- пара, $\pi\colon Y\to X$ --- лог-разрешение с исключительными дивизорами
$E_1,\dots,E_r$. Поскольку $\pi$ --- изоморфизм вне исключительного множества, дивизор
$K_Y+\Delta_Y-\pi^*(K_X+\Delta)$ сосредоточен на $\bigcup E_i$ (здесь $K_Y$ и $K_X$ выбраны
согласованно: $K_Y$ --- дивизор формы $\pi^*\omega$, $K_X$ --- дивизор формы $\omega$), и мы можем
записать
\[
K_Y+\Delta_Y=\pi^*(K_X+\Delta)+\sum_{i=1}^ra_{(X,\Delta)}(E_i)\,E_i,\qquad a_{(X,\Delta)}(E_i)\in\Q .
\]
Число $a(E)=a_{(X,\Delta)}(E)$ называется \emph{дискрепантностью} дивизора $E$ относительно
пары; оно зависит только от дивизориального нормирования $\ord_E$ (т.\,е. от $E$ как дивизора над
$X$, определение~\ref{def:divisor-over}), поскольку обе части равенства --- дивизоры одной и той же
рациональной формы с поправкой. \emph{Лог-дискрепантность} --- $A_{(X,\Delta)}(E)=1+a_{(X,\Delta)}(E)$.
Для неисключительного дивизора $E=D_j$ полагают $a(D_j)=-b_j$, $A(D_j)=1-b_j$; для простого
дивизора $E\subset X$, не входящего в $\Supp\Delta$, $A(E)=1$. Наконец,
\[
\discrep(X,\Delta)=\inf\{a_{(X,\Delta)}(E):E\ \text{исключительный дивизор над }X\},
\]
где инфимум берётся по всем исключительным дивизорам на всех бирациональных моделях.
\end{definition}

\begin{definition}[{\cite[2.34]{KM}}]\label{def:klt}
Пара $(X,\Delta)$ имеет
\begin{itemize}
\item \emph{терминальные} особенности, если $\discrep(X,\Delta)>0$;
\item \emph{канонические} особенности, если $\discrep(X,\Delta)\ge0$;
\item \emph{лог-терминальные по Кавамате} (klt) особенности, если $\discrep(X,\Delta)>-1$ и все
$b_j<1$;
\item \emph{чисто лог-терминальные} (plt) особенности, если $\discrep(X,\Delta)>-1$;
\item \emph{лог-канонические} (lc) особенности, если $\discrep(X,\Delta)\ge-1$.
\end{itemize}
В терминах лог-дискрепантностей: klt $\iff$ $A(E)>0$ для всех $E$ над $X$ (включая
неисключительные), lc $\iff$ $A(E)\ge0$ для всех $E$. При $\Delta=0$ говорят об особенностях
многообразия $X$ (терминальных, канонических, klt --- в этом случае klt = plt, «лог-терминальных»).
\end{definition}

Дискрепантности --- мера «плохости» особенности: чем меньше $a(E)$, тем хуже. Иерархия:
гладкие $\subset$ терминальные $\subset$ канонические $\subset$ klt $\subset$ plt $\subset$ lc.

\begin{lemma}[Почему коэффициенты $\le1$]\label{lem:coeff-le-1}
Если какой-нибудь коэффициент $b_j>1$, то $\discrep(X,\Delta)=-\infty$. Если все $b_j\le1$, то
$\discrep(X,\Delta)\in\{-\infty\}\cup[-1,1]$.
\end{lemma}
\begin{proof}
Пусть $b_j>1$, $p\in D_j$ --- гладкая точка $X$, общая на $D_j$, и $Z\subset D_j$ --- гладкое
подмногообразие коразмерности~2 в $X$ через $p$. Раздутие $\pi_1$ центра $Z$ даёт исключительный
дивизор $E_1$ с $K_{Y_1}=\pi_1^*K_X+E_1$ (теорема~\ref{thm:blowup-K}) и $\pi_1^*D_j=\tilde D_j+E_1$,
так что $a(E_1)=1-b_j<0$. Теперь раздуем $\tilde D_j\cap E_1$ (коразмерность 2 в $Y_1$): новый
дивизор $E_2$ имеет $a(E_2)=1+a(E_1)-b_j=2(1-b_j)$; продолжая, $a(E_k)=k(1-b_j)\to-\infty$.
Второе утверждение --- \cite[2.31]{KM}: $\discrep\le1$ всегда (раздутие коразмерности 2 в гладкой
точке вне $\Delta$ даёт $a=1$), а если $\discrep<-1$, то тем же приёмом получаем $-\infty$.
\end{proof}

Как вычислять $\discrep$? Инфимум берётся по бесконечному множеству дивизоров, но его можно найти
на одном лог-разрешении.

\begin{theorem}[{\cite[2.31, 2.32]{KM}}]\label{thm:discrep-one-resolution}
Пусть $\pi\colon Y\to X$ --- лог-разрешение пары $(X,\Delta)$, $\Delta=\sum b_jD_j$ с $b_j\le1$, и
$a_i=a_{(X,\Delta)}(E_i)$ --- дискрепантности его исключительных дивизоров.
\begin{enumerate}
\item $(X,\Delta)$ lc $\iff$ все $a_i\ge-1$; $(X,\Delta)$ klt $\iff$ все $a_i>-1$ и все $b_j<1$.
\item Если все $a_i\ge-1$ и компоненты $\Delta_Y$ попарно не пересекаются, то
$\discrep(X,\Delta)=\min\{\min_ia_i,\ \min_j(1-b_j),\ 1\}$.
\end{enumerate}
\end{theorem}
\begin{proof}[Идея]
Любой дивизор $E$ над $X$ получается из $Y$ последовательностью раздутий гладких центров,
лежащих в snc-дивизоре $\mathrm{Exc}(\pi)\cup\Supp\Delta_Y$, с последующим выбором компоненты
(\cite[2.45]{KM}). Для раздутия центра $Z$ коразмерности $c$, лежащего ровно в компонентах
$F_1,\dots,F_s$ snc-дивизора с «коэффициентами» $-a(F_1),\dots,-a(F_s)$, новый исключительный
дивизор имеет дискрепантность $c-1+\sum_{l\le s}a(F_l)$ (теорема~\ref{thm:blowup-K} плюс
$\pi^*F_l=\tilde F_l+E$). Если все $a(F_l)\ge-1$ и $s\le c$, то это число $\ge-1$; индукция даёт (1).
Для (2) нужно проверить, что дискрепантности при этом не опускаются ниже указанного минимума;
условие непересечения компонент $\Delta_Y$ нужно, чтобы центр не лежал сразу в двух компонентах
границы с положительными коэффициентами.
\end{proof}

\subsection{Примеры вычисления дискрепантностей}

\begin{example}[Гладкие многообразия]
Если $X$ гладко, $\dim X=n\ge2$, то $\discrep(X)=1$: раздутие гладкого центра коразмерности $c$
даёт $a=c-1\ge1$, а дальнейшие раздутия лишь увеличивают дискрепантности. Для кривых исключительных
дивизоров нет, и по соглашению $\discrep=1$. Гладкие многообразия терминальны.
\end{example}

\begin{example}[Квадратичный конус; пример~2.5 обзора]\label{ex:quadric-cone-discrep}
Пусть $X=\{xy=z^2\}\subset\A^3$, $\pi\colon Z=\Bl_0\A^3\to\A^3$ --- раздутие начала координат с
исключительным дивизором $E\cong\PP^2$, $Y\subset Z$ --- собственный прообраз $X$,
$\pi_Y=\pi|_Y\colon Y\to X$. В карте $x=u$, $y=uv$, $z=uw$ раздутия уравнение $uv\cdot u=u^2w^2$
даёт $Y=\{v=w^2\}$ --- гладкая поверхность; аналогично в других картах. Итак, $\pi_Y$ --- разрешение
особенностей, и $E_Y:=E\cap Y$ --- исключительная кривая: в $E=\PP^2$ с координатами $[u:v:w]$
(касательные направления) она задаётся уравнением $uv=w^2$ --- гладкая коника. Её самопересечение
на $Y$: $E_Y^2=(E|_Y)^2=E\cdot E\cdot Y=(E|_E)\cdot(Y|_E)=\cO_{\PP^2}(-1)\cdot\cO_{\PP^2}(2)=-2$,
так как $\cO_E(E)=\cO_{\PP^2}(-1)$ (теорема~\ref{thm:blowup-K}) и $Y|_E=E_Y$ --- коника. (Конечно,
$X=\frac12(1,1)$, и это согласуется с примером~\ref{ex:discrep-P11n}.)

Теперь найдём $a=a_X(E_Y)$ из $K_Y=\pi_Y^*K_X+aE_Y$ ($X$ горенштейново, пример~\ref{ex:hyp-Gor},
так что $\pi_Y^*K_X$ --- обычный обратный образ). Прибавим $E_Y$ и умножим на $E_Y$:
\[
(K_Y+E_Y)\cdot E_Y=\pi_Y^*K_X\cdot E_Y+(a+1)E_Y^2 .
\]
Слева по формуле присоединения $\deg K_{E_Y}=-2$ ($E_Y\cong\PP^1$); справа
$\pi_Y^*K_X\cdot E_Y=K_X\cdot(\pi_Y)_*E_Y=0$ (формула проекции: $E_Y$ стягивается в точку), и
$E_Y^2=-2$. Итак, $-2=-2(a+1)$, $a=0$. Поскольку других исключительных дивизоров у минимального
разрешения нет, а остальные дивизоры над $X$ получаются раздутиями точек $Y$ и имеют
дискрепантности $\ge0$ (теорема~\ref{thm:discrep-one-resolution}), $\discrep(X)=0$: особенность
$A_1$ \emph{канонична}, но не терминальна (ср. теорему~\ref{thm:lt-surface}).
\end{example}

\begin{example}[Конус над эллиптической кривой; пример~2.6 обзора]\label{ex:elliptic-cone}
Пусть $C\subset\PP^2$ --- гладкая плоская кубика, $X\subset\A^3$ --- аффинный конус над ней
(особая точка кубической поверхности $\{F(x_0,x_1,x_2)=0\}\subset\PP^3$, пример~\ref{ex:cubic-normal-forms}(5)).
Раздутие вершины даёт разрешение $\pi\colon Y\to X$, где $Y$ --- тотальное пространство расслоения
$\cO_C(-1)$ (или, проективно, $\PP(\cO\oplus\cO(-1))$ над $C$), а исключительная кривая $E\cong C$
--- нулевое сечение с $E^2=-\deg\cO_C(1)=-3$. По примеру~\ref{ex:hyp-Gor} $X$ горенштейново;
пишем $K_Y=\pi^*K_X+aE$ и, как выше,
\[
(K_Y+E)\cdot E=\pi^*K_X\cdot E+(a+1)E^2,\qquad\text{т.\,е.}\qquad 0=0+(a+1)(-3),
\]
где слева $\deg K_E=2g(E)-2=0$. Итак, $a=-1$: $\discrep(X)=-1$, особенность \emph{лог-каноническая,
но не klt}. (Дальнейшие раздутия точек на $E$ дают дискрепантности $\ge-1$ по
теореме~\ref{thm:discrep-one-resolution}.) Это простая эллиптическая особенность $\tilde E_6$.
\end{example}

\begin{example}[Конусы над кривыми рода $g\ge2$; упражнение~2.1(5)]
Пусть $X$ --- конус над гладкой кривой $C$ рода $g\ge2$, вложенной расслоением $L$ степени $d$,
и предположим, что $K_X$ $\Q$-картьеров (по примеру~\ref{ex:cone-canonical} это значит
$K_C\sim_\Q rL$, $r=(2g-2)/d>0$). Тот же счёт: $2g-2=(a+1)E^2=-(a+1)d$, откуда
$a=-1-\frac{2g-2}d<-1$. Значит, $X$ \emph{не} лог-канонично. Вывод: конусы над кривыми
лог-терминальны при $g=0$, лог-каноничны, но не klt при $g=1$, и не lc при $g\ge2$ --- в полном
соответствии с «положительностью» $-K_C$.
\end{example}

\begin{example}[Конус над рациональной нормальной кривой с образующей; упражнение~2.1(4)]
Пусть $X=\PP(1,1,d)$ (точнее, его аффинная окрестность вершины $\frac1d(1,1)$), $D$ ---
образующая конуса (прямая через вершину), $Y=\mathbb F_d\to X$, $E$ --- $(-d)$-кривая,
$\tilde D=F$ --- слой. По примеру~\ref{ex:P11n-intersections} $\pi^*D=F+\frac1dE$, а по
примеру~\ref{ex:discrep-P11n} $K_Y=\pi^*K_X+(-1+\frac2d)E$. Поэтому
\[
K_Y+\tilde D=\pi^*(K_X+D)+\Big(-1+\frac2d-\frac1d\Big)E=\pi^*(K_X+D)+\Big(-1+\frac1d\Big)E .
\]
Раздутие точки $F\cap E$ даёт дивизор $E'$ с $a(E')=1+a(E)+a(\tilde D)=1+(-1+\frac1d)-1=-1+\frac1d$,
и по теореме~\ref{thm:discrep-one-resolution}(2) (компоненты границы: одна, $\tilde D$, с $b=1$)
\[
\discrep(X,D)=\min\{-1+\tfrac1d,\ 1-1,\ 1\}=-1+\tfrac1d .
\]
Пара $(X,D)$ plt (дискрепантность $>-1$), но не klt (коэффициент $b=1$).
\end{example}

\begin{example}[Терминальные особенности поверхностей гладкие; упражнение~2.1(1)]\label{ex:terminal-smooth}
Пусть $0\in X$ --- нормальная терминальная особенность поверхности, $\pi\colon Y\to X$ ---
минимальное разрешение. По лемме~\ref{lem:minimal-discrep} все дискрепантности исключительных
кривых $\le0$, а терминальность требует $>0$; значит, исключительных кривых нет, $\pi$ ---
изоморфизм, и $X$ гладка. Так же доказывается, что канонические особенности поверхностей ---
в точности особенности Дю Валя: из $a_i\le0$ и $a_i\ge0$ следует $K_Y=\pi^*K_X$, т.\,е. $K_Y\cdot E_i=0$,
$E_i^2=-2$, $p_a(E_i)=0$.
\end{example}

\subsection{Лог-канонический порог}

\begin{definition}[{\cite[8.1]{KolSing}}]\label{def:lct}
Пусть $X$ --- нормальное многообразие с lc-особенностями и $\Q$-картьеровым $K_X$, $D\ne0$ ---
эффективный $\Q$-картьеров дивизор. \emph{Лог-канонический порог}
\[
\lct(X,D)=\sup\{c\ge0:\ (X,cD)\ \text{лог-канонична}\}.
\]
Если $X$ ясно из контекста, пишут $\lct(D)$. Более общо, для пары $(X,\Delta)$ и $D$ определяется
$\lct(X,\Delta;D)=\sup\{c:(X,\Delta+cD)\text{ lc}\}$.
\end{definition}

По лемме~\ref{lem:coeff-le-1} $\lct(X,D)\le1/\max_jd_j$, где $d_j$ --- коэффициенты $D$; в
частности, $\lct\le1$ для целочисленных $D$. Супремум достигается (множество $c$, при которых пара
lc, замкнуто), т.\,е. $(X,\lct(D)\cdot D)$ --- lc, а при большем коэффициенте --- нет.

\begin{proposition}[Вычисление на лог-разрешении]\label{prop:lct-formula}
Пусть $\pi\colon Y\to X$ --- лог-разрешение пары $(X,D)$, $K_Y=\pi^*K_X+\sum a_iE_i$ и
$\pi^*D=\sum m_iE_i$ (суммы по всем простым дивизорам $Y$: исключительным и компонентам $\tilde D$,
для которых $a_i=0$). Тогда
\[
\lct(X,D)=\min_i\frac{a_i+1}{m_i}=\min_{E\ \text{над}\ X}\frac{A_X(E)}{\ord_E(D)} .
\]
\end{proposition}
\begin{proof}
$K_Y=\pi^*(K_X+cD)+\sum(a_i-cm_i)E_i$, и по теореме~\ref{thm:discrep-one-resolution}(1)
пара $(X,cD)$ lc $\iff$ $a_i-cm_i\ge-1$ для всех $i$ $\iff$ $c\le(a_i+1)/m_i$ для всех $i$.
Второе равенство --- та же теорема: минимум по всем дивизорам над $X$ достигается уже на
лог-разрешении.
\end{proof}

\begin{corollary}\label{cor:lct-mult}
Если $X$ гладко в точке $p$, $\dim X=n$, и $D$ --- гиперповерхность кратности $m=\mult_pD$, то
$\frac1m\le\lct_p(X,D)\le\frac nm$, где $\lct_p$ --- порог в окрестности $p$. В частности, если
$D$ гладка в $p$, то $\lct_p=1$.
\end{corollary}
\begin{proof}
Верхняя оценка: для исключительного дивизора $E$ раздутия точки $p$ имеем $A(E)=n$ и
$\ord_E(D)=m$ (пример~\ref{ex:blowup-point-ord}), так что $\lct_p\le n/m$ по
предложению~\ref{prop:lct-formula}. Нижняя оценка --- \cite[8.10]{KolPairs}: пара $(X,D)$
лог-канонична в гладкой точке $p$, если $\mult_pD\le1$ (индукция по размерности с помощью общих
гиперплоских сечений и обращения присоединения); применяя это к $\frac1mD$, получаем $\lct_p\ge1/m$.
\end{proof}

\begin{example}[Плоские кривые; упражнение~2.1(2)]\label{ex:lct-curves}
Для кривых $C=\{f=0\}\subset\A^2$ удобны \emph{взвешенные раздутия}: торическое раздутие с весами
$(w_1,w_2)$ (взаимно простыми) даёт дивизор $E$ над $0$ с $A(E)=w_1+w_2$ и
$\ord_E(f)=\min\{w_1i+w_2j: x^iy^j\in f\}$ (пример в \S\ref{ss:divisors-over}), откуда
$\lct\le(w_1+w_2)/\ord_E(f)$. Для невырожденных относительно многоугольника Ньютона $f$ (в
частности, для всех примеров ниже) эта оценка точна при подходящих весах \cite[\S\,8.3]{KolSing}.
\begin{enumerate}
\item \emph{Касп} $y^2=x^3$: веса $(2,3)$, $\ord_E=6$, $A=5$: $\lct=5/6$.
\item \emph{Тэкнода} $y^2=x^4$: веса $(1,2)$, $\ord_E=4$, $A=3$: $\lct=3/4$.
\item $y^2=x^n$: при нечётном $n$ веса $(2,n)$, $\ord_E=2n$, $A=n+2$; при чётном --- $(1,n/2)$,
$\ord_E=n$, $A=1+n/2$. В обоих случаях $\lct=\frac12+\frac1n$. При $n=2$ (узел) получаем $1$.
\item $n$ прямых через начало координат: обычное раздутие, $\ord_E=n$, $A=2$: $\lct=\min\{1,2/n\}$
(при $n\le2$ кривая snc, и порог равен 1; при $n\ge3$ --- $2/n$).
\item \emph{Гладкая кривая и узловая кривая}: $\lct=1$ (snc-пары lc).
\end{enumerate}
Чем «хуже» особенность кривой, тем меньше порог: $1$ (узел) $>\frac56$ (касп) $>\frac34$ (тэкнода)
$>\frac23$ (три прямые через точку) $>\dots$
\end{example}

\begin{example}[Гиперплоские сечения кубической поверхности]\label{ex:lct-cubic-sections}
Пусть $X\subset\PP^3$ --- гладкая кубика, $H$ --- плоскость, $D=H\cap X$ --- плоская кубика на
$X$. Так как $X$ гладка, $\lct(X,D)$ вычисляется по особенностям кривой $D$: если $H$ не касается
$X$, $D$ гладка и $\lct=1$; если $H=\overline T_pX$ --- касательная плоскость в общей точке $p$, то
$D$ --- узловая кубика ($\lct=1$); в точках, где $D$ каспидальна, $\lct=5/6$; если $H$ ---
тритангенциальная плоскость, $D$ --- треугольник ($\lct=1$); наконец, если $p$ --- \emph{точка
Эккардта} (три прямые на $X$ через $p$), то $D$ --- три прямые через точку и $\lct=2/3$. Других особенностей у плоских сечений гладкой кубики не бывает (кроме узлов,
каспов, тэкнод $A_3$ --- «прямая плюс касающаяся её коника», $\lct=3/4$, --- и перечисленных выше). По
теореме Чельцова (теорема~\ref{thm:Cheltsov}) $\inf_H\lct(X,H\cap X)$ равен $2/3$, если на $X$ есть
точка Эккардта, и $3/4$ иначе; это и есть $\alpha$-инвариант кубической поверхности.
\end{example}

\begin{remark}[Аналитическая интерпретация]
Для $X$ гладкого и $D=\{f=0\}$: $\lct_p(D)=\sup\{c:|f|^{-2c}\text{ локально интегрируема около }p\}$
(\cite[8.5]{KolPairs}). В таком виде порог появляется у Демайи в теореме~\ref{thm:Demailly}:
$\alpha$-инвариант Тяня --- это инфимум порогов интегрируемости по всем антиканоническим
«потенциалам».
\end{remark}

% ============================================================
\section{K-стабильность через тест-конфигурации}\label{sec:tc}
% ============================================================

Это подробное изложение \S3.1 обзора. Мы рассматриваем нормальные проективные многообразия и, как
и в обзоре, не ограничиваемся гладким случаем.

\subsection{Вырождения вдоль однопараметрических подгрупп}

Пусть $(X,L)$ --- \emph{поляризованное} многообразие: $X$ нормально, проективно, $\dim X=n$, $L$ ---
обильное линейное расслоение. При $m\gg0$ полная линейная система $|L^m|$ задаёт вложение
$X\hookrightarrow\PP^N=\PP(H^0(X,L^m)^*)$. Пусть $\lambda\colon\Gm\to\GL_{N+1}$ --- однопараметрическая
подгруппа. По \S\ref{sec:families} существует плоский предел
\[
X_0=\lim_{t\to0}\lambda(t)\cdot X\in\Hilb(\PP^N),
\]
и семейство $\cX=\overline{\{(x,t):x\in\lambda(t)X,\ t\ne0\}}\subset\PP^N\times\A^1$ плоско над $\A^1$
со слоями $X_t\cong X$ при $t\ne0$ и $X_0$ при $t=0$. Группа $\Gm$ действует на $\cX$ по правилу
$s\cdot(x,t)=(\lambda(s)x,st)$ (это переводит слой над $t$ в слой над $st$), а значит, и на
центральном слое $X_0$ и на $\cL_0=\cO_{\PP^N}(1)|_{X_0}$.

\begin{example}[Коника вырождается в пару прямых]\label{ex:conic-tc}
Пусть $X=\{xz=y^2\}\subset\PP^2$ --- коника ($X\cong\PP^1$, $L=\cO_{\PP^1}(2)$) и
$\lambda(t)=\mathrm{diag}(1,t^{-1},1)$. Тогда $\lambda(t)X=\{(x,t^{-1}y,z):xz=y^2\}=\{xz=t^2y^2\}$,
$\cX=\{xz=t^2y^2\}\subset\PP^2\times\A^1$, и $X_0=\{xz=0\}$ --- пара прямых, пересекающихся в
точке $[0:1:0]$. Центральный слой приведён, но \emph{не нормален}. Действие $s\cdot(x,y,z,t)=
(x,s^{-1}y,z,st)$ сохраняет $\cX$ и на $X_0$ сводится к $y\mapsto s^{-1}y$.
\end{example}

\begin{example}[Деформация к нормальному конусу]\label{ex:normal-cone}
Пусть $p\in X$, $\cX=\Bl_{p\times0}(X\times\A^1)$ с исключительным дивизором $\mathcal E$, и
$\cL_c=\pi^*L-c\mathcal E$ при рациональном $c>0$ (для целочисленности заменяем $L$ и $\cL_c$ на
кратные). При $0<c<\epsilon(L;p)$ (постоянная Сешадри) $\cL_c$ относительно обилен, и
$\cX\to\A^1$ --- семейство с $X_t\cong X$ при $t\ne0$ и $X_0=\tilde X\cup\mathcal E$, где
$\tilde X=\Bl_pX$ и $\mathcal E\cong\PP^n$ (проективизация $T_pX\oplus\C$), склеенные по
исключительному дивизору $\tilde X$; $\Gm$ действует на $\A^1$ умножением, тривиально на $X$.
Это вырождение не задаётся 1-ПП в фиксированном $\PP^N$ --- в нём участвует «переход к высшим
степеням $L$»; см. замечание~\ref{rem:tc-vs-1ps}.
\end{example}

\subsection{Определение тест-конфигурации}

\begin{definition}[{Тянь, Дональдсон; \cite[2.1]{Xu}}]\label{def:tc}
\emph{Тест-конфигурация} поляризованного многообразия $(X,L)$ --- это набор $(\cX,\cL)$, где
\begin{enumerate}
\item $\cX$ --- нормальное многообразие с действием $\Gm$ и $\Gm$-эквивариантным плоским проективным
морфизмом $\pi\colon\cX\to\A^1$ ($\Gm$ действует на $\A^1$ умножением);
\item $\cL$ --- $\pi$-обильное $\Gm$-линеаризованное линейное расслоение (или $\Q$-расслоение) на
$\cX$;
\item задан $\Gm$-эквивариантный изоморфизм $(\cX,\cL)|_{\pi^{-1}(\Gm)}\cong(X\times\Gm,\ p_1^*L^{\otimes r})$,
где $\Gm$ действует на $X\times\Gm$ умножением на второй сомножитель; число $r>0$ --- \emph{показатель}
тест-конфигурации.
\end{enumerate}
Слой $X_0=\pi^{-1}(0)$ с $\cL_0=\cL|_{X_0}$ --- \emph{центральный слой}; на нём действует $\Gm$.
Тест-конфигурация \emph{тривиальна}, если $\cX\cong X\times\A^1$ с тривиальным действием $\Gm$ на
$X$; она \emph{произведение} (product test configuration), если $\cX\cong X\times\A^1$ с действием
$s\cdot(x,t)=(\sigma(s)x,st)$ для некоторого $\sigma\colon\Gm\to\Aut(X,L)$.
\end{definition}

Условие нормальности $\cX$ существенно (Ли--Сюй): без него возникают «патологические»
тест-конфигурации с нулевым инвариантом Футаки, но нетривиальным центральным слоем. Переход к
нормализации не меняет общих слоёв, но меняет центральный (нормализует его, грубо говоря, и
заменяет неприведённые компоненты).

\begin{remark}[Тест-конфигурации и однопараметрические подгруппы]\label{rem:tc-vs-1ps}
Конструкция начала параграфа даёт тест-конфигурацию (после нормализации) показателя $m$; обратно,
всякая тест-конфигурация получается так: при $k\gg0$ прямой образ $\pi_*\cL^k$ --- свободный
$\C[t]$-модуль с $\Gm$-действием, его тривиализация задаёт вложение $\cX\hookrightarrow\PP^N\times\A^1$
и 1-ПП $\lambda$ \cite[2.?]{Xu}. Разница с ГИТ (замечание~3.2 обзора): в ГИТ фиксировано объемлющее
$\PP(V)$, а здесь можно переходить к сколь угодно большим степеням $L^m$ и к большим $\PP^N$.
Поэтому условий K-стабильности «больше», чем условий ГИТ-стабильности в любом фиксированном
вложении, и, вообще говоря, K-стабильность сильнее (ср. теорему~\ref{thm:PT}).
\end{remark}

\subsection{Веса и инвариант Футаки}

Пусть $(\cX,\cL)$ --- тест-конфигурация $(X,L)$ показателя $r$; заменяя $L$ на $L^r$, можно
считать $r=1$. Для $k\gg0$ положим
\[
d_k=h^0(X,L^k)=h^0(X_0,\cL_0^k)
\]
(равенство --- плоскость и обращение в нуль высших когомологий при $k\gg0$). По асимптотической
формуле Римана--Роха (теорема~\ref{thm:aRR})
\[
d_k=a_0k^n+a_1k^{n-1}+O(k^{n-2}),\qquad a_0=\frac{L^n}{n!},\quad a_1=-\frac{L^{n-1}\cdot K_X}{2(n-1)!} .
\]
Группа $\Gm$ действует на конечномерном пространстве $H^0(X_0,\cL_0^k)$; пусть $w_k$ --- сумма
весов (след инфинитезимального генератора). Соглашение о знаке: $\Gm$ действует на сечениях
естественным образом, $(s\cdot\sigma)(x)=s\cdot\sigma(s^{-1}x)$; при этом функция $t$ на $\A^1$
имеет вес $-1$ (так как $(s\cdot t)(x)=t(s^{-1}x)=s^{-1}t$).

\begin{proposition}[{\cite[2.?]{Xu}, \cite{Don02}}]\label{prop:weight-poly}
При $k\gg0$ число $w_k$ --- многочлен от $k$ степени $\le n+1$:
$w_k=b_0k^{n+1}+b_1k^n+O(k^{n-1})$.
\end{proposition}
\begin{proof}[Идея]
Склеим $\cX$ с тривиальным семейством $X\times(\PP^1\setminus0)$ вдоль $\pi^{-1}(\Gm)\cong X\times\Gm$;
получится $\Gm$-эквивариантное проективное семейство $\bar\cX\to\PP^1$ с расслоением $\bar\cL$
(\emph{компактифицированная тест-конфигурация}). По эквивариантной теореме Римана--Роха
$w_k=\chi_{\Gm}$-вес, и
\[
w_k=\frac{\bar\cL^{\,n+1}}{(n+1)!}k^{n+1}-\frac{\bar\cL^{\,n}\cdot K_{\bar\cX/\PP^1}}{2\,n!}k^n+O(k^{n-1}),
\]
где индексы пересечения берутся на $(n+1)$-мерном многообразии $\bar\cX$ \cite[2.1.2]{Xu}.
\end{proof}

\begin{definition}[Инвариант Футаки]\label{def:Futaki}
Запишем
\[
F(k):=\frac{w_k}{k\,d_k}=F_0+F_1k^{-1}+F_2k^{-2}+\dots\qquad(k\to\infty).
\]
\emph{(Обобщённый) инвариант Футаки} тест-конфигурации --- $\Fut(\cX,\cL):=-F_1$. Явно:
\[
\frac{b_0k^{n+1}+b_1k^n+\dots}{a_0k^{n+1}+a_1k^n+\dots}=\frac{b_0}{a_0}+\frac{a_0b_1-a_1b_0}{a_0^2}k^{-1}+\dots,
\qquad\Fut(\cX,\cL)=\frac{a_1b_0-a_0b_1}{a_0^2}.
\]
\end{definition}

Инвариант не меняется при замене $L$ на $L^m$ (с соответствующей заменой $\cL$) и однороден
степени~1 по действию $\Gm$: замена $\lambda\mapsto\lambda^m$ умножает $w_k$, а с ним и $\Fut$,
на $m$. Для гладких
$X_0$ с $L=-K$ он совпадает (с точностью до положительного множителя) с классическим инвариантом
Футаки векторного поля, порождающего действие (теорема~\ref{thm:Donaldson-Futaki}). Для тривиальной
тест-конфигурации $w_k=0$ и $\Fut=0$.

\begin{example}[Произведение $\PP^1$]\label{ex:product-P1}
Пусть $X=\PP^1$, $L=\cO(2)$, $\cX=X\times\A^1$ с действием $s\cdot[x:y]=[sx:y]$. На
$H^0(\cO(2k))$ --- мономы $x^ay^{2k-a}$, $0\le a\le2k$, с весами $a$ (с точностью до общего знака).
Тогда $d_k=2k+1$, $w_k=\sum_{a=0}^{2k}a=k(2k+1)$, $F(k)=1$, $F_1=0$: $\Fut=0$. Это общее явление
для $\PP^n$: инвариант Футаки всех произведений равен нулю (у $\PP^n$ есть КЭ-метрика, и классический
инвариант Футаки обращается в нуль, теорема~\ref{thm:Futaki-classical}).
\end{example}

\begin{example}[Коника и пара прямых]\label{ex:conic-Futaki}
В примере~\ref{ex:conic-tc} $X_0=\{xz=0\}$, $\cL_0=\cO(1)$, и
$H^0(X_0,\cO(k))=(\C[x,y,z]/(xz))_k$ имеет базис $y^bz^{k-b}$ ($0\le b\le k$) и $x^{k-b}y^b$
($0\le b\le k-1$): $d_k=2k+1$ --- столько же, сколько у коники, как и должно быть. Действие
$s\cdot(x,y,z)=(x,s^{-1}y,z)$ на точках даёт на функциях $y\mapsto sy$, т.\,е. моном с $y^b$ имеет
вес $b$. Итак,
\[
w_k=\sum_{b=0}^kb+\sum_{b=0}^{k-1}b=k^2,\qquad F(k)=\frac{k^2}{k(2k+1)}=\frac k{2k+1}=\frac12-\frac1{4k}+\dots,
\]
$F_1=-1/4$ и $\Fut=1/4>0$. Вырождение коники в пару прямых «дестабилизирует» в правильную сторону:
$\PP^1$ K-полустабильно, и все его тест-конфигурации имеют $\Fut\ge0$ (хотя эта тест-конфигурация
не нормальна, и строго говоря, в определении~\ref{def:Kstab} она не участвует: её нормализация ---
две непересекающиеся прямые, что уже не тест-конфигурация $\PP^1$; см. обсуждение в \cite[2.1]{Xu}).
\end{example}

\begin{example}[Нормальный конус точки на $\PP^1$]\label{ex:normal-cone-Futaki}
Пусть $X=\PP^1$, $L=\cO(1)$, $p=\{x=0\}$, $\cX=\Bl_{(p,0)}(\PP^1\times\A^1)$,
$\cL=\pi^*\cO(2)-\mathcal E$ (это $2\cL_{1/2}$ в обозначениях примера~\ref{ex:normal-cone}; заменим
$L$ на $\cO(2)$). Сечения $\cL^{\otimes m}$ на $\cX$ --- это многочлены $f=\sum_it^ig_i(x,y)$,
$g_i\in H^0(\cO(2m))$, обращающиеся в нуль с порядком $\ge m$ в точке $(p,0)$, т.\,е. $\ord_pg_i\ge m-i$.
Модуль $H^0(X_0,\cL_0^m)=H^0(\cX,\cL^m)/t\cdot H^0(\cX,\cL^m)$ порождён классами мономов $x^at^i$
(где $x^a$ означает $x^ay^{2m-a}$) с $a+i\ge m$, не делящихся на $t$ «внутри модуля»: это
$x^at^{m-a}$ при $0\le a\le m$ и $x^a$ при $m<a\le2m$ --- всего $2m+1$ элементов, как и положено.
Вес $t^i$ равен $-i$, вес $x^a$ равен $0$. Итак,
\[
w_m=-\sum_{a=0}^m(m-a)=-\frac{m(m+1)}2,\qquad
F(m)=\frac{-m(m+1)/2}{m(2m+1)}=-\frac14-\frac1{8m}+\dots,
\]
$F_1=-1/8$, $\Fut=1/8>0$. Центральный слой $X_0=\tilde X\cup\mathcal E$ --- две прямые, склеенные в
точке, с $\deg\cL_0=1$ на каждой.
\end{example}

\subsection{Многообразия Фано и специальные тест-конфигурации}

Пусть теперь $X$ --- $\Q$-Фано: нормальное проективное с klt-особенностями и обильным $-K_X$. В
качестве поляризации берём $L=-mK_X$ ($m$ --- такое, что $-mK_X$ --- дивизор Картье; инвариант
Футаки от выбора $m$ не зависит). При $L=-mK_X$ имеем $a_0=\frac{m^n(-K_X)^n}{n!}$,
$a_1=\frac{m^{n-1}(-K_X)^n}{2(n-1)!}$ (теорема~\ref{thm:aRR}), и формула
предложения~\ref{prop:weight-poly} даёт \cite[2.17]{Xu}:
\begin{equation}\label{eq:Fut-Fano}
\Fut(\cX,\cL)=\frac1{2(-K_X)^n}\Big(\frac1{m^n}\bar\cL^{\,n}\cdot K_{\bar\cX/\PP^1}+\frac n{n+1}\cdot\frac1{m^{n+1}}\bar\cL^{\,n+1}\Big).
\end{equation}

\begin{definition}[{\cite[LX14]{LX}}]\label{def:special-tc}
Тест-конфигурация $(\cX,\cL)$ многообразия Фано $X$ называется \emph{специальной}, если
$\cL\sim_\Q-mK_{\cX/\A^1}$ (т.\,е. поляризация антиканоническая) и центральный слой $X_0$ ---
$\Q$-Фано многообразие (нормальное с klt-особенностями). Иначе говоря, $\cX\to\A^1$ ---
$\Q$-горенштейново семейство $\Q$-Фано многообразий с $\Gm$-действием.
\end{definition}

Для специальной тест-конфигурации формула~\eqref{eq:Fut-Fano} с $\bar\cL=-mK_{\bar\cX/\PP^1}$ даёт
\begin{equation}\label{eq:Fut-special}
\Fut(\cX,\cL)=-\frac{1}{2(n+1)(-K_X)^n}\big(-K_{\bar\cX/\PP^1}\big)^{n+1}
\end{equation}
(замечание~3.7 обзора): \emph{знак инварианта Футаки специальной тест-конфигурации определяется знаком
самопересечения относительного антиканонического класса компактифицированного семейства}.

\begin{theorem}[{Ли--Сюй, \cite{LX}}]\label{thm:LX}
Для проверки K-(полу/поли)стабильности $\Q$-Фано многообразия достаточно рассматривать только
специальные тест-конфигурации: для любой тест-конфигурации $(\cX,\cL)$ существует специальная
$(\cX^s,-mK_{\cX^s})$ с $\Fut(\cX^s)\le c\cdot\Fut(\cX)$ для некоторого $c>0$, причём равенство
возможно лишь в хорошо контролируемых случаях.
\end{theorem}

\begin{proof}[Идея]
Это применение программы минимальных моделей (ПММ) к тотальному пространству $\cX$: после конечной
замены базы $t\mapsto t^d$ и нормализации, lc-модификации и серии бирациональных операций ПММ
(дивизориальных стягиваний и флипов) центральный слой становится klt-многообразием Фано, а
инвариант Футаки на каждом шаге не возрастает. Детали --- в \cite{LX}, \cite[гл.~2]{Xu}.
\end{proof}

\begin{definition}[Тянь, Дональдсон]\label{def:Kstab}
Пусть $X$ --- $\Q$-Фано многообразие. $X$ называется
\begin{enumerate}
\item \emph{K-полустабильным}, если $\Fut(\cX,\cL)\ge0$ для всех специальных тест-конфигураций;
\item \emph{K-полистабильным}, если оно K-полустабильно и $\Fut(\cX,\cL)=0$ лишь для
произведений ($\cX\cong X\times\A^1$);
\item \emph{K-стабильным}, если оно K-полустабильно и $\Fut(\cX,\cL)=0$ лишь для тривиальной
тест-конфигурации;
\item \emph{K-нестабильным} (неустойчивым), если не K-полустабильно.
\end{enumerate}
По теореме~\ref{thm:LX} те же понятия получатся, если брать все (нормальные) тест-конфигурации.
\end{definition}

\begin{remark}\label{rem:BX}
\begin{enumerate}
\item Если $X$ K-стабильно, то $\Aut(X)$ конечна (иначе нетривиальное произведение имело бы
$\Fut=0$: инвариант Футаки произведения $\sigma$ и произведения $\sigma^{-1}$ противоположны по
знаку, а оба $\ge0$). Обратно, если $\Aut(X)$ конечна, то K-стабильность и K-полистабильность
совпадают \cite[BX19]{BX}. Аналогия с ГИТ: полистабильность --- «замкнутость орбиты»,
стабильность --- замкнутость плюс конечность стабилизатора.
\item \textbf{Красный флаг} (замечание~3.10 обзора). Чтобы проверить K-полустабильность по
определению, нужно перебрать \emph{все} специальные тест-конфигурации --- все $\Q$-горенштейновы
вырождения $X$ в klt-многообразия Фано с $\Gm$-действием. Это бесконечный и плохо обозримый набор
(как если бы в ГИТ нужно было проверять вес Гильберта--Мамфорда во всех вложениях сразу). Разделы
\ref{sec:alpha}--\ref{sec:az} посвящены тому, как с этим бороться.
\end{enumerate}
\end{remark}

\begin{theorem}[{Одака, \cite{Oda13}}]\label{thm:Odaka}
Пусть $X$ --- нормальное проективное многообразие с $\Q$-картьеровым обильным $-K_X$. Если $X$
K-полустабильно (относительно всех тест-конфигураций), то $X$ имеет klt-особенности.
\end{theorem}

\begin{proof}[Идея]
Если $X$ не klt, то по lc-модификации/разрешению строится специальная «дестабилизирующая»
тест-конфигурация (деформация к нормальному конусу подходящего идеала), у которой инвариант Футаки
отрицателен --- вычисление сводится к отрицательности подходящей дискрепантности.
\end{proof}

Это первое указание на связь K-стабильности с особенностями: условие на инвариант Футаки,
в котором нет ни слова об особенностях, влечёт лог-терминальность.

\subsection{Примеры специальных тест-конфигураций}

\begin{example}[Вырождение $\PP^2$ в $\PP(1,1,4)$]\label{ex:P2-to-P114}
Семейство $\cX=\{xy=z^2+tw\}\subset\PP(1,1,1,2)\times\A^1_t$ из \S\ref{sec:families} с действием
$s\cdot[x:y:z:w;t]=[x:y:z:s^{-1}w;st]$ --- специальная тест-конфигурация $\PP^2$ (поляризация
$\cO(3)=-K$): $\cX$ нормально, $K_{\cX/\A^1}$ --- $\Q$-Картье, центральный слой
$\PP(1,1,4)$ --- klt (особенность $\frac14(1,1)$). Так как $\PP^2$ K-полистабильно
(пример~\ref{ex:Pn-polystable}) и это не произведение, $\Fut>0$. Вычислить его можно по
формуле~\eqref{eq:Fut-special} или через $\beta$-инвариант (раздел~\ref{sec:beta}): специальной
тест-конфигурации отвечает дивизориальное нормирование $\ord_{X_0}|_{\C(X)}$ на $X$, и
$\Fut=c\cdot\beta(\ord_{X_0})$, $c>0$ (Фудзита \cite{Fuj19b}).
\end{example}

\begin{example}[Вырождение кубики в $3A_2$]\label{ex:cubic-to-3A2}
Пусть $X\subset\PP^3$ --- кубика с особенностью $A_2$ (или гладкая кубика, вырождающаяся к такой).
Из доказательства теоремы~\ref{thm:Hilbert}(3): $\lambda(t)=\mathrm{diag}(t,t,1,t^{-2})$ переводит
$F=x_3x_0x_1+cx_2^3+C'$ в семейство с центральным слоем $x_0x_1x_3+cx_2^3=0$. Это специальная
тест-конфигурация: тотальное пространство --- гиперповерхность в $\PP^3\times\A^1$ (нормальная,
горенштейнова), центральный слой --- $3A_2$-кубика, klt. Инвариант Футаки $\ge0$, если $X$
K-полустабильна; мы увидим (пример~\ref{ex:3A2-polystable}), что $3A_2$-кубика K-полистабильна, а
все кубики с особенностями $A_1,A_2$ K-полустабильны.
\end{example}

\begin{example}[Произведения и торические многообразия]\label{ex:toric-product-tc}
Пусть $X=X_P$ --- гладкое торическое многообразие Фано с антиканоническим многогранником
$P=P_{-K}\subset M_\R$ \cite[4.22]{AG}, $u\in N$ --- однопараметрическая подгруппа тора. Произведение
$\cX=X\times\A^1$ с действием $\lambda^u$ --- тест-конфигурация, и её инвариант Футаки
вычисляется по формуле Дональдсона \cite{Don02}: сечения $H^0(-kK_X)$ --- характеры $\chi^m$,
$m\in kP\cap M$, с весами $\langle m,u\rangle$, поэтому
\[
d_k=\#(kP\cap M)=\vol(P)k^n+\tfrac12\vol(\partial P)k^{n-1}+\dots,\qquad
w_k=\sum_{m\in kP\cap M}\langle m,u\rangle=k^{n+1}\int_P\langle x,u\rangle\,dx+\tfrac12k^n\int_{\partial P}\langle x,u\rangle\,d\sigma+\dots,
\]
где $\vol$ --- евклидов объём, а $d\sigma$ --- мера на границе, нормированная решёткой (вес
характера $\chi^m$ равен $\langle m,u\rangle$ с точностью до общего знака). Для рефлексивного $P$
(грани на решёточном расстоянии $1$ от начала координат) разбиение $P$ на конусы над гранями даёт
$\vol(\partial P)=n\vol(P)$ и $\int_{\partial P}\langle x,u\rangle d\sigma=(n+1)\int_P\langle x,u\rangle dx$.
Подставляя $a_0=\vol P$, $a_1=\frac n2\vol P$, $b_0=\vol(P)\langle b_P,u\rangle$,
$b_1=\frac{n+1}2\vol(P)\langle b_P,u\rangle$, где $b_P=\frac1{\vol P}\int_Px\,dx$ ---
\emph{барицентр} $P$, в определение~\ref{def:Futaki}, получаем
\[
\Fut(X\times\A^1,\lambda^u)=\frac{a_1b_0-a_0b_1}{a_0^2}=\frac{n-(n+1)}2\langle b_P,u\rangle=-\frac12\langle b_P,u\rangle .
\]
Вывод (Ван--Чжу \cite{WZ}): торическое многообразие Фано K-полистабильно (равносильно: допускает
КЭ-метрику) тогда и только тогда, когда барицентр антиканонического многогранника --- начало
координат; если $b_P\ne0$, то, беря $u$ с $\langle b_P,u\rangle>0$, получаем произведение с
$\Fut<0$, и $X$ K-нестабильно. Мы вернёмся к этому в примере~\ref{ex:toric-beta}, где та же величина
возникнет из $\beta$-инварианта, и там же приведём барицентры торических поверхностей дель Пеццо.
\end{example}

\subsection{Связь с ГИТ для гиперповерхностей}

\begin{theorem}[{Пол--Тянь, \cite{PT}}]\label{thm:PT}
Пусть $X\subset\PP^N$ --- гиперповерхность степени $d$ с $-K_X=\cO_X(N+1-d)$ обильным, и
$\lambda\colon\Gm\to\SL_{N+1}$ --- 1-ПП. Тест-конфигурация $(\cX_\lambda,\cL_\lambda)$, индуцированная
$\lambda$ (с поляризацией $\cO(1)$), имеет инвариант Футаки
\[
\Fut(\cX_\lambda,\cL_\lambda)=a\cdot\mu(F,\lambda),
\]
где $\mu(F,\lambda)$ --- вес Гильберта--Мамфорда (с подходящим знаком) точки $[F]\in\PP(\Sym^d)$,
а $a>0$ зависит только от $N$ и $d$. Иначе говоря, инвариант Футаки --- это вес действия
на \emph{CM-расслоении} (расслоении Чжоу--Мамфорда) на пространстве гиперповерхностей, а оно
для гиперповерхностей Фано --- положительная степень $\cO(1)$.
\end{theorem}

Следствие: \emph{K-полустабильная гиперповерхность Фано ГИТ-полустабильна} (все веса
Гильберта--Мамфорда $\ge0$, так как они пропорциональны инвариантам Футаки конкретных
тест-конфигураций), и аналогично для стабильности и полистабильности. Обратное неверно в общей
ситуации, но верно для кубических поверхностей (теорема~\ref{thm:GIT=K}), кубических трёхмерных и
четырёхмерных многообразий.

% ============================================================
\section{$\alpha$-инвариант Тяня}\label{sec:alpha}
% ============================================================

Это \S3.2 обзора. $\alpha$-инвариант --- первый практический инструмент: он даёт \emph{достаточное}
условие K-стабильности, которое часто удаётся проверить.

\subsection{Определение и первые примеры}

\begin{definition}[{Тянь \cite{Tia87}; алгебраическая форма --- теорема~\ref{thm:Demailly}}]\label{def:alpha}
Пусть $X$ --- $\Q$-Фано многообразие. \emph{$\alpha$-инвариант} ---
\[
\alpha(X)=\inf\{\lct(X,D):\ D\ge0,\ D\sim_\Q-K_X\}.
\]
Если на $X$ действует алгебраическая группа $G$, то
$\alpha_G(X)=\inf\{\lct(X,D):\ D\ge0,\ D\sim_\Q-K_X,\ D\ \text{$G$-инвариантен}\}$ (инфимум по
пустому множеству равен $+\infty$).
\end{definition}

Инфимум берётся по $\Q$-дивизорам $D=\frac1mY$, $Y\in|-mK_X|$, $m\ge1$, и, поскольку
$\lct(X,\frac1mY)=m\,\lct(X,Y)$,
\[
\alpha(X)=\inf_{m\ge1}\ \inf_{Y\in|-mK_X|}m\cdot\lct(X,Y)
=\inf_{m}\ \inf_{Y\in|-mK_X|}\ \min_{E}\frac{m\,A_X(E)}{\ord_E(Y)} .
\]
Чем больше $\alpha(X)$, тем «менее особыми» могут быть антиканонические дивизоры; это глобальная
мера «положительности» $-K_X$.

\begin{example}[$\alpha(\PP^n)=\frac1{n+1}$]\label{ex:alpha-Pn}
Так как $-K_{\PP^n}=(n+1)H$, дивизор $D=(n+1)H$ (гиперплоскость с коэффициентом $n+1$) даёт
$\lct(\PP^n,D)=\frac1{n+1}\lct(H)=\frac1{n+1}$. Обратно, пусть $D=\frac1mY$, $Y$ --- гиперповерхность
степени $m(n+1)$. В любой точке $p$ кратность $\mult_pY\le\deg Y=m(n+1)$ (пересечём $Y$ с общей
прямой через $p$), поэтому по следствию~\ref{cor:lct-mult} $\lct(Y)\ge\frac1{m(n+1)}$ и
$\lct(D)=m\,\lct(Y)\ge\frac1{n+1}$. Итак, $\alpha(\PP^n)=\frac1{n+1}$. При $n=1$: $\alpha(\PP^1)=\frac12$.
\end{example}

\begin{example}
$\alpha(\PP^1\times\PP^1)=\frac12$ (дивизор $2C$, $C$ --- диагональ или любая гладкая кривая
бистепени $(1,1)$, даёт $\lct=\frac12$; нижняя оценка --- как выше, с учётом того, что кривая
бистепени $(2m,2m)$ имеет кратность $\le2m$ в точках вне... требует небольшого разбора; см.
\cite{Che08}). Вообще, для торического многообразия Фано $\alpha(X)$ вычисляется по многограннику
\cite{Song}.
\end{example}

\subsection{Критерий Тяня}

\begin{theorem}[{Тянь \cite{Tia87}; алгебраически --- Одака--Сано \cite{OS12}, Фудзита \cite{Fuj19a}}]\label{thm:Tian-alpha}
Пусть $X$ --- $\Q$-Фано многообразие размерности $n$. Если $\alpha(X)>\frac n{n+1}$, то $X$
K-стабильно; если $\alpha(X)\ge\frac n{n+1}$, то $X$ K-полустабильно.
\end{theorem}

\begin{proof}[Идея (в терминах раздела~\ref{sec:beta})]
По валюативному критерию достаточно проверить $\beta(E)=A(E)-S(E)>0$ для всех $E$. Из определения
$\alpha$ следует $A(E)\ge\alpha\cdot\ord_E(D)$ для всех $D\sim_\Q-K_X$, в частности для
$D=\frac1mY$ с $Y\in|-mK_X|$ --- общим элементом, обращающимся в нуль вдоль $E$ с максимальным
порядком $\approx m\tau$, где $\tau$ --- псевдоэффективный порог: $A(E)\ge\alpha\,\tau(E)$. С другой
стороны, $S(E)\le\frac n{n+1}\tau(E)$ (функция $\vol(-K_X-tE)^{1/n}$ ---
вогнутая функция на $[0,\tau]$, и интеграл оценивается через значение в $0$; см.
\cite{Fuj19b}, \cite[\S\,7]{BJ}). Итак, $\beta(E)\ge(\alpha-\frac n{n+1})\tau(E)$.
\end{proof}

\begin{example}[Теорема ничего не говорит о $\PP^n$ при $n\ge2$]
$\alpha(\PP^1)=\frac12=\frac n{n+1}$, и $\PP^1$ K-полустабильно. Но $\alpha(\PP^n)=\frac1{n+1}<\frac n{n+1}$
при $n\ge2$, и критерий не применим; K-полустабильность $\PP^n$ будет доказана через
$G$-инвариантную версию (пример~\ref{ex:Pn-polystable}).
\end{example}

\begin{definition}
\emph{Степень} поверхности дель Пеццо $X$ --- $d=(-K_X)^2$. Для $\PP^2$ $d=9$, для $\PP^1\times\PP^1$
и $\mathbb F_1$ --- $8$, для $\Bl_r\PP^2$ --- $9-r$ \cite[7.7]{AG}.
\end{definition}

\begin{example}[Дель Пеццо степени 1; пример~3.17 обзора]\label{ex:dP1-alpha}
Пусть $X=\Bl_{p_1,\dots,p_8}\PP^2$ --- поверхность дель Пеццо степени 1 (точки в общем положении).
Тогда $-K_X=\pi^*(-K_{\PP^2})-\sum_{i=1}^8E_i=3L-\sum E_i$, и элементы $|-K_X|$ --- собственные
прообразы кубик, проходящих через $p_1,\dots,p_8$. Кубики через 8 общих точек образуют пучок
($h^0(-K_X)=2$ по Риману--Роху: $1+\frac12(-K)\cdot(-K-K)=1+1=2$), все его члены проходят ещё через
девятую точку $p_9$ (базисную), и через каждую точку $x\ne p_9$ проходит ровно один член $C_x$.
Все члены пучка неприводимы и приведены (три точки не на прямой, шесть не на конике), а
неприводимая плоская кубика --- гладкая, узловая или каспидальная.

Пусть $D\sim_\Q-K_X$, $D\ge0$. Оценим $\lct(X,D)$ в точке $x$.
\begin{itemize}
\item Если $C_x\not\subset\Supp D$ (или $x=p_9$ и некоторый член пучка не входит в $D$), то
$\mult_xD\le D\cdot C_x=(-K_X)^2=1$ (локальная кратность пересечения не меньше произведения
кратностей), и $(X,D)$ лог-канонична в $x$ (следствие~\ref{cor:lct-mult}).
\item Если $D=aC_x+D''$ с $a>0$, $C_x\not\subset\Supp D''$, то $D''\sim_\Q(1-a)(-K_X)$,
$\mult_xD''\le D''\cdot C_x=1-a$, и $\mult_xD\le a\,\mult_xC_x+1-a$. В гладкой точке $C_x$ это
$\le1$; в особой точке (узел или касп) нужен более тонкий анализ, который даёт: $(X,cD)$ lc при
$c\le\lct(X,C_x)$, а $\lct(X,C_x)=1$ для гладкой и узловой $C_x$ и $\frac56$ для каспидальной
(пример~\ref{ex:lct-curves}).
\end{itemize}
Итог \cite[Thm~1.7]{Che08}: $\alpha(X)=\frac56$, если в пучке $|-K_X|$ есть каспидальная кривая, и
$\alpha(X)=1$ иначе. В любом случае $\alpha(X)\ge\frac56>\frac23$, и по теореме~\ref{thm:Tian-alpha}
\emph{поверхность дель Пеццо степени 1 K-стабильна}.
\end{example}

\begin{theorem}[{Чельцов, \cite{Che08}}]\label{thm:Cheltsov}
$\alpha$-инварианты гладких поверхностей дель Пеццо:
\[
\begin{array}{c|c|c|c|c|c|c|c|c|c}
X & \PP^2 & \PP^1\times\PP^1 & \mathbb F_1 & d=7 & d=6 & d=5 & d=4 & d=3 & d=2,\ 1\\\hline
\alpha(X) & \frac13 & \frac12 & \frac13 & \frac13 & \frac12 & \frac12 & \frac23 &
\frac23\ \text{или}\ \frac34 & \frac34\ \text{или}\ \frac56;\ \ \frac56\ \text{или}\ 1
\end{array}
\]
При $d=3$ значение $\frac23$ отвечает кубикам с точкой Эккардта (пример~\ref{ex:lct-cubic-sections}),
при $d=2$ --- поверхностям с тэкнодальной кривой в $|-K|$, при $d=1$ --- с каспидальной.
\end{theorem}

\begin{theorem}[{Фудзита, \cite{Fuj19a}}]\label{thm:Fujita-alpha}
Если $X$ --- поверхность дель Пеццо или гладкое многообразие Фано размерности $\ge3$ и
$\alpha(X)\ge\frac n{n+1}$, то $X$ K-стабильно.
\end{theorem}

\begin{corollary}\label{cor:dP-le4}
Все гладкие поверхности дель Пеццо степени $\le4$ K-стабильны. В частности, \emph{все гладкие
кубические поверхности K-стабильны}.
\end{corollary}

\begin{example}[Гиперповерхности]
Если $X\subset\PP^{n+1}$ --- гладкая гиперповерхность степени $n+1$ (Фано индекса 1), то
$\alpha(X)\ge\frac n{n+1}$ (Чельцов--Пак \cite{CP02}), и по теореме~\ref{thm:Fujita-alpha} $X$
K-стабильна. При $n=2$ это кубические поверхности (снова), при $n=3$ --- квартики в $\PP^4$.
\end{example}

\subsection{Эквивариантный $\alpha$-инвариант}

\begin{theorem}[{\cite{DS16}, \cite{LX20}, \cite{LZ22}, \cite{Zhu21}}]\label{thm:alphaG}
Пусть $X$ --- $\Q$-Фано многообразие с действием алгебраической группы $G$.
\begin{enumerate}
\item Если $\alpha_G(X)\ge\frac n{n+1}$, то $X$ K-полустабильно.
\item Если $G$ редуктивна и $\alpha_G(X)>\frac n{n+1}$, то $X$ K-полистабильно.
\end{enumerate}
\end{theorem}

Смысл: $G$-инвариантных дивизоров гораздо меньше, чем всех, так что $\alpha_G\ge\alpha$ и
$\alpha_G$ бывает велик даже при маленьком $\alpha$. Это отражает общий принцип: K-стабильность
достаточно проверять на $G$-эквивариантных тест-конфигурациях (Датар--Секейхиди, Чжуан).

\begin{example}[$\PP^n$ K-полистабильно]\label{ex:Pn-polystable}
$G=\PGL_{n+1}$ действует на $\PP^n$ транзитивно, и $G$-инвариантных эффективных дивизоров нет
(носитель инвариантного дивизора --- собственное замкнутое $G$-инвариантное подмножество, а
таких нет). Значит, $\alpha_G(\PP^n)=+\infty>\frac n{n+1}$, $G$ редуктивна, и $\PP^n$ K-полистабильно.
Оно не K-стабильно (бесконечная группа автоморфизмов, замечание~\ref{rem:BX}).
\end{example}

\begin{example}[$\PP^1\times\PP^1$ и $\Bl_3\PP^2$; упражнение~3.4(4)]\label{ex:dP8-dP6-polystable}
\emph{$X=\PP^1\times\PP^1$}: $G=\Aut(X)\supset\PGL_2\times\PGL_2$ действует транзитивно, инвариантных
дивизоров нет, $\alpha_G=\infty$, $X$ K-полистабильно.

\emph{$X=\Bl_3\PP^2$} (дель Пеццо степени 6): $G=\Aut(X)=T\rtimes(S_3\times\Z/2)$, где
$T=(\C^*)^2$ --- тор (поверхность торическая, \cite[4.11(4)]{AG}), а конечная группа переставляет
шесть $(-1)$-кривых $D_1,\dots,D_6$ (шестиугольник). $G$ редуктивна. $T$-инвариантный эффективный
дивизор сосредоточен на дополнении к открытой орбите, т.\,е. на $\bigcup D_i$: $D=\sum a_iD_i$;
группа $S_3\times\Z/2$ действует на $\{D_i\}$ транзитивно, так что $G$-инвариантный $D$ имеет вид
$a\sum D_i$, а $\sum D_i\sim-K_X$ (торическая формула \cite[4.20]{AG}) даёт $a=1$. Пара
$(X,\sum D_i)$ --- торическая пара с полной границей, она лог-каноническая (snc-дивизор с
коэффициентами 1), и $\lct=1$. Итак, $\alpha_G(X)=1>\frac23$, и $\Bl_3\PP^2$ K-полистабильно.
(Сравните: $\alpha(X)=\frac12<\frac23$ --- без эквивариантности критерий Тяня не работает.)
\end{example}

\begin{example}[Кубика $3A_2$ K-полистабильна]\label{ex:3A2-polystable}
Пусть $X=\{x_0x_1x_2=x_3^3\}$ (пример~\ref{ex:cubic-normal-forms}(2)). $X$ --- $\Q$-Фано: нормальная
горенштейнова кубика с особенностями $A_2$ (klt), $-K_X=\cO_X(1)$ обилен. Группа
$G=\Aut(X)\supset T\rtimes S_3$, где $T=\{(x_0,x_1,x_2,x_3)\mapsto(ax_0,bx_1,cx_2,x_3):abc=1\}\cong(\C^*)^2$,
редуктивна. Открытая $T$-орбита --- $\{x_0x_1x_2\ne0\}$, её дополнение --- три прямые
$D_i=\{x_i=x_3=0\}$. Как в предыдущем примере, $G$-инвариантный эффективный $D\sim_\Q-K_X$ имеет вид
$a(D_1+D_2+D_3)$; плоское сечение $\{x_3=0\}\cap X=\{x_0x_1x_2=0\}=D_1+D_2+D_3$, так что
$D_1+D_2+D_3\sim-K_X$ и $a=1$. Пара $(X,D_1+D_2+D_3)$ --- торическая пара с полной границей,
т.\,е. лог-каноническая (это общий факт для торических пар $(X_\Sigma,\sum D_\rho)$: на торическом
разрешении $K+\sum D=\pi^*(K+\sum D)$ и все дискрепантности равны $-1$). Значит, $\lct=1$,
$\alpha_G(X)=1>\frac23$, и по теореме~\ref{thm:alphaG}(2) \emph{кубика $3A_2$ K-полистабильна}.
Она не K-стабильна: $\Aut$ бесконечна. Это в точности то, что предсказывает ГИТ
(теорема~\ref{thm:Hilbert}(3)).
\end{example}

\begin{example}[Кубика Ферма]
$X=\{x_0^3+x_1^3+x_2^3+x_3^3=0\}$ имеет 18 точек Эккардта (например, $[1:-1:0:0]$: плоскость
$x_0+x_1=0$ высекает $x_2^3+x_3^3=0$ --- три прямые через точку), поэтому $\alpha(X)=\frac23$
(теорема~\ref{thm:Cheltsov}), и по теореме~\ref{thm:Fujita-alpha} $X$ K-стабильна. Группа
автоморфизмов $(\mu_3)^3\rtimes S_4$ конечна, так что K-стабильность и K-полистабильность
совпадают. Это один из способов решить упражнение~5.1(3) обзора: раз одна кубика K-(полу)стабильна,
то по открытости (теорема~\ref{thm:openness}) и общая кубика K-(полу)стабильна; другой способ ---
пример~\ref{ex:3A2-polystable}.
\end{example}

% ============================================================
\section{$\beta$- и $\delta$-инварианты. Валюативный критерий}\label{sec:beta}
% ============================================================

Это \S3.3 обзора. $\alpha$-инвариант даёт лишь достаточные условия. Здесь --- критерий,
являющийся «тогда и только тогда», и притом сформулированный в терминах бирациональной геометрии
$X$: K-стабильность $\iff$ неравенство для \emph{каждого} дивизора над $X$. Он позволяет доказывать
\emph{не}стабильность, предъявляя один дивизор.

\subsection{Определения}

\begin{definition}[{Фудзита \cite{Fuj19b}, Ли \cite{Li17}}]\label{def:beta}
Пусть $X$ --- $\Q$-Фано многообразие размерности $n$, $E$ --- простой дивизор над $X$,
реализованный на $\mu\colon Y\to X$.
\begin{itemize}
\item \emph{Лог-дискрепантность} $A_X(E)=1+\ord_E(K_Y-\mu^*K_X)=1+a_X(E)$ (определение~\ref{def:discrep}).
\item \emph{Ожидаемый порядок обращения в нуль}
\[
S_X(E)=\frac1{(-K_X)^n}\int_0^\infty\vol(\mu^*(-K_X)-tE)\,dt=\frac1{(-K_X)^n}\int_0^{\tau(E)}\vol(-K_X-tE)\,dt ,
\]
где $\tau(E)$ --- псевдоэффективный порог (\S\ref{ss:volume}); результат не зависит от выбора
$Y$, поэтому пишут $\vol(-K_X-tE)$.
\item \emph{$\beta$-инвариант} $\beta_X(E)=A_X(E)-S_X(E)$ и \emph{$\delta$-инвариант}
$\delta_X(E)=A_X(E)/S_X(E)$.
\item \emph{$\delta$-инвариант многообразия} (Фудзита--Одака, Блюм--Йонссон)
$\delta(X)=\inf_E\delta_X(E)$ по всем простым дивизорам над $X$.
\end{itemize}
Для пары $(X,\Delta)$ всё то же с $-K_X-\Delta$ вместо $-K_X$ и $A_{(X,\Delta)}$ вместо $A_X$:
$\beta_{(X,\Delta)}(E)=A_{(X,\Delta)}(E)-S_{(X,\Delta)}(E)$.
\end{definition}

Почему $S(E)$ --- «ожидаемый порядок обращения в нуль»? По определению объёма
\[
S_X(E)=\lim_{m\to\infty}\frac1{m\,h^0(-mK_X)}\sum_{j\ge1}\dim\{s\in H^0(-mK_X):\ \ord_E(s)\ge j\}
=\lim_m\frac{1}{m}\cdot\frac{\sum_{s\in\text{базис}}\ord_E(s)}{h^0(-mK_X)}
\]
(для базиса, согласованного с фильтрацией по $\ord_E$) --- среднее значение $\frac1m\ord_E$ по
сечениям $-mK_X$. Сравните с инвариантом Футаки: $w_k/(kd_k)$ --- среднее значение веса. Это не
случайно: специальная тест-конфигурация с центральным слоем $X_0$ задаёт нормирование
$v=\ord_{X_0}|_{\C(X)}$, и $\Fut(\cX)=c\cdot\beta_X(v)$, $c>0$ \cite{Fuj19b, Li17}; в этом смысле
$\beta$ --- «инвариант Футаки, записанный на языке нормирований».

\begin{remark}
\begin{enumerate}
\item $\beta$ и $\delta$ определены для \emph{всех} дивизоров над $X$, а не только для
дивизоров на $X$: исключительные дивизоры раздутий, взвешенных раздутий и т.\,д.
\item $S_X(E)\le\frac n{n+1}\tau(E)$ и $S_X(E)\ge\frac1{n+1}\tau(E)$ (вогнутость $\vol^{1/n}$).
\item Если $\delta(X)<\frac{n+1}n$, то инфимум достигается: существует дивизор $E$ над $X$ с
$\delta_X(E)=\delta(X)$ \cite{LXZ}. Связь с $\alpha$: $\frac{n+1}n\alpha(X)\le\delta(X)\le(n+1)\alpha(X)$ \cite{BJ}.
\end{enumerate}
\end{remark}

\subsection{Валюативный критерий}

\begin{theorem}[{Фудзита \cite{Fuj19b}, Ли \cite{Li17}, Фудзита--Одака \cite{FO18}, Блюм--Йонссон \cite{BJ}, Ли--Сюй--Чжуан \cite{LXZ}}]\label{thm:valuative}
Пусть $X$ --- $\Q$-Фано многообразие (или $(X,\Delta)$ --- лог-Фано пара). Тогда
\begin{enumerate}
\item $X$ K-полустабильно $\iff$ $\beta_X(E)\ge0$ для всех простых дивизоров $E$ над $X$
$\iff$ $\delta(X)\ge1$;
\item $X$ K-стабильно $\iff$ $\beta_X(E)>0$ для всех $E$ $\iff$ $\delta(X)>1$.
\end{enumerate}
\end{theorem}

История: Фудзита и Ли доказали «$\Rightarrow$» и «$\Leftarrow$» для равномерной K-стабильности
($\delta>1$); равносильность равномерной K-стабильности и K-стабильности для $\Q$-Фано ---
теорема Ли--Сюя--Чжуана о конечной порождённости \cite{LXZ}. Для K-полистабильности критерий в
этих терминах формулируется через $\beta$ и $G$-эквивариантность (\cite{Zhu21}): $X$
K-полистабильно $\iff$ $\beta(E)\ge0$ для всех $E$ и $\beta(E)=0$ лишь для $E$, индуцированных
$\Gm$-действиями (произведениями).

\begin{proof}[Идея «$\Rightarrow$» в (1)]
По дивизору $E$ строится последовательность тест-конфигураций: фильтрация $H^0(-mK_X)$ по
$\ord_E$ задаёт алгебру Риса $\bigoplus_{m,j}t^{-j}\{s:\ord_E(s)\ge j\}$; если она конечно
порождена, её $\Proj$ над $\A^1$ --- тест-конфигурация с $\Fut=c\,\beta(E)$, а в общем случае
аппроксимируем. Поэтому $\beta(E)<0$ даёт дестабилизирующую (почти) тест-конфигурацию.
«$\Leftarrow$»: всякая специальная тест-конфигурация происходит из одного дивизора $E$ (центральный
слой $X_0$ определяет нормирование на $\C(X)$), и $\Fut=c\,\beta(E)$.
\end{proof}

\textbf{Красный флаг (меньший).} Критерий всё ещё требует перебрать \emph{все} дивизоры над $X$.
Но теперь (а) доказать \emph{не}стабильность можно, предъявив один дивизор с $\beta<0$; (б) есть
методы (раздел~\ref{sec:az}) оценивать $\delta$ снизу через индукцию по размерности.

\subsection{Как считать $S(E)$: примеры на поверхностях}

Схема вычисления $S(E)$ на поверхности: (1) найти $A(E)$ (для кривой $E\subset X$: $A=1$; для
исключительной кривой раздутия гладкой точки: $A=2$; для $(-n)$-кривой над $\frac1n(1,1)$:
$A=\frac2n$); (2) найти, при каких $t$ дивизор $-K_X-tE$ nef (пересечения с образующими конуса
Мори), там $\vol=(-K_X-tE)^2$; (3) за nef-порогом найти разложение Зарисского
(теорема~\ref{thm:Zariski}) и $\vol=P_t^2$ вплоть до псевдоэффективного порога $\tau$; (4) взять
интеграл.

\begin{example}[Пример~3.30 обзора: раздутие точки на $\PP^2$]\label{ex:beta-P2-point}
$X=\PP^2$, $\mu\colon Y=\Bl_p\PP^2\to\PP^2$, $E$ --- исключительная кривая. $A_{\PP^2}(E)=2$
($K_Y=\mu^*K_{\PP^2}+E$). Далее, $\mu^*(-K_{\PP^2})-tE=3L-tE=-K_Y+(1-t)E$. Конус Мори $Y$ порождён
$E$ и слоем $F=L-E$ (пример~\ref{ex:cones}): $(3L-tE)\cdot E=t\ge0$, $(3L-tE)\cdot F=3-t$. Итак,
$3L-tE$ nef при $0\le t\le3$, а при $t>3$ не псевдоэффективен (отрицателен на nef-классе $F$).
Поэтому $\vol(3L-tE)=9-t^2$ на $[0,3]$ и
\[
S_{\PP^2}(E)=\frac19\int_0^3(9-t^2)\,dt=\frac19\cdot18=2,\qquad\beta_{\PP^2}(E)=2-2=0 .
\]
Само по себе это не доказывает K-полустабильности $\PP^2$ (нужны все $E$), но согласуется с ней
и показывает, что $\PP^2$ не K-стабильно: $\beta=0$ для нетривиального $E$, отвечающего
произведению с действием $\mathrm{diag}(t,t,1)$ (его нормирование на $\C(\PP^2)$ --- порядок
обращения в нуль в точке $p=[0:0:1]$, т.\,е. $\ord_E$).
\end{example}

\begin{example}[Прямая на $\PP^2$, слои и диагональ на $\PP^1\times\PP^1$]
\begin{enumerate}
\item $X=\PP^2$, $E=$ прямая: $A=1$, $\vol(3L-tL)=(3-t)^2$ на $[0,3]$, $S=\frac19\int_0^3(3-t)^2dt=1$,
$\beta=0$.
\item $X=\PP^1\times\PP^1$, $E=F$ --- слой (класс $(1,0)$): $A=1$, $-K-tF=(2-t,2)$ nef при
$t\le2$, $\vol=2\cdot(2-t)\cdot2=4(2-t)$, $S=\frac18\int_0^24(2-t)dt=1$, $\beta=0$.
\item $X=\PP^1\times\PP^1$, $E=\Delta$ --- диагональ (класс $(1,1)$): $-K-t\Delta=(2-t,2-t)$,
$\vol=2(2-t)^2$ на $[0,2]$, $S=\frac18\int_0^22(2-t)^2dt=\frac18\cdot\frac{16}3=\frac23$,
$\beta=\frac13>0$.
\end{enumerate}
Нулевые $\beta$ в (1), (2) отвечают $\Gm$-действиям (произведениям); положительное $\beta$ в (3)
--- тому, что диагональ не отвечает никакому произведению.
\end{example}

\begin{example}[$\mathbb F_1$ K-нестабильна; упражнения~3.4(5,6)]\label{ex:beta-F1}
$X=\mathbb F_1=\Bl_p\PP^2$, $E\subset X$ --- исключительная кривая, $A_X(E)=1$. По
примеру~\ref{ex:vol-F1} $\vol(-K_X-tE)=9-(1+t)^2$ на $[0,2]$ и $0$ далее; $(-K_X)^2=8$. Поэтому
\[
S_X(E)=\frac18\int_0^2\big(9-(1+t)^2\big)dt=\frac18\Big(18-\frac{27-1}3\Big)=\frac18\cdot\frac{28}3=\frac76,
\qquad\beta_X(E)=1-\frac76=-\frac16<0 .
\]
По валюативному критерию \emph{$\mathbb F_1$ K-нестабильна}. Это согласуется с теоремой Мацусимы
(пример~\ref{ex:Aut-F1}) и с торическим критерием (пример~\ref{ex:toric-beta}).
\end{example}

\begin{example}[Дель Пеццо степени 7 K-нестабильна; упражнение~3.4(7)]\label{ex:beta-dP7}
$X=\Bl_{p_1,p_2}\PP^2$, $E=E_1$, $A=1$, $(-K)^2=7$. По примеру~\ref{ex:vol-dP7}
\[
S_X(E_1)=\frac17\Big(\int_0^1\big(8-(1+t)^2\big)dt+\int_1^2(8-4t)\,dt\Big)
=\frac17\Big(8-\frac73+2\Big)=\frac17\cdot\frac{23}3=\frac{23}{21},\qquad
\beta_X(E_1)=-\frac2{21}<0 .
\]
Итак, дель Пеццо степени 7 K-нестабильна.
\end{example}

\begin{example}[Дель Пеццо степени 6: $\beta(E_1)=0$]
$X=\Bl_3\PP^2$, $-K-tE_1=3L-(1+t)E_1-E_2-E_3$. Образующие конуса: $E_i$, $\ell_{ij}=L-E_i-E_j$.
Пересечения: с $E_1$: $1+t$; с $E_2,E_3$: $1$; с $\ell_{12},\ell_{13}$: $1-t$; с $\ell_{23}$: $1$.
Nef при $t\le1$: $\vol=9-(1+t)^2-2=7-(1+t)^2$. При $1\le t\le2$ отрицательная часть
$N=(t-1)(\ell_{12}+\ell_{13})$ (две непересекающиеся $(-1)$-кривые), и
$P_t=(5-2t)L-(3-t)E_1-(2-t)E_2-(2-t)E_3$, $P_t^2=(2-t)(4-t)$; при $t=2$ объём обращается в нуль.
\[
S_X(E_1)=\frac16\Big(\int_0^1\big(7-(1+t)^2\big)dt+\int_1^2(2-t)(4-t)\,dt\Big)=\frac16\Big(\frac{14}3+\frac43\Big)=1,
\qquad\beta_X(E_1)=0 .
\]
Нуль --- потому что $E_1$ инвариантна относительно тора, и ей отвечает произведение; в
примере~\ref{ex:dP8-dP6-polystable} мы видели, что $\Bl_3\PP^2$ K-полистабильна.
\end{example}

\begin{example}[Торическая формула для $\beta$]\label{ex:toric-beta}
Пусть $X=X_P$ --- торическое $\Q$-Фано многообразие с антиканоническим многогранником
$P=\{m:\langle m,u_\rho\rangle\ge-1\ \forall\rho\}$ и $D_\rho$ --- инвариантный дивизор луча $\rho$.
Тогда $A(D_\rho)=1$, а $H^0(-mK_X-jD_\rho)$ натянуто на характеры $\chi^{m'}$,
$m'\in mP\cap\{\langle\cdot,u_\rho\rangle\ge-m+j\}$ (теорема~4.16 из \cite{AG}), так что
$\vol(-K_X-tD_\rho)=n!\vol(P\cap\{\langle x,u_\rho\rangle\ge-1+t\})$ и
\[
S(D_\rho)=\frac{1}{n!\vol P}\int_0^\infty n!\vol\big(P\cap\{\langle x,u_\rho\rangle\ge t-1\}\big)dt
=\frac1{\vol P}\int_P\big(\langle x,u_\rho\rangle+1\big)dx=1+\langle b_P,u_\rho\rangle ,
\]
где $b_P$ --- барицентр (теорема Фубини: $\int_0^\infty\vol\{x\in P:\langle x,u\rangle+1\ge t\}dt=\int_P(\langle x,u\rangle+1)dx$). Итак,
\[
\boxed{\ \beta(D_\rho)=-\langle b_P,u_\rho\rangle\ }.
\]
Это согласуется с формулой для инварианта Футаки произведений из примера~\ref{ex:toric-product-tc}
(с точностью до множителя $\frac12$). Для пяти гладких торических поверхностей дель Пеццо
(веера --- \cite[рис.~2]{AG}; многоугольники $P=\{\langle m,u_i\rangle\ge-1\}$):
\[
\begin{array}{l|l|c|c}
X & \text{лучи }u_i & \text{вершины }P & b_P\\\hline
\PP^2 & (1,0),(0,1),(-1,-1) & (-1,-1),(2,-1),(-1,2) & (0,0)\\
\PP^1\times\PP^1 & (1,0),(0,1),(-1,0),(0,-1) & (\pm1,\pm1) & (0,0)\\
\mathbb F_1 & (1,0),(0,1),(-1,1),(0,-1) & (-1,-1),(0,-1),(2,1),(-1,1) & (\tfrac1{12},\tfrac16)\\
\Bl_2\PP^2 & (1,0),(1,1),(0,1),(-1,0),(0,-1) & (-1,0),(0,-1),(1,-1),(1,1),(-1,1) & (\tfrac2{21},\tfrac2{21})\\
\Bl_3\PP^2 & (1,0),(1,1),(0,1),(-1,0),(-1,-1),(0,-1) & \text{шестиугольник} & (0,0)
\end{array}
\]
Для $\mathbb F_1$ кривая $E$ отвечает лучу $u=(0,1)$ ($u_1+u_3=u_2$, т.\,е. $D_2^2=-1$), и
$\beta(E)=-\langle b_P,u\rangle=-\frac16$ --- ровно то, что мы получили в примере~\ref{ex:beta-F1}
прямым вычислением. Для $\Bl_2\PP^2$ кривая $E_1$ отвечает лучу $(1,0)$, и $\beta(E_1)=-\frac2{21}$
--- как в примере~\ref{ex:beta-dP7}. Это хорошая проверка обоих вычислений. Общий вывод:
\emph{торическое многообразие Фано K-полустабильно тогда и только тогда, когда $b_P=0$}
(в этом случае оно даже K-полистабильно), и барицентр --- полный ответ для торических
многообразий.
\end{example}

\begin{example}[$\PP(1,1,n)$ K-нестабильно при $n>1$; упражнения~3.4(8), 6.3(10d)]\label{ex:beta-P11n}
Пусть $X=\PP(1,1,n)$, $\pi\colon Y=\mathbb F_n\to X$ --- раздутие вершины, $E$ --- $(-n)$-кривая,
$F$ --- слой. По примеру~\ref{ex:discrep-P11n} $A_X(E)=\frac2n$. Далее, $-K_X=(n+2)\ell$,
$\pi^*\ell=F+\frac1nE$ (пример~\ref{ex:P11n-intersections}), так что
\[
\pi^*(-K_X)-tE=\frac{n+2}n(E+nF)-tE,\qquad (E+nF)\cdot E=0,\ (E+nF)\cdot F=1,\ (E+nF)^2=n .
\]
Пересечение с $E$: $-tE^2=tn\ge0$; с $F$: $\frac{n+2}n-t$. Дивизор nef при $t\le\frac{n+2}n=:\tau$
и не псевдоэффективен далее ($F$ nef). Объём: $\big(\frac{n+2}n\big)^2n-nt^2=\frac{(n+2)^2}n-nt^2$.
\[
S_X(E)=\frac n{(n+2)^2}\int_0^{(n+2)/n}\Big(\frac{(n+2)^2}n-nt^2\Big)dt
=\frac n{(n+2)^2}\cdot\frac{(n+2)^3}{n^2}\cdot\frac23=\frac{2(n+2)}{3n},
\]
\[
\beta_X(E)=\frac2n-\frac{2(n+2)}{3n}=\frac{6-2n-4}{3n}=-\frac{2(n-1)}{3n}.
\]
При $n=1$ ($X=\PP^2$, $E$ --- исключительная кривая точки) получаем $\beta=0$ --- пример~\ref{ex:beta-P2-point}.
При $n\ge2$: $\beta<0$, и \emph{$\PP(1,1,n)$ K-нестабильно}; в частности, квадратичный конус
$\PP(1,1,2)$ K-нестабилен ($\beta=-\frac13$). Торически: барицентр многоугольника $\PP(1,1,n)$ при
$n\ge2$ не в нуле.
\end{example}

\begin{example}[Раздутие $\PP^3$ вдоль плоской кубики; упражнение~3.4(9)]\label{ex:Bl-cubic-P3}
Пусть $C\subset H\cong\PP^2\subset\PP^3$ --- гладкая плоская кубика, $\pi\colon X=\Bl_C\PP^3\to\PP^3$,
$E$ --- исключительный дивизор, $\tilde H$ --- собственный прообраз плоскости. По
теореме~\ref{thm:blowup-K} $-K_X=4\pi^*H-E$, а $\tilde H=\pi^*H-E$, так что $-K_X=3\pi^*H+\tilde H$.
\emph{$X$ --- Фано}: конус Мори $X$ порождён классом $f$ слоя $E\to C$ и классом $\tilde\ell$
собственного прообраза прямой $\ell\subset H$ ($\pi^*H\cdot\tilde\ell=1$, $E\cdot\tilde\ell=3$);
$-K_X\cdot f=1$, $-K_X\cdot\tilde\ell=4-3=1$; обильность по критерию Клеймана. Индексы пересечения:
$(\pi^*H)^3=1$, $(\pi^*H)^2E=0$, $\pi^*H\cdot E^2=-\deg C=-3$, $E^3=-\deg N_{C/\PP^3}=-(3+9)=-12$
($N_{C/\PP^3}=\cO_C(1)\oplus\cO_C(3)$); отсюда $(-K_X)^3=64-9\cdot4\cdot1+12=40$ (семейство
№2.28 в списке Мори--Мукаи).

Вычислим $\beta(\tilde H)$: $A=1$ ($\tilde H$ --- дивизор на $X$). $D_t=-K_X-t\tilde H=(4-t)\pi^*H-(1-t)E$.
При $0\le t\le1$: $D_t\cdot f=1-t\ge0$, $D_t\cdot\tilde\ell=(4-t)-3(1-t)=1+2t>0$ --- nef, и
\[
\vol(D_t)=D_t^3=(4-t)^3-9(4-t)(1-t)^2+12(1-t)^3 .
\]
При $1\le t\le4$: $D_t=(4-t)\pi^*H+(t-1)E$ с $t-1\ge0$, и $\pi_*\cO_X(kE)=\cO_{\PP^3}$ при $k\ge0$,
так что $H^0(X,mD_t)=H^0(\PP^3,\cO(m(4-t)))$ и $\vol(D_t)=(4-t)^3$; при $t>4$ объём нулевой.
Интегрируем:
\[
\int_0^1\!\big((4-t)^3-9(4-t)(1-t)^2+12(1-t)^3\big)dt=\frac{175}4-\frac{45}4+3=\frac{71}2,\qquad
\int_1^4(4-t)^3dt=\frac{81}4,
\]
\[
S(\tilde H)=\frac1{40}\Big(\frac{71}2+\frac{81}4\Big)=\frac{223}{160},\qquad
\beta(\tilde H)=1-\frac{223}{160}=-\frac{63}{160}<0 .
\]
Итак, $\Bl_C\PP^3$ K-нестабильно. (Подсказка обзора верна: $\beta(E)$ здесь положителен и ничего
не даёт.) Это пример в размерности~3, где разложения Зарисского нет, но объём удаётся найти
прямо --- через структуру базисного множества.
\end{example}

\begin{example}[K-стабильность пар: $(\PP(1,1,2),cQ)$; упражнение~3.4(10)]\label{ex:P112-pair}
Пусть $X=\PP(1,1,2)$, $Q=\{z=0\}$ --- коника «на бесконечности» (пример~\ref{ex:P112}), не
проходящая через вершину, $Q\sim\cO(2)$, $-K_X=\cO(4)$, $-K_X-cQ=\cO(4-2c)$ (обилен при $0<c<2$).
Вычислим $\beta$ для двух дивизоров.

\emph{$E$ --- исключительная кривая разрешения $Y=\mathbb F_2\to X$.} $A_{(X,cQ)}(E)=A_X(E)=\frac22=1$
($Q$ не проходит через вершину). $\pi^*(-K_X-cQ)=\frac{4-2c}2(E+2F)=(2-c)(E+2F)$;
$\pi^*(-K_X-cQ)-tE$ nef при $t\le2-c$, объём $2(2-c)^2-2t^2$, $\vol(-K_X-cQ)=\frac{(4-2c)^2}2=2(2-c)^2$:
\[
S(E)=\frac1{2(2-c)^2}\int_0^{2-c}\big(2(2-c)^2-2t^2\big)dt=\frac23(2-c),\qquad
\beta_{(X,cQ)}(E)=1-\frac{2(2-c)}3=\frac{2c-1}3 .
\]
\emph{$Q$ сама.} $A_{(X,cQ)}(Q)=1-c$; $-K_X-cQ-tQ=\cO(4-2c-2t)$, объём $2(2-c-t)^2$ на $[0,2-c]$:
\[
S(Q)=\frac1{2(2-c)^2}\int_0^{2-c}2(2-c-t)^2dt=\frac{2-c}3,\qquad
\beta_{(X,cQ)}(Q)=1-c-\frac{2-c}3=\frac{1-2c}3 .
\]
Итак, $\beta(E)<0$ при $c<\frac12$ и $\beta(Q)<0$ при $c>\frac12$: \emph{пара $(\PP(1,1,2),cQ)$
K-нестабильна при всех $c\ne\frac12$}. При $c=\frac12$ оба инварианта равны нулю; в
примере~\ref{ex:AZ-P112} мы докажем, что $(\PP(1,1,2),\frac12Q)$ K-полустабильна. Геометрический
смысл: $(\PP(1,1,2),\frac12Q)$ --- фактор $(\PP^1\times\PP^1,0)$ по инволюции $(x,y)\mapsto(y,x)$
(с ветвлением вдоль диагонали, дающим коэффициент $\frac12$ у образа диагонали $Q$), а
$\PP^1\times\PP^1$ K-полистабильно.
\end{example}

\subsection{Сводка: K-стабильность гладких поверхностей дель Пеццо}

\begin{theorem}[Тянь, 1990; алгебраически --- по результатам этого и предыдущего разделов]\label{thm:dP-table}
\[
\begin{array}{c|c|l}
\text{степень} & \text{K-стабильность} & \text{причина}\\\hline
9\ (\PP^2) & \text{полистабильна} & \alpha_G\ (\text{пример~\ref{ex:Pn-polystable}})\\
8\ (\PP^1\times\PP^1) & \text{полистабильна} & \alpha_G\ (\text{пример~\ref{ex:dP8-dP6-polystable}})\\
8\ (\mathbb F_1) & \text{нестабильна} & \beta(E)=-\tfrac16\ (\text{пример~\ref{ex:beta-F1}})\\
7 & \text{нестабильна} & \beta(E_1)=-\tfrac2{21}\ (\text{пример~\ref{ex:beta-dP7}})\\
6 & \text{полистабильна} & \alpha_G\ (\text{пример~\ref{ex:dP8-dP6-polystable}})\\
5 & \text{стабильна} & \alpha_G>\tfrac23\ \text{для }G=\Aut(X)\cong S_5\ \text{и конечность }\Aut\ \cite{Che08}\\
\le4 & \text{стабильна} & \alpha\ge\tfrac23\ (\text{следствие~\ref{cor:dP-le4}})
\end{array}
\]
В частности, гладкая поверхность дель Пеццо K-полустабильна $\iff$ K-полистабильна $\iff$ её группа
автоморфизмов редуктивна $\iff$ она допускает КЭ-метрику (Тянь).
\end{theorem}

% ============================================================
\section{Теория допустимых флагов Аббана--Чжуана}\label{sec:az}
% ============================================================

Это \S4 обзора. Валюативный критерий хорош для доказательства \emph{не}стабильности; чтобы доказать
стабильность, нужно оценить $\delta(X)$ снизу, т.\,е. справиться со \emph{всеми} дивизорами над $X$
сразу. Метод Аббана--Чжуана \cite{AZ} --- «формула присоединения для $\delta$-инварианта»: оценка
$\delta$ многообразия через $\delta$-инварианты, вычисленные на дивизоре $E\subset X$ (на единицу
меньшей размерности). Индукция по размерности сводит всё к вычислениям на кривых.

\subsection{Градуированные линейные системы и их фильтрации}

Пусть $(X,\Delta)$ --- лог-Фано пара, $L=-K_X-\Delta$ (можно считать $\Delta=0$).

\begin{definition}\label{def:graded-series}
\emph{Полная градуированная линейная система} $V_\bullet=\{V_m\}_{m\in\mathbb N}$, $V_m=H^0(X,mL)$;
её объём $\vol(V_\bullet)=\lim\dim V_m/(m^n/n!)=\vol(L)$. Для дивизора $E$ над $X$ и $t\ge0$ положим
\[
\mathcal F_E^tV_m=\{s\in V_m:\ \ord_E(s)\ge mt\},\qquad
\vol(\mathcal F_EV_\bullet)^t=\lim_{m\to\infty}\frac{\dim\mathcal F_E^tV_m}{m^n/n!},\qquad
S(V_\bullet;E)=\frac1{\vol V_\bullet}\int_0^\infty\vol(\mathcal F_EV_\bullet)^t\,dt .
\]
\end{definition}

Так как $\mathcal F_E^tV_m=H^0(Y,m\mu^*L-\lceil mt\rceil E)$ (сечения $mL$, обращающиеся в нуль
вдоль $E$ с порядком $\ge mt$), имеем $\vol(\mathcal F_EV_\bullet)^t=\vol(\mu^*L-tE)$ и
$S(V_\bullet;E)=S_{(X,\Delta)}(E)$ --- это определение~\ref{def:beta}, записанное на языке фильтраций.
Преимущество языка в том, что он переносится на \emph{неполные} и \emph{многоградуированные} системы.

\begin{definition}
$\delta(X,\Delta;V_\bullet)=\inf_E\frac{A_{(X,\Delta)}(E)}{S(V_\bullet;E)}$ по всем $E$ над $X$;
для подмногообразия $Z\subset X$
\[
\delta_Z(X,\Delta;V_\bullet)=\inf_{E:\,Z\subset c_X(E)}\frac{A_{(X,\Delta)}(E)}{S(V_\bullet;E)}
\]
--- инфимум по дивизорам, центр которых содержит $Z$.
\end{definition}

Очевидно, $\delta(X,\Delta;V_\bullet)=\inf_Z\delta_Z=\inf_{p\in X}\delta_p$ (достаточно точек: центр
любого $E$ содержит точку). По теореме~\ref{thm:valuative} $(X,\Delta)$ K-полустабильна, если
$\delta_p(X,\Delta;V_\bullet)\ge1$ для всех $p\in X$.

\subsection{Ограничение на дивизор}

Пусть $E$ --- простой дивизор над $X$ \emph{plt-типа}: существует $\pi\colon\tilde X\to X$,
извлекающий ровно $E$ (т.\,е. $\mathrm{Exc}(\pi)=E$), причём $-E$ $\pi$-обилен и пара
$(\tilde X,\tilde\Delta+E)$ plt, где $\tilde\Delta$ определено равенством
$K_{\tilde X}+\tilde\Delta+(1-A_{(X,\Delta)}(E))E=\pi^*(K_X+\Delta)$. (Для дивизора $E\subset X$ на
самом $X$ берём $\pi=\id$; тогда условие --- plt-ость пары $(X,\Delta+E)$.) По формуле
присоединения с different (определение~\ref{def:different})
\[
K_E+\Delta_E=(K_{\tilde X}+\tilde\Delta+E)|_E ,
\]
и $(E,\Delta_E)$ --- klt-пара.

\begin{definition}\label{def:refined-series}
Определим \emph{многоградуированную линейную систему} на $E$:
\[
W^E_{m,j}=\im\Big(H^0(\tilde X,m\pi^*L-jE)\to H^0(E,(m\pi^*L-jE)|_E)\Big),\qquad m\ge0,\ j\ge0,
\]
и её объём $\vol(W^E_{\bullet\bullet})=\lim_m\sum_{j\ge0}\dim W^E_{m,j}/(m^n/n!)$. Для дивизора $F$ над
$E$ и $t\ge0$ --- фильтрацию $\mathcal F_F^tW^E_{m,j}=\{s\in W^E_{m,j}:\ord_F(s)\ge mt\}$, объёмы
$\vol(\mathcal F_FW^E_{\bullet\bullet})^t$ и
\[
S(W^E_{\bullet\bullet};F)=\frac1{\vol W^E_{\bullet\bullet}}\int_0^\infty\vol(\mathcal F_FW^E_{\bullet\bullet})^t\,dt,\qquad
\delta_{Z'}(E,\Delta_E;W^E_{\bullet\bullet})=\inf_{F:\,Z'\subset c_E(F)}\frac{A_{(E,\Delta_E)}(F)}{S(W^E_{\bullet\bullet};F)} .
\]
\end{definition}

Смысл: $W^E_{m,j}$ --- «$j$-й слой» фильтрации $\mathcal F_E$, ограниченный на $E$; из точной
последовательности $0\to\cO(m\pi^*L-(j+1)E)\to\cO(m\pi^*L-jE)\to\cO_E(\dots)\to0$ видно, что
$\sum_j\dim W^E_{m,j}=\dim V_m$, откуда $\vol(W^E_{\bullet\bullet})=\vol(V_\bullet)$.

\begin{theorem}[{Аббан--Чжуан, \cite[3.2]{AZ}}]\label{thm:AZ}
Пусть $E$ --- дивизор plt-типа над $X$, $Z\subset c_X(E)$. Тогда
\[
\delta_Z(X,\Delta;V_\bullet)\ \ge\ \min\Big\{\frac{A_{(X,\Delta)}(E)}{S(V_\bullet;E)},\ \inf_{Z'}\delta_{Z'}(E,\Delta_E;W^E_{\bullet\bullet})\Big\},
\]
где инфимум берётся по подмногообразиям $Z'\subset E$ с $\pi(Z')=Z$.
\end{theorem}

Итак, $\delta$ в точке оценивается снизу «$\delta$-значением самого $E$» и $\delta$-инвариантом
\emph{меньшей размерности} --- пары $(E,\Delta_E)$ с неполной многоградуированной системой
$W^E_{\bullet\bullet}$. Применяя теорему повторно, доходим до кривых, где всё считается явно.

\subsection{Случай поверхностей: явная формула}

Пусть $X$ --- поверхность, $E$ --- кривая plt-типа (тогда $E$ гладкая), $Z=p\in X$, $Z'$ --- точки
$q\in E$ над $p$ (если $E\subset X$, то $q=p$). Дивизоры над кривой $E$ --- это её точки $F=q$,
$A_{(E,\Delta_E)}(q)=1-\ord_q\Delta_E$. Пусть $\tau$ --- псевдоэффективный порог
$\pi^*L-uE$, и $\pi^*L-uE=P(u)+N(u)$ --- разложение Зарисского при $0\le u\le\tau$.

\begin{theorem}[{\cite[1.106]{ACCFKMSV}}]\label{thm:AZ-surface}
В этих обозначениях положим $t(u)=P(u)\cdot E$ (степень $\Q$-дивизора $P(u)|_E$ на кривой $E$).
Тогда для $q\in E$
\[
S(W^E_{\bullet\bullet};q)=\frac2{\vol L}\int_0^\tau\Big(t(u)\cdot\ord_q\big(N(u)|_E\big)
+\int_0^{t(u)}\big(t(u)-v\big)dv\Big)du
=\frac2{\vol L}\int_0^\tau\Big(t(u)\,\ord_q\big(N(u)|_E\big)+\frac{t(u)^2}2\Big)du .
\]
Если $q\notin\Supp N(u)$ при всех $u$ (в частности, если $N(u)=0$), первое слагаемое отсутствует.
\end{theorem}

Пояснение. При $u=j/m$ пространство $W^E_{m,j}$ --- это (с точностью до поправки на $N(u)$, дающей
первое слагаемое) образ сечений $mP(u)$ на $E$, т.\,е. почти полная линейная система $\Q$-дивизора
$mP(u)|_E$ степени $m\,t(u)$ на кривой $E$. Для полной линейной системы дивизора степени $d$ на
кривой подпространство сечений, обращающихся в нуль в $q$ с порядком $\ge v$, имеет размерность
$\approx d-v$ при $v\le d$ (Риман--Рох), откуда $\vol(\mathcal F_qW)^v=\max\{t(u)-v,0\}$ и
внутренний интеграл $\int_0^{t(u)}(t(u)-v)dv=t(u)^2/2$. В обзоре этот интеграл записан как
$\int_0^{t(u)}\max\{\ord_q(P(u)|_E)-v,0\}dv$, где $\ord_q(P(u)|_E)$ обозначает степень $P(u)|_E$.
Множитель $2=n$ --- из нормировки $m^n/n!$ при переходе от $n=2$ к $n=1$.

\begin{example}[Гладкая кубическая поверхность K-полустабильна; пример~4.10 обзора]\label{ex:AZ-cubic}
Пусть $X\subset\PP^3$ --- гладкая кубика, $p\in X$, $\Delta=0$, $L=-K_X$, $(-K_X)^2=3$. Выберем
гладкую кривую $E\in|-K_X|$ через $p$ (плоское сечение; гладкое через данную точку существует, так
как $-K_X$ очень обилен). Пара $(X,E)$ plt (даже snc), $E$ --- plt-типа с $\pi=\id$, $\Delta_E=0$.

\emph{Шаг 1: $A/S$ для $E$.} $A_X(E)=1$, $-K_X-tE=(1-t)(-K_X)$ nef при $t\le1$, объём $3(1-t)^2$:
\[
S_X(E)=\frac13\int_0^13(1-t)^2dt=\frac13,\qquad\frac{A_X(E)}{S_X(E)}=3 .
\]
\emph{Шаг 2: $\delta$ на кривой $E$.} $P(u)=(1-u)(-K_X)$, $N(u)=0$, $\tau=1$;
$P(u)|_E=(1-u)(-K_X)|_E$ --- дивизор степени $3(1-u)$ на $E$ (кубика степени 3), т.\,е.
$\ord_p(P(u)|_E)=t(u)=3(1-u)$ и
\[
S(W^E_{\bullet\bullet};p)=\frac23\int_0^1\int_0^{3(1-u)}\big(3(1-u)-v\big)dv\,du
=\frac23\int_0^1\frac{9(1-u)^2}2du=\frac23\cdot\frac32=1 .
\]
$A_{(E,0)}(p)=1$, так что $\delta_p(E,0;W^E_{\bullet\bullet})=1$.

\emph{Шаг 3: теорема~\ref{thm:AZ}.} $\delta_p(X;-K_X)\ge\min\{3,1\}=1$ для всех $p\in X$, откуда
$\delta(X)\ge1$: \emph{всякая гладкая кубическая поверхность K-полустабильна}. (На самом деле
K-стабильна: это следствие~\ref{cor:dP-le4}, а также последняя фраза теоремы~1.2 в \cite{AZ}:
при строгом неравенстве в одном из членов и некоторых условиях получается $\delta>1$.)
\end{example}

Отметим, насколько этот способ «дешевле» критерия Тяня: вместо вычисления $\lct$ всех
антиканонических дивизоров --- два интеграла.

\begin{example}[Пара $(\PP(1,1,2),\frac12Q)$ K-полустабильна; пример~4.11 и упражнение~4.1(1) обзора]\label{ex:AZ-P112}
Пусть $X=\PP(1,1,2)$, $\Delta=\frac12Q$, $L=-K_X-\Delta=\cO(3)$, $\vol L=\frac92$. В
примере~\ref{ex:P112-pair} мы видели, что это единственный кандидат на K-полустабильность в семействе
$(X,cQ)$.

\emph{Гладкая точка $p$.} Пусть $E$ --- образующая конуса через $p$ (дивизор класса $\cO(1)$,
$E\cong\PP^1$, $E^2=\frac12$, $E\cdot Q=1$; $E$ проходит через вершину $p_0$). Пара
$(X,\Delta+E)$ plt: вне вершины она snc, а в вершине пара $(X,E)$ plt с дискрепантностью
$-\frac12$ (раздел~\ref{sec:mmp}, пример с образующей конуса). $E\not\subset\Supp\Delta$,
$A_{(X,\Delta)}(E)=1$. Дивизор $L-tE=\cO(3-t)$ обилен при
$t<3$, $\tau=3$, $\vol(L-tE)=\frac{(3-t)^2}2$:
\[
S(E)=\frac29\int_0^3\frac{(3-t)^2}2dt=\frac19\cdot9=1,\qquad\frac{A}{S}=1 .
\]
Далее, $P(u)=L-uE$, $N(u)=0$, $P(u)|_E$ имеет степень $(3-u)E\cdot E=\frac{3-u}2$:
\[
S(W^E_{\bullet\bullet};p)=\frac2{9/2}\int_0^3\frac12\Big(\frac{3-u}2\Big)^2du=\frac49\cdot\frac12\cdot\frac14\cdot9=\frac12 .
\]
Different: $\Delta_E=\frac12(Q\cap E)+\frac12p_0$ --- точка $Q\cap E$ с коэффициентом $\frac12$
плюс вклад особой точки $\frac12(1,1)$ с коэффициентом $1-\frac12$ (определение~\ref{def:different}).
Поскольку $p\ne p_0$, имеем $A_{(E,\Delta_E)}(p)\ge\frac12$ и $\delta_p(E,\Delta_E;W)\ge\frac{1/2}{1/2}=1$. По
теореме~\ref{thm:AZ} $\delta_p(X,\Delta;L)\ge1$.

\emph{Вершина $p_0$.} Теперь возьмём $E$ --- исключительную кривую разрешения
$\pi\colon Y=\mathbb F_2\to X$ ($E^2=-2$); $-E$ $\pi$-обилен, $(Y,\tilde\Delta+E)$ plt ($Y$ гладко,
$\tilde\Delta=\frac12\tilde Q$ не пересекает $E$), $\Delta_E=0$. Из примера~\ref{ex:P112-pair} (при
$c=\frac12$) $A_{(X,\Delta)}(E)=1$, $S(E)=\frac23\cdot\frac32=1$, $\frac AS=1$. Далее
$\pi^*L-uE=\frac32(E+2F)-uE$, nef при $0\le u\le\frac32=\tau$ (пересечение с $F$ равно $\frac32-u$,
с $E$ равно $2u$), $N(u)=0$, и $P(u)|_E$ имеет степень $P(u)\cdot E=2u$. Для $q\in E$:
\[
S(W^E_{\bullet\bullet};q)=\frac2{9/2}\int_0^{3/2}\frac{(2u)^2}2du=\frac49\cdot2\cdot\frac{(3/2)^3}3=\frac49\cdot\frac94=1,
\qquad A_{(E,0)}(q)=1,\qquad\delta_q(E,0;W)=1 .
\]
Итак, $\delta_{p_0}(X,\Delta;L)\ge\min\{1,1\}=1$.

\emph{Вывод.} $\delta(\PP(1,1,2),\tfrac12Q)\ge1$, пара K-полустабильна. Это решает упражнение~4.1(1).
Как отмечалось, $(\PP(1,1,2),\frac12Q)=(\PP^1\times\PP^1)/\langle\sigma\rangle$ --- фактор по
инволюции перестановки сомножителей; K-полустабильность пары равносильна K-полустабильности
$\PP^1\times\PP^1$ (эквивариантная K-стабильность и факторы \cite{LZ22}).
\end{example}

\begin{remark}
Теория Аббана--Чжуана --- основной инструмент книги \cite{ACCFKMSV}, в которой определена
K-(полу/поли)стабильность общего представителя каждого из 105 семейств гладких трёхмерных
многообразий Фано; во многих случаях --- всех представителей. Для кубических поверхностей она даёт
короткое доказательство K-полустабильности гладких кубик; но для описания \emph{границы}
пространства модулей (особых K-полустабильных кубик) нужен ещё один инструмент --- нормализованный
объём (раздел~\ref{sec:cubic}).
\end{remark}

% ============================================================
\section{Модули K-полустабильных многообразий Фано}\label{sec:moduli}
% ============================================================

Это \S5 обзора. Мы объясним, почему пространства модулей \emph{всех} многообразий Фано безнадёжно
плохи, и сформулируем теорему о K-модулях, которая все эти проблемы решает.

\subsection{Чего мы хотим и почему это не получается}

Напомним (\S\ref{ss:moduli-def}): от пространства модулей хочется ограниченности, отделимости,
собственности и редуктивности (а лучше конечности) групп автоморфизмов. Семейства --- 
$\Q$-горенштейновы (определение~\ref{def:QGor-family}), так что обильность $-K_{X_t}$ --- открытое
условие (пример~5.2 обзора), а $(-K_{X_t})^n$ постоянен.

\begin{question}[Собственность]
Для многообразий общего типа (с обильным $K_X$) пределы семейств строятся канонически: предел ---
$\Proj$ канонического кольца $R(K_{\cX})$, конечно порождённого по БКХМ \cite{BCHM}. Для Фано
$-K_X$ обилен, и аналогичной канонической конструкции нет: предел семейства гладких Фано можно
«подправить» бирациональным преобразованием (флопом, или просто другим выбором семейства над
проколотой кривой). Как же строить пределы функториально?
\end{question}

\begin{example}[Неотделимость; пример~5.4 обзора]\label{ex:nonseparated}
Пусть $X=\PP^1$, вложенная коникой в $\PP^2$. Семейство $\{xz=t^2y^2\}$ (пример~\ref{ex:conic-tc})
имеет все слои при $t\ne0$ изоморфными $\PP^1$, а при $t=0$ --- $X_0=\PP^1\cup\PP^1$ (две прямые,
склеенные в точке). Это \emph{изотривиальное} вырождение: общие слои изоморфны, а предел ---
другой. Если бы мы допускали в модули все (даже ненормальные) вырождения с фиксированным объёмом
--- а $\vol(-K_{X_0})=2\vol(-K_{\PP^1}-\text{точка})=2=\vol(-K_{\PP^1})$ (нормализация $X_0$ --- две
копии $(\PP^1,\{\text{точка}\})$ с кондуктором-точкой), --- то у постоянного семейства $\PP^1$ было бы
два разных предела: $\PP^1$ и $X_0$. Пространство модулей неотделимо. (Для многообразий с обильным
$K$ такого не бывает.)
\end{example}

\begin{example}[Неограниченность: тройки Маркова; пример~5.5 обзора]\label{ex:Markov}
Чтобы построить проективное пространство модулей, объекты надо вложить в одно $\PP^N$ и
воспользоваться схемой Гильберта. Но лог-терминальных многообразий Фано с фиксированными
размерностью и объёмом бесконечно много деформационных типов. Пусть $(a,b,c)$ --- решение
\emph{уравнения Маркова}
\[
a^2+b^2+c^2=3abc
\]
в натуральных числах (\emph{тройка Маркова}). Все решения получаются из $(1,1,1)$ перестановками и
\emph{мутациями} $(a,b,c)\mapsto(a,b,3ab-c)$; начало дерева Маркова:
\[
(1,1,1)\to(1,1,2)\to(1,2,5)\to\{(1,5,13),\,(2,5,29)\}\to\{(1,13,34),\,(5,13,194),\,(5,29,433),\,(2,29,169)\}\to\dots
\]
Рассмотрим $X_{a,b,c}=\PP(a^2,b^2,c^2)$. Числа $a,b,c$ попарно взаимно просты (если $p\mid a,b$,
то $p\mid c^2$, и т.\,д.), так что веса нормализованы, и
\[
(-K_X)^2=\frac{(a^2+b^2+c^2)^2}{a^2b^2c^2}=\frac{9a^2b^2c^2}{a^2b^2c^2}=9 .
\]
Получаем бесконечную серию $\PP^2,\ \PP(1,1,4),\ \PP(1,4,25),\ \PP(1,25,169),\ \PP(4,25,841),\dots$
лог-терминальных (торических, с циклическими факторособенностями) поверхностей дель Пеццо степени~9.
Более того, все они --- \emph{$\Q$-горенштейновы вырождения $\PP^2$} (Хакинг--Прохоров \cite{HP}:
это в точности все $\Q$-горенштейновы вырождения $\PP^2$ с факторособенностями). Итак, множество
klt-дель Пеццо степени 9 неограничено, и у $\PP^2$ «бесконечно много» разных пределов.

Решим упражнения~5.1(1,2) обзора.

\emph{(1) Только тройки Маркова.} Пусть $\PP(p,q,r)$ нормализовано (т.\,е. $p,q,r$ попарно
взаимно просты) и $(-K)^2=\frac{(p+q+r)^2}{pqr}=9$. Тогда $9pqr=(p+q+r)^2$ --- точный квадрат; так
как $9$ --- квадрат, а $p,q,r$ попарно взаимно просты, каждое из $p,q,r$ --- точный квадрат:
$p=a^2$, $q=b^2$, $r=c^2$, и $3abc=a^2+b^2+c^2$.

\emph{(2a) Особенности $\PP(a^2,b^2,c^2)$ --- $T$-особенности.} Вершина с весом $a^2$ --- особенность
$\frac1{a^2}(b^2,c^2)\cong\frac1{a^2}(1,q)$, $q\equiv c^2b^{-2}\pmod{a^2}$. Из уравнения Маркова
$c^2\equiv3abc-b^2\pmod{a^2}$, поэтому $q\equiv3acb^{-1}-1=ak-1$, где $k\equiv3cb^{-1}\pmod a$.
Так как $3\nmid a$ (если $3\mid a$, то $3\mid b^2+c^2$, откуда $3\mid b,c$ --- противоречие) и
$\gcd(c,a)=1$, имеем $\gcd(k,a)=1$. Итак, особенность имеет вид $\frac1{a^2}(1,ak-1)$ ---
$T$-особенность с $d=1$, $n=a$ (теорема~\ref{thm:KSB}), и она $\Q$-горенштейново сглаживаема
(пример~\ref{ex:T-sing}(3)). Проверка: $\PP(1,1,4)$ --- $\frac14(1,1)$; $\PP(1,4,25)$ ---
$\frac14(1,1)$ и $\frac1{25}(1,4)$; $\PP(4,25,841)$ --- $\frac14(1,1)$, $\frac1{25}(1,4)$,
$\frac1{841}(1,637)$ ($637=29\cdot22-1$).

\emph{(2b) Нет локально-глобальных препятствий.} Для нормальной поверхности $X$ с
факторособенностями препятствия к продолжению локальных деформаций до глобальных лежат в
$H^2(X,T_X)$ ($T_X$ --- рефлексивный касательный пучок). Для дель Пеццо $H^2(X,T_X)=H^0(X,\Omega^{[1]}_X\otimes\cO(K_X))^*=0$
(двойственность Серра и отрицательность $K_X$; см. \cite[Prop.~3.1]{HP}). Поэтому локальные
$\Q$-горенштейновы сглаживания особенностей склеиваются в глобальное $\Q$-горенштейново сглаживание
$X\rightsquigarrow X_t$, где $X_t$ --- гладкая дель Пеццо с $K^2=9$ (постоянство $K^2$), т.\,е.
$X_t\cong\PP^2$. Явное семейство для $\PP(1,1,4)$ дано в \S\ref{sec:families}: $\{xy=z^2+tw\}\subset\PP(1,1,1,2)$.
\end{example}

\begin{example}[Нередуктивные группы автоморфизмов; пример~5.6 обзора]
Для «хорошего фактора» нужны по крайней мере редуктивные стабилизаторы. Но $\Aut(\mathbb F_1)$ и
$\Aut(\Bl_2\PP^2)$ нередуктивны (пример~\ref{ex:Aut-F1}); для многообразий общего типа группы
автоморфизмов конечны.
\end{example}

Мораль: пространства модулей всех многообразий Фано не существует в хорошем смысле. Все четыре
проблемы снимает K-стабильность.

\subsection{Теорема о K-модулях}

\begin{theorem}[{Ограниченность; Цзян \cite{Jia20}, Биркар \cite{Bir19}}]\label{thm:boundedness}
K-полустабильные $\Q$-Фано многообразия размерности $n$ и объёма $(-K)^n=V$ образуют ограниченное
семейство: существует $N$ такое, что все они вкладываются в $\PP^N$ системой $|-mK|$ при одном $m$.
\end{theorem}

Идея: K-полустабильность даёт нижнюю оценку $\alpha$-инварианта ($\alpha\ge\frac1{n+1}\delta\ge\frac1{n+1}$,
\cite{BJ}), а Фано с $\alpha$ и объёмом, отделёнными от нуля, ограничены (Цзян), что в свою очередь
опирается на теорему Биркара об ограниченности (BAB). Пример~\ref{ex:Markov} не противоречит:
все $\PP(a^2,b^2,c^2)$, кроме $\PP^2$, K-нестабильны (пример~\ref{ex:wps-unstable}).

\begin{theorem}[{Открытость; Блюм--Лю--Сюй \cite{BLX}, Сюй \cite{Xu20}}]\label{thm:openness}
В $\Q$-горенштейновом семействе $\Q$-Фано многообразий $\cX\to T$ множество $\{t: X_t\ \text{K-полустабильно}\}$
открыто по Зарисскому; то же для K-стабильности.
\end{theorem}

Идея: функция $t\mapsto\delta(X_t)$ полунепрерывна снизу, и её значения можно вычислять по
конечному числу дивизоров.

\begin{theorem}[{Собственность и проективность; \cite{LWX}, \cite{BX}, \cite{BHLLX}, \cite{LXZ}, \cite{XZ}}]\label{thm:properness}
Стек K-полустабильных $\Q$-Фано многообразий размерности $n$ и объёма $V$ собственный: для
$\Q$-горенштейнова семейства K-полустабильных Фано над проколотой кривой $C\setminus0$ после
конечной замены базы существует заполнение K-полустабильным слоем над $0$, причём K-полистабильное
заполнение единственно с точностью до изоморфизма (\cite{BX}). Пространство модулей
K-полистабильных Фано проективно (обильность CM-расслоения, \cite{XZ}).
\end{theorem}

\begin{theorem}[{Редуктивность; \cite{ABHLX}}]\label{thm:reductivity}
Группа автоморфизмов K-полистабильного $\Q$-Фано многообразия редуктивна.
\end{theorem}

Это алгебраический аналог теоремы Мацусимы~\ref{thm:Matsushima}. Всё вместе даёт:

\begin{theorem}[Теорема о K-модулях]\label{thm:Kmoduli}
Для $n\ge1$ и $V>0$ существует алгебраический стек Артина конечного типа $\cM^{\mathrm{K}}_{n,V}$,
параметризующий $\Q$-горенштейновы семейства K-полустабильных $\Q$-Фано многообразий размерности
$n$ и объёма $V$, и его хорошее пространство модулей $M^{\mathrm{K}}_{n,V}$ --- проективная схема,
точки которой взаимно однозначно соответствуют K-полистабильным $\Q$-Фано многообразиям размерности
$n$ и объёма $V$.
\end{theorem}

В терминах проблем из \S\ref{ss:moduli-def}: вырождение в ненормальные многообразия исключено
(K-полустабильные klt, в частности нормальны: теорема~\ref{thm:Odaka}); неограниченность снята
теоремой~\ref{thm:boundedness}; собственность и отделимость (единственность полистабильного предела)
--- теорема~\ref{thm:properness}; редуктивность --- теорема~\ref{thm:reductivity}. Каждая
K-полустабильная $X$ изотривиально вырождается в единственную K-полистабильную $X^{\mathrm{ps}}$
(аналог $S$-эквивалентности), и точка $[X^{\mathrm{ps}}]\in M^{\mathrm{K}}$ --- образ $X$.

\begin{remark}
Всё сказанное переносится на лог-Фано пары $(X,D)$ с заменой $-K_X$ на $-(K_X+D)$; для пар
$(X,cD)$ с переменным $c$ пространства модулей меняются при переходе $c$ через «стены»
(\emph{wall-crossing}, \cite{ADL}).
\end{remark}

\begin{example}[Общая кубика K-полустабильна; упражнение~5.1(3)]
Семейство всех кубических поверхностей $\PP^{19}$ (точнее, его открытое подмножество нормальных
кубик) $\Q$-горенштейново. Кубика $3A_2$ K-полистабильна (пример~\ref{ex:3A2-polystable}), а кубика
Ферма K-стабильна; по теореме~\ref{thm:openness} K-(полу)стабильные кубики образуют открытое
непустое, значит, плотное подмножество. (На самом деле все гладкие кубики K-стабильны,
следствие~\ref{cor:dP-le4}.)
\end{example}

\begin{example}[K-полустабильное, но не K-полистабильное]
Упражнение~5.1(4) просит найти гладкое трёхмерное многообразие Фано, K-полустабильное, но не
K-полистабильное. Рецепт: взять K-полистабильное $X_0$ с бесконечной группой автоморфизмов и
гладкое $X$ с изотривиальным вырождением $X\rightsquigarrow X_0$ (например, специальную
тест-конфигурацию с $\Fut=0$, индуцированную $\Gm\subset\Aut(X_0)$, «развёрнутую» в деформацию $X_0$).
По открытости $X$ K-полустабильно, а по единственности полистабильного предела его полистабильная
модель --- $X_0\not\cong X$. Примеры такого рода среди гладких трёхмерных Фано разобраны в \cite{ACCFKMSV}; в размерности 2 для гладких дель Пеццо такого не бывает (теорема~\ref{thm:dP-table}),
а среди особых кубик пример --- любая кубика с особенностью $A_2$, отличная от $3A_2$
(теорема~\ref{thm:GIT=K}).
\end{example}

% ============================================================
\section{Нормализованный объём и K-модули кубических поверхностей}\label{sec:cubic}
% ============================================================

Это \S6 обзора --- кульминация текста. Теорема о K-модулях говорит, что пространство
$M^{\mathrm K}_{2,3}$ существует; мы хотим \emph{увидеть} его. Для этого нужно ответить на вопрос:
какие особые поверхности дель Пеццо степени 3 K-полустабильны? Инструмент --- \emph{нормализованный
объём} Ли, связывающий локальные инварианты особенностей с глобальным объёмом $(-K_X)^n$.

\subsection{Оценка объёма}

\begin{theorem}[{Фудзита \cite{Fuj18}, Лю \cite{Liu18}}]\label{thm:volume-bound}
Пусть $X$ --- K-полустабильное $\Q$-Фано многообразие размерности $n$. Тогда
\[
(-K_X)^n\le(n+1)^n ,
\]
причём равенство достигается тогда и только тогда, когда $X\cong\PP^n$.
\end{theorem}

\begin{proof}[Доказательство неравенства]
Выберем гладкую точку $x\in X$ и пусть $\mu\colon Y=\Bl_xX\to X$ --- раздутие с исключительным
дивизором $E\cong\PP^{n-1}$. По валюативному критерию (теорема~\ref{thm:valuative})
$\beta_X(E)\ge0$, т.\,е.
\begin{equation}\label{eq:beta-ineq}
A_X(E)\cdot(-K_X)^n\ \ge\ \int_0^\infty\vol(-K_X-tE)\,dt .
\end{equation}
\emph{Лог-дискрепантность.} По теореме~\ref{thm:blowup-K} $K_Y=\mu^*K_X+(n-1)E$, откуда
$A_X(E)=1+(n-1)=n$.

\emph{Оценка объёма снизу.} Пусть $t\in\Q_{\ge0}$ и $m$ таково, что $mt\in\Z$. По определению
\[
\vol(\mu^*(-K_X)-tE)=\lim_{m\to\infty}\frac{h^0\big(Y,\cO_Y(m\mu^*(-K_X)-mtE)\big)}{m^n/n!} .
\]
Так как $\mu_*\cO_Y(-kE)=\mm_x^k$ для $k\ge0$ (функции, обращающиеся в нуль с порядком $\ge k$ в $x$;
пример~\ref{ex:blowup-point-ord}) и $\mu_*\mu^*\cO(-mK_X)=\cO(-mK_X)$ (формула проекции),
\[
h^0\big(Y,m\mu^*(-K_X)-mtE\big)=h^0\big(X,\cO_X(-mK_X)\otimes\mm_x^{mt}\big)
\ \ge\ h^0\big(X,-mK_X\big)-\length\big(\cO_{X,x}/\mm_x^{mt}\big) .
\]
Последнее неравенство --- из точной последовательности $0\to\mm_x^{k}\to\cO_X\to\cO_X/\mm_x^k\to0$,
тензорно умноженной на $\cO(-mK_X)$ (локально свободный около $x$ пучок, так что третий член не
меняется): размерность пространства сечений первого пучка не меньше разности размерностей сечений
второго и третьего. Далее, $\length(\cO_{X,x}/\mm_x^k)=\binom{n+k-1}{n}=\frac{k^n}{n!}+O(k^{n-1})$
(число мономов степени $<k$ от $n$ переменных; это упражнение~6.3(4) обзора). Делим на
$m^n/n!$ и переходим к пределу:
\[
\vol(-K_X-tE)\ \ge\ (-K_X)^n-\lim_{m\to\infty}\frac{(mt)^n/n!}{m^n/n!}=(-K_X)^n-t^n .
\]
(В общем случае вместо $t^n$ появляется $\vol(\ord_E)\cdot t^n$ --- объём нормирования, см. ниже;
здесь $\vol(\ord_E)=1$.)

\emph{Подстановка.} Так как $\vol\ge0$, из \eqref{eq:beta-ineq}:
\[
n\,(-K_X)^n\ \ge\ \int_0^\infty\max\{(-K_X)^n-t^n,0\}\,dt=\int_0^{V^{1/n}}(V-t^n)\,dt
=V^{1+\frac1n}\Big(1-\frac1{n+1}\Big)=\frac n{n+1}V^{\frac{n+1}n},
\]
где $V=(-K_X)^n$. Итак, $V\ge\frac1{n+1}V^{1+1/n}$, т.\,е. $V^{1/n}\le n+1$, $V\le(n+1)^n$.
\end{proof}

Для равенства нужны все неравенства как равенства, в частности $\vol(-K_X-tE)=(-K_X)^n-t^n$ на всём
промежутке $[0,n+1]$; отсюда Лю \cite{Liu18} выводит, что $X\cong\PP^n$.

\subsection{Нормирования и нормализованный объём}

До сих пор мы работали с дивизориальными нормированиями $\ord_E$. Нормализованный объём естественно
определять для всех нормирований, но читатель может всюду думать о дивизорах.

\begin{definition}
Пусть $x\in X=\Spec R$ --- точка нормального многообразия (аффинного около $x$). \emph{Нормирование
с центром в $x$} --- отображение $v\colon R\setminus0\to\R_{\ge0}$ с $v(fg)=v(f)+v(g)$,
$v(f+g)\ge\min\{v(f),v(g)\}$, $v(\C^*)=0$ и $v>0$ ровно на $\mm_x$; множество таких $v$ ---
$\Val_{x,X}$. Для дивизора $E$ над $X$ с $c_X(E)=x$ нормирование $\ord_E$ лежит в $\Val_{x,X}$.
\emph{Объём нормирования} \cite{ELS}:
\[
\vol(v)=\lim_{k\to\infty}\frac{\length(R/\fa_k(v))}{k^n/n!},\qquad\fa_k(v)=\{f\in R:\ v(f)\ge k\} .
\]
\emph{Лог-дискрепантность} $A_X(v)$ (Йонссон--Мустацэ \cite{JM}) --- продолжение функции
$E\mapsto A_X(E)$ на все нормирования (для дивизориальных $A_X(\ord_E)=A_X(E)$; для $v=c\,\ord_E$
$A_X(v)=cA_X(E)$).
\end{definition}

\begin{definition}[{Ли \cite{Li18}}]\label{def:nvol}
Пусть $x\in X$ --- klt-особенность. \emph{Нормализованный объём} нормирования $v\in\Val_{x,X}$ ---
\[
\hvol(v)=A_X(v)^n\cdot\vol(v),
\]
а \emph{локальный объём} особенности ---
$\hvol(x,X)=\inf_{v\in\Val_{x,X}}\hvol(v)$.
\end{definition}

Нормировка $A^n\cdot\vol$ выбрана так, чтобы $\hvol(cv)=\hvol(v)$: $A(cv)=cA(v)$,
$\vol(cv)=c^{-n}\vol(v)$. В доказательстве теоремы~\ref{thm:volume-bound} мы использовали
$A_X(\ord_E)=n$, $\vol(\ord_E)=1$ для раздутия гладкой точки: $\hvol(\ord_E)=n^n$.

\begin{remark}[Нормализованный объём и K-стабильность конусов; замечание~6.4 обзора]\label{rem:cone-nvol}
Пусть $V$ --- $\Q$-Фано, $r$ таково, что $-rK_V$ --- Картье, $X=C(V,-rK_V)$ --- аффинный конус с
вершиной $x$. Вершина --- klt-точка (klt-ость конуса равносильна тому, что $V$ Фано, т.\,е. $-K_V$
обилен, и $V$ klt). Раздутие вершины $\mu\colon Y\to X$ имеет исключительный дивизор $V_0\cong V$.
Теорема (Ли \cite{Li17}, Ли--Сюй \cite{LX20}): \emph{$V$ K-полустабильно тогда и только тогда, когда
$\ord_{V_0}$ минимизирует $\hvol$ на $\Val_{x,X}$}. Так K-стабильность многообразия Фано
превращается в \emph{локальное} свойство особенности --- вершины конуса над ним; это ещё одна
алгебраическая формулировка K-полустабильности (обобщающая теорему Фудзиты--Лю).
\end{remark}

\begin{theorem}[{Локально-глобальная теорема; Ли--Лю \cite{LL19}}]\label{thm:local-global}
Пусть $X$ --- K-полустабильное $\Q$-Фано многообразие размерности $n$. Тогда для любой точки
$x\in X$
\[
(-K_X)^n\le\Big(1+\frac1n\Big)^n\hvol(x,X) .
\]
\end{theorem}
\begin{proof}
Повторяем доказательство теоремы~\ref{thm:volume-bound} для произвольного дивизора $E$ над $X$ с
центром $x$ (или нормирования $v$), не заменяя $A=A_X(E)$ и $\vol(\ord_E)$ на $n$ и $1$:
из $\beta(E)\ge0$ и $\vol(-K_X-tE)\ge V-\vol(\ord_E)t^n$ получаем
$A\,V\ge\frac n{n+1}V^{1+1/n}\vol(\ord_E)^{-1/n}$, т.\,е. $V\le\big(\frac{n+1}n\big)^nA^n\vol(\ord_E)=\big(1+\frac1n\big)^n\hvol(\ord_E)$,
и берём инфимум по $E$ (дивизориальные нормирования плотны в нужном смысле).
\end{proof}

Это \emph{локально-глобальный принцип}: глобальный инвариант $(-K_X)^n$ ограничен локальными
инвариантами особенностей. Чем хуже особенность (меньше $\hvol$), тем меньше должен быть объём; в
частности, K-полустабильные многообразия Фано большого объёма имеют мягкие особенности.

\begin{theorem}[{Свойства локального объёма}]\label{thm:nvol-properties}
Пусть $\dim X=n$.
\begin{enumerate}
\item (\cite{dFEM}, \cite{Li18}) Если $x\in X$ гладкая, то $\hvol(x,X)=n^n$.
\item (\cite{LX19}) Всегда $\hvol(x,X)\le n^n$, и равенство --- только в гладких точках. Вместе с
теоремой~\ref{thm:local-global} это даёт теорему~\ref{thm:volume-bound}.
\item (\cite{Liu18}) Если $x\in X=\A^n/G$ --- факторособенность по конечной $G\subset\GL_n(\C)$,
действующей свободно в коразмерности~1, то $\hvol(x,X)=n^n/|G|$. Более общо, для фактора
произвольной особенности по $G$: $\hvol(x,X/G)=\hvol(x,X)/|G|$ при свободном в коразмерности 1
действии; в частности, $\hvol(x,X/G)\le n^n/|G|$. Следствие: \emph{если $X$ K-полустабильно и
$x\in X$ --- факторособенность $\A^n/G$, то}
\begin{equation}\label{eq:quotient-bound}
(-K_X)^n\le\frac{(n+1)^n}{|G|} .
\end{equation}
\item (\cite{LX19}) Если $n=2$ или $3$ и $x\in X$ --- особая точка, то $\hvol(x,X)\le2(n-1)^n$,
с равенством только для обыкновенной двойной точки ($A_1$: $x_1^2+\dots+x_{n+1}^2=0$).
Гипотетически обыкновенная двойная точка имеет второй по величине локальный объём во всех
размерностях.
\end{enumerate}
\end{theorem}

Пункт (3) делает всё вычислимым для поверхностей: лог-терминальные особенности поверхностей ---
факторособенности (теорема~\ref{thm:lt-surface}), и $\hvol(x,X)=4/|G|$.

\begin{definition}
Напомним (определение~\ref{def:cyclic-quot}): $\frac1n(a,b)$ --- фактор $\A^2/\mu_n$ по действию
$\zeta\cdot(x,y)=(\zeta^ax,\zeta^by)$. Для неё $|G|=n$ и $\hvol=4/n$.
\end{definition}

\begin{example}[$\PP(1,1,2)$ K-нестабильно: второе доказательство; пример~6.8 обзора]\label{ex:P112-nvol}
Пусть $X=\PP(1,1,2)$, $n=2$. Особая точка --- $\frac12(1,1)$, $|G|=2$. Далее, $X$ --- квадратичный
конус в $\PP^3$, и по формуле присоединения $K_X=(K_{\PP^3}+X)|_X=\cO_{\PP^3}(-2)|_X$, откуда
$K_X^2=\cO(-2)\cdot\cO(-2)\cdot\cO(2)=8$ (пересечение на $\PP^3$). Или: $-K_X=\cO(4)$ на $\PP(1,1,2)$
и $(-K_X)^2=16/2=8$. Если бы $X$ было K-полустабильно, то по \eqref{eq:quotient-bound}
$8=(-K_X)^2\le\frac{3^2}{2}=\frac92$ --- противоречие. Итак, $\PP(1,1,2)$ K-нестабильно (ср.
пример~\ref{ex:beta-P11n}, где $\beta(E)=-\frac13$).
\end{example}

\begin{example}[Все взвешенные проективные пространства, кроме $\PP^n$, K-нестабильны; упражнение~6.3(3)]\label{ex:wps-unstable}
Пусть $X=\PP(a_0,\dots,a_n)$, нормализованное, $a_{\max}=\max a_i>1$ (иначе $X=\PP^n$). Вершина
с весом $a_{\max}$ --- факторособенность $\A^n/\mu_{a_{\max}}$ (нормализованность гарантирует
свободу действия в коразмерности 1). Если $X$ K-полустабильно, то по \eqref{eq:quotient-bound}
$(-K_X)^n\le(n+1)^n/a_{\max}$. С другой стороны, по неравенству о средних
\[
(-K_X)^n=\frac{(\sum a_i)^n}{\prod a_i}\ \ge\ \frac{(n+1)^n(\prod a_i)^{n/(n+1)}}{\prod a_i}
=\frac{(n+1)^n}{(\prod a_i)^{1/(n+1)}} .
\]
Но $(\prod a_i)^{1/(n+1)}<a_{\max}$, если не все $a_i$ равны (а равными они быть не могут, кроме
случая $a_i=1$, в силу нормализованности). Значит, $(-K_X)^n>(n+1)^n/a_{\max}$ --- противоречие.
В частности, все $\PP(a^2,b^2,c^2)$ из примера~\ref{ex:Markov}, кроме $\PP^2$, K-нестабильны, и
неограниченность семейства дель Пеццо степени 9 исчезает при ограничении на K-полустабильные.
\end{example}

\begin{example}[Нет нетривиальных K-полустабильных вырождений $\PP^2$ и $\PP^1\times\PP^1$; упражнение~6.3(1)]
Пусть $X$ --- K-полустабильная klt поверхность дель Пеццо с $(-K_X)^2=9$ (кандидат на вырождение
$\PP^2$). По теореме~\ref{thm:local-global} для любой точки $x$: $9\le\frac94\hvol(x,X)$, т.\,е.
$\hvol(x,X)\ge4=2^2$, и по теореме~\ref{thm:nvol-properties}(2) $x$ гладкая. Значит, $X$ --- гладкая
дель Пеццо степени 9, т.\,е. $\PP^2$. (Это и есть случай равенства в теореме~\ref{thm:volume-bound}.)
Для $(-K_X)^2=8$: $8\le\frac94\hvol(x,X)$ даёт $\hvol(x,X)\ge\frac{32}9>2=2(n-1)^n$, и по
пункту (4) $x$ гладкая; гладкие дель Пеццо степени 8 --- $\PP^1\times\PP^1$ и $\mathbb F_1$, вторая
K-нестабильна. Итак, единственное K-полустабильное вырождение $\PP^1\times\PP^1$ --- оно само.
Отсюда: $M^{\mathrm K}_{2,9}$ и $M^{\mathrm K}_{2,8}$ --- по одной точке.
\end{example}

\begin{example}[Особенности $\frac14(1,1)$ и $\frac14(1,3)$; упражнение~6.3(2)]
Обе особенности $\Q$-горенштейново сглаживаемы ($\frac14(1,1)$ --- $T$-особенность,
$\frac14(1,3)=A_3$ --- Дю Валя), так что могут появляться на K-полустабильных вырождениях гладких
дель Пеццо. По \eqref{eq:quotient-bound} с $|G|=4$: $(-K_X)^2\le\frac94$. Для $A_3$ поверхность
горенштейнова, $K^2$ целое, и $d\le2$; пример: дель Пеццо степени 2 --- двулистное накрытие
$\PP^2$, разветвлённое в квартике с тэкнодой $y^2=x^4$, имеет над ней особенность
$z^2=y^2-x^4$, т.\,е. $A_3$. Для $\frac14(1,1)$: $K^2\le\frac94$; такие поверхности встречаются
среди K-полустабильных дель Пеццо степени 1 и 2 (см. \cite{OSS}).
\end{example}

\begin{example}[Пары на $\PP(1,1,4)$; упражнение~6.3(5)]\label{ex:P114-pair}
Пусть $X=\PP(1,1,4)$ с координатами $x,y$ (степени 1), $z$ (степени 4), особая точка $p_0=[0:0:1]$
типа $\frac14(1,1)$, $-K_X=\cO(6)$, $D\in|\cO(4d)|$, $c>0$; пара $(X,cD)$ лог-Фано при $6-4dc>0$,
т.\,е. $c<\frac3{2d}$.

\emph{(a) $D$ проходит через $p_0$.} Уравнение $D$: $f=\sum_{j\le d}z^jg_{4(d-j)}(x,y)$;
$p_0\in D$ означает, что коэффициент при $z^d$ равен нулю, и в орбифолдной карте $z=1$
(координаты $(x,y)$ на $\A^2/\mu_4$) функция $f$ имеет порядок $\ge4$ в начале координат. Пусть
$E$ --- исключительная кривая разрешения $Y=\mathbb F_4\to X$ ($E^2=-4$). Так как инвариантный
моном степени $4k$ обращается в нуль вдоль $E$ с порядком $k$ (пример~\ref{ex:P11n}),
$\ord_E(D)\ge1$, и $A_{(X,cD)}(E)=A_X(E)-c\,\ord_E(D)\le\frac24-c=\frac12-c$. Далее, $L=-K_X-cD=\cO(6-4dc)$,
и по вычислению примера~\ref{ex:beta-P11n} с $\ell=6-4dc$ вместо $n+2$:
$S(E)=\frac23\cdot\frac{\ell}4=1-\frac23dc$. Поэтому
\[
\beta_{(X,cD)}(E)\le\frac12-c-1+\frac23dc=-\frac12+c\Big(\frac{2d}3-1\Big)<0
\quad\text{при}\quad c<\frac{3}{2(2d-3)}\ (d\ge2)\ \text{или при всех }c\ (d=1).
\]
Так как $\frac3{2(2d-3)}>\frac3{2d}$, пара K-нестабильна при всех допустимых $c\in(0,\frac3{2d})$.

\emph{(b) $D$ не проходит через $p_0$.} Тогда $\hvol(p_0,X,cD)=\hvol(p_0,X)=\frac44=1$, и
локально-глобальная оценка для пар даёт
\[
(-K_X-cD)^2=\frac{(6-4dc)^2}4\le\frac94\cdot1\quad\Longleftrightarrow\quad6-4dc\le3\quad\Longleftrightarrow\quad c\ge\frac3{4d}.
\]
Итак, $(X,cD)$ может быть K-полустабильной лишь при $c\ge\frac3{4d}$. Это типичное рассуждение
при изучении «стен» в K-модулях пар.
\end{example}

\subsection{K-модули кубических поверхностей}\label{ss:Kmoduli-cubic}

Гладкая кубика $X\subset\PP^3$ --- дель Пеццо степени 3 с $-K_X=\cO_X(1)$ (пример~\ref{ex:hyp-Gor}).
Обратно, всякая гладкая дель Пеццо степени 3 --- кубика (теорема~\ref{thm:Gor-dP}; упражнение~6.3(6)).

\begin{question}
Как устроено пространство модулей $M^{\mathrm K}_{2,3}$ K-полистабильных дель Пеццо степени 3?
Обзор рассматривает его компоненту $M^{\sm}_{2,3}$, состоящую из \emph{сглаживаемых} поверхностей
(пределов гладких кубик); по результатам группы летней школы 2025~г. других компонент нет.
\end{question}

\begin{theorem}[{Теорема~6.10 обзора}]\label{thm:cubic-list}
Пусть $X$ --- K-полустабильная $\Q$-Фано поверхность степени $(-K_X)^2=3$, являющаяся
$\Q$-горенштейновым пределом гладких кубик. Тогда $X$ --- кубическая поверхность в $\PP^3$,
все особые точки которой имеют тип $A_1$ или $A_2$.
\end{theorem}
\begin{proof}
\emph{Шаг 1: особенности --- факторособенности порядка 2 или 3.} Пусть $x\in X$ --- особая точка.
По теореме~\ref{thm:Odaka} $X$ klt, а klt-особенности поверхностей --- факторособенности
$\A^2/G$, $G\subset\GL_2$ малая (теорема~\ref{thm:lt-surface}). По \eqref{eq:quotient-bound}
\[
3=(-K_X)^2\le\frac9{|G|},\qquad\text{т.\,е.}\qquad|G|\le3 ;
\]
так как $x$ особа, $|G|\ge2$. Итак, $|G|\in\{2,3\}$.

\emph{Шаг 2: перечисление.} Группы порядка 2 и 3 циклические; малая циклическая группа
$\mu_m\subset\GL_2$ диагонализуема, и её образующая действует весами $(\zeta^a,\zeta^b)$ с
$\gcd(a,m)=\gcd(b,m)=1$ (иначе есть псевдоотражения). Варианты:
\begin{enumerate}
\item $|G|=2$: $(-1,-1)$ --- особенность $\frac12(1,1)=A_1$;
\item $|G|=3$, веса $(\zeta,\zeta^{-1})$ (или $(\zeta^2,\zeta)$ --- та же группа): $\frac13(1,2)=A_2$;
\item $|G|=3$, веса $(\zeta,\zeta)$ (или $(\zeta^2,\zeta^2)$): $\frac13(1,1)$ --- конус над скрученной
кубикой.
\end{enumerate}

\emph{Шаг 3: исключаем $\frac13(1,1)$.} Поверхность $X$ --- $\Q$-горенштейнов предел гладких
поверхностей, поэтому каждая её особенность $\Q$-горенштейново сглаживаема, а по классификации
Коллара--Шеперд-Баррона (теорема~\ref{thm:KSB}, пример~\ref{ex:T-sing}(2)) $\frac13(1,1)$ таковой
не является. (Заметим: $\frac13(1,1)$ сглаживаема \emph{не} $\Q$-горенштейново, и при этом $K^2$
меняется; именно $\Q$-горенштейновость семейств в теореме о K-модулях исключает этот случай.)

\emph{Шаг 4: горенштейновость и антиканоническая модель.} Особенности $A_1$, $A_2$ горенштейновы
(пример~\ref{ex:An-invariants}; упражнение~6.3(8)), поэтому $K_X$ --- дивизор Картье, $X$ ---
горенштейнова дель Пеццо степени 3. По теореме~\ref{thm:Gor-dP} $-K_X$ очень обилен и
$|-K_X|$ вкладывает $X$ в $\PP^3$ как кубическую поверхность (упражнение~6.3(7); это результат
Хидаки--Ватанабе и Фудзиты \cite{Fuj90}; см. также \cite[III.3]{Kol96}). Особенности этой кубики ---
те же $A_1$, $A_2$.
\end{proof}

\begin{remark}
Какие именно конфигурации $A_1$, $A_2$ реализуются на кубиках: $A_1$, $2A_1$, $3A_1$, $4A_1$,
$A_2$, $A_1A_2$, $2A_1A_2$, $2A_2$, $A_12A_2$, $3A_2$ (\S\ref{ss:dP}). Теорема пока не утверждает, что
все они K-полустабильны; это --- содержание следующей теоремы, и ответ: да.
\end{remark}

Итак, объекты $M^{\sm}_{2,3}$ --- поверхности в $\PP^3$. Можно ли тогда в определении K-стабильности
ограничиться тест-конфигурациями, у которых и центральный слой --- кубика в $\PP^3$, т.\,е.
однопараметрическими подгруппами $\PGL_4$? Это в точности ГИТ!

\begin{theorem}[{Одака--Спотти--Сунь \cite{OSS}; ГИТ = K для кубических поверхностей}]\label{thm:GIT=K}
Кубическая поверхность $X\subset\PP^3$ K-полустабильна (соответственно K-полистабильна, K-стабильна)
тогда и только тогда, когда она ГИТ-полустабильна (полистабильна, стабильна) относительно $\SL_4$.
Следовательно,
\[
M^{\sm}_{2,3}\ \cong\ \PP^{19}/\!/\SL_4\ \cong\ \PP(1,2,3,4,5),
\]
и K-полустабильные дель Пеццо степени 3 --- это в точности кубики с особенностями типов $A_1$, $A_2$;
K-стабильные --- кубики с особенностями $A_1$; единственная K-полистабильная, но не K-стабильная ---
кубика $x_0x_1x_2=x_3^3$.
\end{theorem}
\begin{proof}
\emph{Шаг 0: K $\Rightarrow$ ГИТ.} Это теорема Пола--Тяня~\ref{thm:PT}: однопараметрические подгруппы
$\lambda\subset\SL_4$ задают тест-конфигурации $(\cX_\lambda,\cL_\lambda)$ кубики $X$ с
$\Fut(\cX_\lambda)=a\,\mu(F,\lambda)$, $a>0$. Если $X$ K-полустабильна, все $\Fut\ge0$, значит, все
веса Гильберта--Мамфорда $\ge0$, и по теореме~\ref{thm:HM} $X$ ГИТ-полустабильна. Аналогично для
стабильности (все $\Fut>0$ для нетривиальных $\lambda$; тривиальность тест-конфигурации
$\cX_\lambda$ равносильна тривиальности $\lambda$ на $X$, а стабилизатор кубики в $\SL_4$
конечен для ГИТ-стабильных и K-стабильных кубик) и полистабильности.

\emph{Шаг 1: существует K-стабильная кубика, и K-стабильность открыта.} Гладкие кубики K-стабильны
(следствие~\ref{cor:dP-le4} или пример~\ref{ex:AZ-cubic}); по теореме~\ref{thm:openness}
K-стабильные кубики образуют открытое по Зарисскому подмножество $\PP^{19}$ (семейство кубик
$\Q$-горенштейново), содержащее все гладкие.

\emph{Шаг 2: ГИТ-полистабильная $\Rightarrow$ K-полистабильная.} Пусть $X=X_0$ ГИТ-полистабильна
(случаи стабильности и полустабильности аналогичны). Возьмём сглаживание $\cX\to C$ над гладкой
пунктированной кривой $0\in C$ с $\cX_0\cong X$ и гладкими (значит, K-стабильными) слоями $X_t$,
$t\ne0$ --- например, прямую в $\PP^{19}$ через $[X]$ общего направления. По собственности K-модулей
(теорема~\ref{thm:properness}) после конечной замены базы $C'\to C$ существует $\Q$-горенштейново
семейство $\cX'\to C'$ с $\cX'|_{C'\setminus0}\cong\cX|_{C\setminus0}\times_CC'$ и K-полистабильным
слоем $X_0':=\cX'_0$.

\emph{Шаг 3: сравнение объёмов.} Слой $X_0'$ --- K-полустабильная $\Q$-Фано поверхность степени 3,
$\Q$-горенштейнов предел гладких кубик, и по теореме~\ref{thm:cubic-list} $X_0'$ --- кубика в
$\PP^3$ с особенностями $A_1,A_2$. Более того, относительный антиканонический пучок
$\omega_{\cX'/C'}^{-1}$ относительно очень обилен (на каждом слое $-K$ очень обилен, а
$h^0(X_t,-K_{X_t})=4$ не зависит от $t$ по Риману--Роху и обращению в нуль; теорема о замене базы),
так что $\cX'\hookrightarrow\PP^3\times C'$ --- семейство кубических поверхностей. По шагу 0
$X_0'$ ГИТ-полистабильна.

\emph{Шаг 4: единственность полистабильного предела в ГИТ.} Над $C'\setminus0$ семейства
$\cX\times_CC'$ и $\cX'$ изоморфны как абстрактные семейства; их антиканонические вложения в $\PP^3$
отличаются на семейство проективных преобразований $g(t)\in\PGL_4$, $t\ne0$ (антиканоническое
вложение единственно с точностью до $\PGL_4$). Итак, $X_0=\lim_{t\to0}F_t$ и
$X_0'=\lim_{t\to0}g(t)\cdot F_t$ --- два ГИТ-полистабильных предела одного семейства
$\{F_t\}\subset\PP^{19}$. По отделимости ГИТ-фактора (замечание после теоремы~\ref{thm:GIT})
они лежат в одной $\SL_4$-орбите: $X\cong X_0\cong X_0'$. Значит, $X$ K-полистабильна.
\end{proof}

\begin{corollary}\label{cor:cubic-moduli}
Все гладкие кубики K-стабильны (они ГИТ-стабильны по теореме~\ref{thm:Hilbert}) --- третье
доказательство следствия~\ref{cor:dP-le4}. Пространство K-модулей $M^{\sm}_{2,3}=\PP(1,2,3,4,5)$
устроено так:
\begin{itemize}
\item открытое плотное подмножество --- гладкие кубики (4 параметра);
\item страты коразмерности $k=1,2,3,4$ --- кубики с $k$ особенностями $A_1$ (при $k=4$ ---
единственная кубика Кэли);
\item единственная точка $[3A_2]$ --- образ всех кубик с особенностями $A_2$ (они K-полустабильны,
но не K-полистабильны и изотривиально вырождаются в $x_0x_1x_2=x_3^3$ вдоль однопараметрической
подгруппы, теорема~\ref{thm:Hilbert}(3)).
\end{itemize}
По \cite{OSS} это же пространство --- громов-хаусдорфова компактификация модулей КЭ кубических
поверхностей: каждая K-полистабильная кубика несёт (орбифолдную) КЭ-метрику, а кубика $3A_2$ ---
фактор метрики Фубини--Штуди на $\PP^2$ по $\mu_3$.
\end{corollary}

\begin{remark}[Почему доказательство не проходит «в лоб»]
Шаг 2 использует собственность K-модулей, шаг 3 --- локально-глобальную оценку
(теорема~\ref{thm:local-global}) через нормализованный объём, шаг 4 --- отделимость ГИТ. Без
теоремы~\ref{thm:cubic-list} неясно было бы, что K-полистабильный предел снова кубика (а не,
скажем, поверхность с особенностью $\frac13(1,1)$ или более высокого индекса, не вкладывающаяся в
$\PP^3$ антиканонически). Именно оценка $|G|\le3$ --- «индексная оценка» --- делает ГИТ
применимой.
\end{remark}

\begin{theorem}[{Лю--Сюй \cite{LX19}, Лю \cite{Liu22}}]
ГИТ = K-стабильность для кубических гиперповерхностей в $\PP^4$ и $\PP^5$ (трёхмерных и
четырёхмерных кубик). Соответствующие K-модули --- ГИТ-факторы Аллкока и Лазы.
\end{theorem}

\begin{conjecture}
ГИТ = K-стабильность для кубических гиперповерхностей любой размерности. Это следовало бы из
гипотезы о втором по величине локальном объёме (теорема~\ref{thm:nvol-properties}(4) в старших
размерностях).
\end{conjecture}

Схема доказательства в старших размерностях та же: индексная оценка через $\hvol$ показывает, что
K-полустабильные пределы --- снова кубики (это требует знать, что все особенности с
$\hvol\ge(-K)^n/(1+1/n)^n$ --- «не хуже», чем у кубик), а дальше работает ГИТ.

% ============================================================
\section{Заключительные замечания}\label{sec:concl}
% ============================================================

\subsection{Итоги}

Мы прошли путь от определения K-стабильности до явного описания пространства K-модулей кубических
поверхностей. Подведём итоги в виде «дорожной карты» доказательства $M^{\sm}_{2,3}\cong\PP(1,2,3,4,5)$:
\begin{enumerate}
\item \emph{Определение.} K-стабильность через тест-конфигурации и инвариант Футаки
(раздел~\ref{sec:tc}); достаточно специальных тест-конфигураций (Ли--Сюй).
\item \emph{Валюативный критерий} (раздел~\ref{sec:beta}): K-полустабильность $\iff\beta(E)\ge0$
для всех дивизоров над $X$.
\item \emph{Гладкие кубики K-стабильны}: три доказательства --- через $\alpha$-инвариант
(Чельцов, Фудзита; раздел~\ref{sec:alpha}), через теорию Аббана--Чжуана (раздел~\ref{sec:az}),
через ГИТ (теорема~\ref{thm:GIT=K}).
\item \emph{Существование K-модулей}: ограниченность, открытость, собственность, проективность,
редуктивность (раздел~\ref{sec:moduli}).
\item \emph{Индексная оценка}: K-полустабильные дель Пеццо степени 3 имеют лишь особенности с
$|G|\le3$, т.\,е. $A_1,A_2,\frac13(1,1)$; последняя исключается теорией $\Q$-горенштейновых
сглаживаний (раздел~\ref{sec:cubic}).
\item \emph{Антиканоническая модель}: такие поверхности --- кубики в $\PP^3$ (Хидака--Ватанабе,
Фудзита).
\item \emph{ГИТ}: теорема Гильберта (раздел~\ref{sec:git}) описывает (полу)стабильные кубики;
теорема Пола--Тяня и отделимость ГИТ-фактора завершают отождествление K-модулей с ГИТ-фактором.
\end{enumerate}

\subsection{Что дальше}

\begin{itemize}
\item \emph{Другие степени.} Для каждой степени $d$ пространство $M^{\mathrm K}_{2,d}$ сглаживаемых
дель Пеццо описано в \cite{OSS}: при $d\ge5$ --- точка, при $d=4$ --- ГИТ-фактор пучков квадрик
($\cong\PP^2/S_5$...), при $d=3$ --- $\PP(1,2,3,4,5)$, при $d=2,1$ --- явные модификации
ГИТ-факторов квартик и секстик, в которых появляются поверхности с особенностями $\frac14(1,1)$
(и это иллюстрирует, что ГИТ $\ne$ K в общем случае).
\item \emph{Трёхмерные Фано.} Для 105 семейств гладких трёхмерных многообразий Фано
K-(полу/поли)стабильность общего представителя определена в \cite{ACCFKMSV}; описания
пространств K-модулей известны лишь в отдельных случаях (\cite{ADL23}, \cite{DJKHQ} и др.).
\item \emph{Гиперповерхности.} Открытый вопрос: всякая ли гладкая гиперповерхность Фано степени
$\ge3$ K-стабильна? Известно для общей (по открытости и примеру Ферма) и для многих частных
случаев.
\item \emph{Стены.} Для пар $(X,cD)$ K-модули зависят от $c$ и меняются при переходе через
конечное число «стен» (пример~\ref{ex:P112-pair}: K-полустабильность лишь при $c=\frac12$) ---
теория wall-crossing \cite{ADL}, \cite{ADL23}.
\end{itemize}

\subsection{Упражнения обзора, не разобранные в тексте}

Для удобства читателя перечислим, где в тексте решены упражнения обзора \cite{DeV}, и что осталось.
\begin{itemize}
\item \S2.1: (1) --- пример~\ref{ex:terminal-smooth}; (2) --- пример~\ref{ex:lct-curves}; (3) ---
пример~\ref{ex:discrep-P11n}; (4), (5) --- раздел~\ref{sec:mmp}; (6) --- пример~\ref{ex:cone-canonical}.
\item \S3.4: (1a,b) --- \cite[\S\,7]{AG}; (1c) --- классификация дель Пеццо \cite[6.10]{AG}, не
доказана здесь; (2) --- пример~\ref{ex:Aut-F1}; (3) --- \cite[7.7]{AG}; (4) ---
пример~\ref{ex:dP8-dP6-polystable}; (5), (6) --- пример~\ref{ex:beta-F1}; (7) --- пример~\ref{ex:beta-dP7};
(8) --- пример~\ref{ex:beta-P11n}; (9) --- пример~\ref{ex:Bl-cubic-P3}; (10) --- примеры~\ref{ex:P112},
\ref{ex:beta-P11n}, \ref{ex:P112-pair}.
\item \S4.1: (1) --- пример~\ref{ex:AZ-P112}.
\item \S5.1: (1), (2) --- пример~\ref{ex:Markov}; (3) --- раздел~\ref{sec:moduli}; (4) --- обсуждено в
разделе~\ref{sec:moduli} без явного примера.
\item \S6.3: (1) --- пример в \S\ref{sec:cubic}; (2) --- там же (без «бонуса» для $\frac14(1,1)$);
(3) --- пример~\ref{ex:wps-unstable}; (4) --- в доказательстве теоремы~\ref{thm:volume-bound}; (5)
--- пример~\ref{ex:P114-pair}; (6) --- теорема~\ref{thm:Gor-dP} (со ссылкой); (7) --- то же;
(8), (9), (10) --- примеры~\ref{ex:An-invariants}, \ref{prop:HJ}, \ref{ex:P11n}, \ref{ex:discrep-P11n},
\ref{ex:beta-P11n}.
\end{itemize}

\subsection*{Упражнения для самостоятельной работы}

\begin{exercise}
Вычислите $\beta(E)$ для $(-1)$-кривой $E$ на поверхности дель Пеццо степени 5 и $\beta(\ell)$ для
прямой $\ell$ на гладкой кубической поверхности. Ответы: $S(E)=\frac{13}{15}$, $\beta(E)=\frac2{15}$
(при $1\le t\le2$ отрицательная часть разложения Зарисского $-K-tE$ сосредоточена на трёх
$(-1)$-кривых, пересекающих $E$); $S(\ell)=\frac59$, $\beta(\ell)=\frac49$ (здесь $-K-t\ell$ nef при
$t\le1$, $\vol=(1-t)(3+t)$, и псевдоэффективный порог равен $1$: класс $-K-\ell$ --- пучок
коник, вычетных к $\ell$).
\end{exercise}

\begin{exercise}
Пользуясь таблицей Гильберта (предложение~\ref{prop:Hilbert-table}), найдите явную 1-ПП,
дестабилизирующую кубику с особенностью $A_3$, и кубику, особую вдоль прямой.
\end{exercise}

\begin{exercise}
Докажите, что кубика Кэли ($4A_1$) K-стабильна, не пользуясь теоремой~\ref{thm:GIT=K}: оцените
$\delta$ методом Аббана--Чжуана в особых точках.
\end{exercise}

\begin{exercise}
Вычислите барицентр антиканонического многоугольника $\PP(1,1,n)$ и убедитесь, что
$\beta(E)=-\langle b_P,u_E\rangle$ совпадает с примером~\ref{ex:beta-P11n}.
\end{exercise}

\begin{exercise}
Покажите, что тест-конфигурация $\{xy=z^2+tw\}\subset\PP(1,1,1,2)\times\A^1$ многообразия $\PP^2$
(пример~\ref{ex:P2-to-P114}) индуцирована нормированием $v=\ord_C$ --- порядком обращения в
нуль вдоль коники $C=\{xy=z^2\}$: фильтрация $\{s\in H^0(\cO(k)):\ord_C(s)\ge j\}$ имеет
присоединённое градуированное кольцо $\C[x,y,z,w]/(xy-z^2)$ с $\deg w=2$ (здесь $w$ --- класс
$xy-z^2$), а его $\Proj$ --- это $\PP(1,1,4)$. Вычислите $\beta(C)$: $S(C)=\frac19\int_0^{3/2}(3-2t)^2dt=\frac12$,
$\beta(C)=\frac12>0$ --- в согласии с K-полистабильностью $\PP^2$.
\end{exercise}


\begin{thebibliography}{ACCFKMSV}
\bibitem[ABHLX]{ABHLX} J.~Alper, H.~Blum, D.~Halpern-Leistner, C.~Xu, \emph{Reductivity of the automorphism group of K-polystable Fano varieties}, Invent. Math. 222 (2020), 995--1032.
\bibitem[ACC$^+$]{ACCFKMSV} C.~Araujo, A.-M.~Castravet, I.~Cheltsov, K.~Fujita, A.-S.~Kaloghiros, J.~Mart\'inez-Garc\'ia, C.~Shramov, H.~S\"u\ss, N.~Viswanathan, \emph{The Calabi problem for Fano threefolds}, LMS Lecture Note Series 485, Cambridge, 2023.
\bibitem[ADL1]{ADL} K.~Ascher, K.~DeVleming, Y.~Liu, \emph{Wall crossing for K-moduli spaces of plane curves}, Proc. London Math. Soc. 128 (2024), e12615.
\bibitem[ADL2]{ADL23} K.~Ascher, K.~DeVleming, Y.~Liu, \emph{K-stability and birational models of moduli of quartic K3 surfaces}, Invent. Math. 232 (2023), 471--552.
\bibitem[AG]{AG} \emph{Дивизоры, канонический класс, многообразия Фано и кубические поверхности. Конспект по классической алгебраической геометрии над $\C$}, файл \texttt{AG\_konspekt.tex} в той же папке репозитория.
\bibitem[Al]{Alper} J.~Alper, \emph{Good moduli spaces for Artin stacks}, Ann. Inst. Fourier 63 (2013), 2349--2402.
\bibitem[AZ]{AZ} H.~Abban, Z.~Zhuang, \emph{K-stability of Fano varieties via admissible flags}, Forum Math. Pi 10 (2022), e15.
\bibitem[BCHM]{BCHM} C.~Birkar, P.~Cascini, C.~Hacon, J.~McKernan, \emph{Existence of minimal models for varieties of log general type}, J. Amer. Math. Soc. 23 (2010), 405--468.
\bibitem[BHLLX]{BHLLX} H.~Blum, D.~Halpern-Leistner, Y.~Liu, C.~Xu, \emph{On properness of K-moduli spaces and optimal degenerations of Fano varieties}, Selecta Math. 27 (2021), 73.
\bibitem[Bi]{Bir19} C.~Birkar, \emph{Anti-pluricanonical systems on Fano varieties}, Ann. of Math. 190 (2019), 345--463; \emph{Singularities of linear systems and boundedness of Fano varieties}, Ann. of Math. 193 (2021), 347--405.
\bibitem[BJ]{BJ} H.~Blum, M.~Jonsson, \emph{Thresholds, valuations, and K-stability}, Adv. Math. 365 (2020), 107062.
\bibitem[BLX]{BLX} H.~Blum, Y.~Liu, C.~Xu, \emph{Openness of K-semistability for Fano varieties}, Duke Math. J. 171 (2022), 2753--2797.
\bibitem[BMRT]{BMRT} M.~Bate, B.~Martin, G.~R\"ohrle, R.~Tange, \emph{Closed orbits and uniform S-instability in geometric invariant theory}, Trans. Amer. Math. Soc. 365 (2013), 3643--3673.
\bibitem[BW]{BW} J.\,W.~Bruce, C.\,T.\,C.~Wall, \emph{On the classification of cubic surfaces}, J. London Math. Soc. (2) 19 (1979), 245--256.
\bibitem[BX]{BX} H.~Blum, C.~Xu, \emph{Uniqueness of K-polystable degenerations of Fano varieties}, Ann. of Math. 190 (2019), 609--656.
\bibitem[Ch]{Che08} I.~Cheltsov, \emph{Log canonical thresholds of del Pezzo surfaces}, Geom. Funct. Anal. 18 (2008), 1118--1144.
\bibitem[CLS]{CLS} D.~Cox, J.~Little, H.~Schenck, \emph{Toric varieties}, GSM 124, AMS, 2011.
\bibitem[CP]{CP02} И.\,А.~Чельцов, Дж.~Пак, \emph{Глобальные лог-канонические пороги и обобщённые точки Эккардта}, Матем. сб. 193:5 (2002), 149--160.
\bibitem[CS]{CS08} И.\,А.~Чельцов, К.\,А.~Шрамов, \emph{Лог-канонические пороги неособых трёхмерных многообразий Фано} (с приложением Ж.-П.~Демайи), УМН 63:5 (2008), 73--180.
\bibitem[De]{De} M.~Demazure, \emph{Surfaces de del Pezzo I--V}, in: S\'eminaire sur les singularit\'es des surfaces, LNM 777, Springer, 1980.
\bibitem[DeV]{DeV} K.~DeVleming, \emph{On moduli of Fano varieties: an introduction to K-stability and K-moduli}, arXiv:2602.20280 (2026).
\bibitem[dFEM]{dFEM} T.~de Fernex, L.~Ein, M.~Musta\c t\u a, \emph{Multiplicities and log canonical threshold}, J. Algebraic Geom. 13 (2004), 603--615.
\bibitem[DJKHQ]{DJKHQ} K.~DeVleming, L.~Ji, P.~Kennedy-Hunt, M.\,H.~Quek, \emph{The K-moduli space of a family of conic bundle threefolds}, arXiv:2403.09557.
\bibitem[Do]{Do} I.~Dolgachev, \emph{Classical algebraic geometry: a modern view}, Cambridge, 2012.
\bibitem[Do2]{Dolg-IT} I.~Dolgachev, \emph{Lectures on invariant theory}, LMS Lecture Note Series 296, Cambridge, 2003.
\bibitem[Do3]{Dolg-WPS} I.~Dolgachev, \emph{Weighted projective varieties}, in: Group actions and vector fields, LNM 956, Springer, 1982, 34--71.
\bibitem[Don]{Don02} S.\,K.~Donaldson, \emph{Scalar curvature and stability of toric varieties}, J. Differential Geom. 62 (2002), 289--349.
\bibitem[DS]{DS14} S.~Donaldson, S.~Sun, \emph{Gromov--Hausdorff limits of K\"ahler manifolds and algebraic geometry}, Acta Math. 213 (2014), 63--106.
\bibitem[DSz]{DS16} V.~Datar, G.~Sz\'ekelyhidi, \emph{K\"ahler--Einstein metrics along the smooth continuity method}, Geom. Funct. Anal. 26 (2016), 975--1010.
\bibitem[Du]{Durfee} A.~Durfee, \emph{Fifteen characterizations of rational double points and simple critical points}, Enseign. Math. 25 (1979), 131--163.
\bibitem[ELS]{ELS} L.~Ein, R.~Lazarsfeld, K.~Smith, \emph{Uniform approximation of Abhyankar valuation ideals in smooth function fields}, Amer. J. Math. 125 (2003), 409--440.
\bibitem[FO]{FO18} K.~Fujita, Y.~Odaka, \emph{On the K-stability of Fano varieties and anticanonical divisors}, Tohoku Math. J. 70 (2018), 511--521.
\bibitem[Fu1]{Fuj18} K.~Fujita, \emph{Optimal bounds for the volumes of K\"ahler--Einstein Fano manifolds}, Amer. J. Math. 140 (2018), 391--414.
\bibitem[Fu2]{Fuj19a} K.~Fujita, \emph{K-stability of Fano manifolds with not small alpha invariants}, J. Inst. Math. Jussieu 18 (2019), 519--530.
\bibitem[Fu3]{Fuj19b} K.~Fujita, \emph{A valuative criterion for uniform K-stability of $\Q$-Fano varieties}, J. Reine Angew. Math. 751 (2019), 309--338.
\bibitem[Fu4]{Fuj90} T.~Fujita, \emph{On singular del Pezzo varieties}, in: Algebraic geometry (L'Aquila, 1988), LNM 1417, Springer, 1990, 117--128.
\bibitem[Ful]{Fulton} W.~Fulton, \emph{Intersection theory}, 2nd ed., Springer, 1998.
\bibitem[GH]{GH} Ф.~Гриффитс, Дж.~Харрис, \emph{Принципы алгебраической геометрии}, т.~1--2, М.: Мир, 1982.
\bibitem[Ha]{Ha} Р.~Хартсхорн, \emph{Алгебраическая геометрия}, М.: Мир, 1981.
\bibitem[HM]{HM} J.~Harris, I.~Morrison, \emph{Moduli of curves}, GTM 187, Springer, 1998.
\bibitem[HP]{HP} P.~Hacking, Yu.~Prokhorov, \emph{Smoothable del Pezzo surfaces with quotient singularities}, Compos. Math. 146 (2010), 169--192.
\bibitem[HW]{HW} F.~Hidaka, K.~Watanabe, \emph{Normal Gorenstein surfaces with ample anti-canonical divisor}, Tokyo J. Math. 4 (1981), 319--330.
\bibitem[Ji]{Jia20} C.~Jiang, \emph{Boundedness of $\Q$-Fano varieties with degrees and alpha-invariants bounded from below}, Ann. Sci. \'Ec. Norm. Sup\'er. 53 (2020), 1235--1248.
\bibitem[JM]{JM} M.~Jonsson, M.~Musta\c t\u a, \emph{Valuations and asymptotic invariants for sequences of ideals}, Ann. Inst. Fourier 62 (2012), 2145--2209.
\bibitem[Ke]{Kempf} G.~Kempf, \emph{Instability in invariant theory}, Ann. of Math. 108 (1978), 299--316.
\bibitem[KM]{KM} J.~Koll\'ar, S.~Mori, \emph{Birational geometry of algebraic varieties}, Cambridge Tracts in Math. 134, Cambridge, 1998.
\bibitem[Ko1]{Kol96} J.~Koll\'ar, \emph{Rational curves on algebraic varieties}, Ergebnisse 32, Springer, 1996.
\bibitem[Ko2]{KolPairs} J.~Koll\'ar, \emph{Singularities of pairs}, Proc. Sympos. Pure Math. 62 (1997), 221--287.
\bibitem[Ko3]{KolSing} J.~Koll\'ar, \emph{Singularities of the minimal model program}, Cambridge Tracts in Math. 200, Cambridge, 2013.
\bibitem[Ko4]{Kol23} J.~Koll\'ar, \emph{Families of varieties of general type}, Cambridge Tracts in Math. 231, Cambridge, 2023.
\bibitem[KSB]{KSB} J.~Koll\'ar, N.\,I.~Shepherd-Barron, \emph{Threefolds and deformations of surface singularities}, Invent. Math. 91 (1988), 299--338.
\bibitem[La]{La} R.~Lazarsfeld, \emph{Positivity in algebraic geometry I, II}, Springer, 2004.
\bibitem[Li1]{Li17} C.~Li, \emph{K-semistability is equivariant volume minimization}, Duke Math. J. 166 (2017), 3147--3218.
\bibitem[Li2]{Li18} C.~Li, \emph{Minimizing normalized volumes of valuations}, Math. Z. 289 (2018), 491--513.
\bibitem[Liu1]{Liu18} Y.~Liu, \emph{The volume of singular K\"ahler--Einstein Fano varieties}, Compos. Math. 154 (2018), 1131--1158.
\bibitem[Liu2]{Liu22} Y.~Liu, \emph{K-stability of cubic fourfolds}, J. Reine Angew. Math. 786 (2022), 55--77.
\bibitem[LL]{LL19} C.~Li, Y.~Liu, \emph{K\"ahler--Einstein metrics and volume minimization}, Adv. Math. 341 (2019), 440--492.
\bibitem[LWX]{LWX} C.~Li, X.~Wang, C.~Xu, \emph{On the proper moduli spaces of smoothable K\"ahler--Einstein Fano varieties}, Duke Math. J. 168 (2019), 1387--1459.
\bibitem[LX1]{LX} C.~Li, C.~Xu, \emph{Special test configuration and K-stability of Fano varieties}, Ann. of Math. 180 (2014), 197--232.
\bibitem[LX2]{LX19} Y.~Liu, C.~Xu, \emph{K-stability of cubic threefolds}, Duke Math. J. 168 (2019), 2029--2073.
\bibitem[LX3]{LX20} C.~Li, C.~Xu, \emph{Stability of valuations and Koll\'ar components}, J. Eur. Math. Soc. 22 (2020), 2573--2627.
\bibitem[LXZ]{LXZ} Y.~Liu, C.~Xu, Z.~Zhuang, \emph{Finite generation for valuations computing stability thresholds and applications to K-stability}, Ann. of Math. 196 (2022), 507--566.
\bibitem[LZ]{LZ22} Y.~Liu, Z.~Zhu, \emph{Equivariant K-stability under finite group action}, Internat. J. Math. 33 (2022), 2250007.
\bibitem[MFK]{MFK} D.~Mumford, J.~Fogarty, F.~Kirwan, \emph{Geometric invariant theory}, 3rd ed., Springer, 1994.
\bibitem[Mu]{Mukai} S.~Mukai, \emph{An introduction to invariants and moduli}, Cambridge, 2003.
\bibitem[Mum]{Mum61} D.~Mumford, \emph{The topology of normal singularities of an algebraic surface and a criterion for simplicity}, Publ. Math. IHES 9 (1961), 5--22.
\bibitem[Od]{Oda13} Y.~Odaka, \emph{The GIT stability of polarized varieties via discrepancy}, Ann. of Math. 177 (2013), 645--661.
\bibitem[OS]{OS12} Y.~Odaka, Y.~Sano, \emph{Alpha invariant and K-stability of $\Q$-Fano varieties}, Adv. Math. 229 (2012), 2818--2834.
\bibitem[OSS]{OSS} Y.~Odaka, C.~Spotti, S.~Sun, \emph{Compact moduli spaces of del Pezzo surfaces and K\"ahler--Einstein metrics}, J. Differential Geom. 102 (2016), 127--172.
\bibitem[PT]{PT} S.\,T.~Paul, G.~Tian, \emph{CM stability and the generalized Futaki invariant I}, arXiv:math/0605278; \emph{II}, Ast\'erisque 328 (2009), 339--354.
\bibitem[PV]{PV} В.\,Л.~Попов, Э.\,Б.~Винберг, \emph{Теория инвариантов}, Итоги науки и техн. Сер. Соврем. пробл. мат. Фундам. направления, т.~55, М.: ВИНИТИ, 1989, 137--309.
\bibitem[So]{Song} J.~Song, \emph{The $\alpha$-invariant on toric Fano manifolds}, Amer. J. Math. 127 (2005), 1247--1259.
\bibitem[Sz]{Sz} G.~Sz\'ekelyhidi, \emph{An introduction to extremal K\"ahler metrics}, GSM 152, AMS, 2014.
\bibitem[Ti1]{Tia87} G.~Tian, \emph{On K\"ahler--Einstein metrics on certain K\"ahler manifolds with $c_1(M)>0$}, Invent. Math. 89 (1987), 225--246.
\bibitem[Ti2]{Tian-book} G.~Tian, \emph{Canonical metrics in K\"ahler geometry}, Birkh\"auser, 2000.
\bibitem[WZ]{WZ} X.-J.~Wang, X.~Zhu, \emph{K\"ahler--Ricci solitons on toric manifolds with positive first Chern class}, Adv. Math. 188 (2004), 87--103.
\bibitem[Xu1]{Xu} C.~Xu, \emph{K-stability of Fano varieties}, монография (2023), \url{https://web.math.princeton.edu/~chenyang/Kstabilitybook.pdf}.
\bibitem[Xu2]{Xu20} C.~Xu, \emph{A minimizing valuation is quasi-monomial}, Ann. of Math. 191 (2020), 1003--1030.
\bibitem[XZ]{XZ} C.~Xu, Z.~Zhuang, \emph{On positivity of the CM line bundle on K-moduli spaces}, Ann. of Math. 192 (2020), 1005--1068.
\bibitem[Za]{Zar} O.~Zariski, \emph{The theorem of Riemann--Roch for high multiples of an effective divisor on an algebraic surface}, Ann. of Math. 76 (1962), 560--615.
\bibitem[Zh]{Zhu21} Z.~Zhuang, \emph{Optimal destabilizing centers and equivariant K-stability}, Invent. Math. 226 (2021), 195--223.
\end{thebibliography}

\end{document}
